% Version 2.20 of 2017/10/04

\documentclass[runningheads]{llncs}
\pagestyle{plain}
% !TEX root = main.tex
\usepackage{amsmath}
\let\proof\relax
\let\endproof\relax
\usepackage{amsthm}
\usepackage{amssymb}
\usepackage[linesnumbered,ruled,vlined]{algorithm2e}
\usepackage{booktabs} %toprule, bottomrule
\usepackage{bm}

%\usepackage{enumerate}
\usepackage{enumitem}

\usepackage{graphicx}
% Sort numerical citations within each group; retain range compression.
\usepackage[sort]{cite}
\usepackage[breaklinks=true]{hyperref}

\usepackage{makecell}
\usepackage{multirow}
\usepackage{mathtools}

\usepackage{stmaryrd}
% St Mary's Road symbols have no bold font; keep their regular shapes.
\SetSymbolFont{stmry}{bold}{U}{stmry}{m}{n}
\usepackage{threeparttable} % e.g. for tnote

\usepackage{xcolor}
\usepackage{xspace}
\usepackage{arydshln}

\usepackage{boxedminipage}
\usepackage[most]{tcolorbox}
\usepackage{mdframed}
%\usepackage{tikz}
%\usetikzlibrary{matrix,decorations.pathreplacing}

%\usepackage{showframe}
\usepackage{todonotes}

%\usepackage{amsfonts}
%\usepackage{float}
%\usepackage{lipsum}
%\usepackage{url}

%\usepackage[showframe]{geometry}

% used for tikz
\usetikzlibrary{positioning, shapes, arrows.meta, calc, backgrounds, fit, calc}
\usetikzlibrary{decorations.pathreplacing}
\renewcommand{\ttdefault}{lmtt} % use Latin Modern mono instead of cmtt

% !TEX root = main.tex
\newcommand{\methodname}{{\tt Mikros}\xspace}
\newcommand{\qafold}{{\tt QAFold}\xspace}
\newcommand{\zkSNARK}{{zkSNARK}\xspace}
\newcommand{\zkSNARKs}{{\tt zkSNARKs}\xspace}
\newcommand{\litebullet}{{\tt LiteBullet}\xspace}
\newcommand{\codeswitch}{{\tt CodeSwitch}\xspace}
% Define custom commands for Dirac notation
\newcommand{\bra}[1]{\langle #1|}
\newcommand{\ket}[1]{|#1\rangle}
\newcommand{\braket}[2]{\langle \, #1 \,|\,	#2 \,\rangle}
\newcommand{\dirac}[3]{\langle \, #1 \,|\,	#2 \,|\,	#3 \,\rangle}
\newcommand{\diracbra}[2]{\langle \, #1 \,|\,	#2 \,|}
\newcommand{\diracket}[2]{|\,	#2 \,|\,	#1 \,\rangle}
\newcommand{\diracbraket}[4]{\langle \, #1 \,|\,	#2 \,|\,	#3 \,|\, #4 \, \rangle}
\newcommand{\ExtensionDegree}[2]{[ \, #1 : #2 \,]}
\newcommand{\ext}[1]{\overline{#1}}
\ExplSyntaxOn
\NewDocumentCommand{\concat}{m}
 {
  \left[\, \clist_use:nn { #1 } { \mid } \,\right]
 }
\ExplSyntaxOff
\newcommand{\blockmat}[4]{
    \begin{pmatrix}
     #1 &\ \ \ & #2 \\
     #3 &\ \ \ & #4 \\
    \end{pmatrix}
}
\newcommand{\foldmat}[2]{
    \begin{pmatrix}
     \begin{pmatrix}
     #1 \\
     #1 \\
     \end{pmatrix}
     &
     \ \ 
     #2
    \begin{pmatrix}
     #1 \\
     #1 \\
     \end{pmatrix}
    \end{pmatrix}
}

\newcommand{\eg}{{\it e.g.}\xspace}
\newcommand{\etc}{{\it etc.}\xspace}
\newcommand{\ie}{{\it i.e.}\xspace}
\newcommand{\cf}{{\it cf.}\xspace}
\newcommand{\etal}{{\it et~al.}\xspace}
\newcommand{\ifandonlyif}{{\it iff}\xspace}
\newcommand{\resp}{{\it resp.}\xspace}

\newcommand{\bb}{\mathsf{bb}}
\newcommand{\be}{\mathsf{be}}
\newcommand{\ee}{\mathsf{ee}}

\newcommand{\trace}{\mathsf{Tr}_{\ExtensionField/\Field}}



\newcommand{\EqOrNot}{\mathrel{\overset{\mathrm{?}}{=}}}
%\DeclareUnicodeCharacter{\ast}{\op}

%\newcommand{\OP}{{\mathcal{O}\mathcal{P}}}
\newcommand{\op}{{\mathbin{\circledast}}}
\newcommand{\gtimes}{\mathbin{\odot}}
\newcommand{\gadd}{\mathbin{\oplus}}
\newcommand{\gminus}{\mathbin{\ominus}}
\newcommand{\gsum}{\mathop{\bigoplus}}
\newcommand{\kronecker}{\mathbin{\otimes}}
\newcommand{\notationmapsto}{\mathrel{\Rrightarrow}}
\newcommand{\notationmapsfrom}{\mathrel{\Lleftarrow}}

\newcommand{\gzero}{{O}}

\newcommand{\Group}{{\mathbb{G}}}
\newcommand{\Field}{{\mathbb{F}}}
\newcommand{\BaseField}{\mathbb{F}}
\newcommand{\ExtensionField}{\mathbb{E}}
\newcommand{\FieldMinusZO}{{\mathbb{F}\setminus\{0,1\}}}
\newcommand{\CodeSpace}{{\mathbb{C}}}
\newcommand{\Code}{{\mathcal{C}}}
\newcommand{\Zp}{{\mathbb{Z}}_p}
\newcommand{\TensorSpace}{\mathbb{T}}
\newcommand{\IndexSet}{{\mathcal{I}}}
\newcommand{\Root}{{\mathcal{R}}}
\newcommand{\Oracle}{{\mathcal{O}}}
\newcommand{\AFunc}{{\mathcal{F}}}


\newcommand{\ComSpace}{\mathbb{C}}
\newcommand{\HashSpace}{\mathbb{H}}
\newcommand{\PolySpace}{\mathbb{P}}
\newcommand{\FieldStar}{{\Field^*}}
\newcommand{\MSpace}{\mathbb{M}}
\newcommand{\QSpace}{\mathbb{Q}}

\newcommand{\ExpressionSpace}{\mathcal{E}} %only used in 4-expression

\newcommand{\security}{{1^{\lambda}}}

\newcommand{\slice}{{\langle:\rangle}}

\newcommand{\ChooseRandom}{\mathrel{\xleftarrow{\$}}}
\newcommand{\sample}{\mathrel{\xleftarrow{\$}}}
\newcommand{\Definition}{\mathrel{\overset{\text{def}}{=}}}

\newcommand{\Partition}{\xleftarrow{{\sf part}}}
\newcommand{\Round}{\xleftarrow{{\sf round}}}
\newcommand{\Flatten}{\xleftarrow{{\sf flat}}}
\newcommand{\Fold}{\xleftarrow{{\sf fold}}}
\newcommand{\Half}{\xleftarrow{{\sf half}}}
\newcommand{\Broadcast}{\xleftarrow{{\sf broadcast}}}
\newcommand{\Adversary}{\mathcal{A}}
\newcommand{\Prover}{\mathcal{P}}
\newcommand{\Verifier}{\mathcal{V}}
\newcommand{\Emulator}{\mathcal{E}}
\newcommand{\Simulator}{\mathcal{S}}
\newcommand{\Relation}{\mathcal{R}}
\newcommand{\Language}{\mathcal{L}}
\newcommand{\Extractor}{\mathcal{E}}
\newcommand{\Algorithm}{\mathcal{O}}
\newcommand{\Expect}{\mathbb{E}}
\newcommand{\Distinguisher}{\mathcal{D}}
\newcommand{\interact}{{\rightleftharpoons}}
\newcommand{\PVinteract}{{(\mathcal{P} \rightleftharpoons \mathcal{V})}}
\newcommand{\PtoV}{{\mathcal{P} \rightarrow \mathcal{V}}}
\newcommand{\VtoP}{{\mathcal{V} \rightarrow \mathcal{P}}}
\newcommand{\myParbox}[1]{\parbox{0.9\textwidth}{#1}}

\newcommand{\RelationExp}[3]{{\left\{~#1~:~#2~\left|~#3~\right.\right\}}}

%\newcommand{\flag}{{\sf flag}}
%\newcommand{\OpSet}{\mathcal{OP}}

\newcommand{\vectorize}[1]{{\sf vec}(#1)}
\newcommand{\Enc}{{\sf Enc}} %encoding
\newcommand{\diag}{{\sf diag}}
\newcommand{\tr}{{\sf tr}}
\newcommand{\rank}{{\sf rank}}
\newcommand{\col}{{\sf Col}}
\newcommand{\im}{{\sf Im}}
% \newcommand{\det}{{\sf det}}
\newcommand{\trans}{{\sf T}}

\newcommand{\Merkle}{{\sf Merkle}}


\newcommand{\trapdoor}{{\sf trap}} % only in appx-preliminary

\newcommand{\muls}{\mathbb{M}}
\newcommand{\hashes}{\mathsf{H}}
\newcommand{\mpath}{\mathsf{MP}}

\newcommand{\pccost}{{\sf pc}}
\newcommand{\zkcost}{{\sf zk}}
\newcommand{\exps}{\mathbb{E}}
\newcommand{\pairing}{\mathbb{P}} % only in table-evalyn-vs-zkml, table-appx-comp and table-appx-pcs

%\newcommand{\st}{\textrm{s.t.}}

\newcommand{\TRUE}{{\sf true}}
\newcommand{\FALSE}{{\sf false}}
\newcommand{\Accept}{{\sf Accept}}

\newcommand{\pc}{{\sf pc}}
\newcommand{\Polycom}{{\sf pc}}
%only keep \Polycom?

\newcommand{\io}{\bm{io}}
\newcommand{\vk}{{\sf vk}}
\newcommand{\pk}{{\sf pk}}

\newcommand{\tensorspan}[1]{ {\sf ts}(#1)}
\newcommand{\FS}{{\sf FS}}
\newcommand{\Setup}{{\sf Setup}}
\newcommand{\Open}{{\sf Open}}
\newcommand{\Query}{{\sf Query}}
\newcommand{\Prove}{{\sf Prove}}
\newcommand{\Verify}{{\sf Verify}}
\newcommand{\Eval}{{\sf EvalProof}}
\newcommand{\Com}{{\sf Comm}}
\newcommand{\Hash}{{\sf Hash}} %only in appx-preliminary
\newcommand{\Poly}{{\sf poly}}
\newcommand{\Negl}{{\sf negl}}
\newcommand{\pp}{{\sf pp}}
\newcommand{\rand}{{\sf rand}}
\newcommand{\bit}{{\sf bit}}

\newcommand{\suchthat}{\textrm{s.t.}}

%\newcommand{\params}{{\sf params}}
%\newcommand{\ID}{{\sf ID}}
%\newcommand{\len}{{\sf len}}
%\newcommand{\thread}{{\sf thread}}
%\newcommand{\Dec}{{\sf Dec}}
%\newcommand{\VerifyEval}{{\sf VerifyEval}}
%\newcommand{\VerifyPoly}{{\sf VerifyPoly}}
%\newcommand{\VerifyQuery}{{\sf VerifyQuery}}
%\newcommand{\VerifyCom}{{\sf VerifyCom}}
%\newcommand{\TP}{{\sf TP}}
%\newcommand{\TV}{{\sf TV}}
%\newcommand{\Solve}{{\sf Solve}}
%\newcommand{\range}{{\sf range}}
%\newcommand{\Coin}{{\sf Coin}}
%\newcommand{\HCom}{{\sf HComm}}
%\newcommand{\Time}{{\sf Time}}
%\newcommand{\Dist}{{\sf Dist}}
%\newcommand{\vp}{{\sf vp}}

\newcommand{\baseline}{{\#[{\sf baseline}]}}
\newcommand{\baselineAlg}[1]{{\#[{\sf baseline, alg{#1}}]}}

\newcommand{\mcircle}[1]{\mathop{\text{\textcircled{#1}}}}

% \makeatletter
% \newcommand\xleftrightarrow[2][]{%
% 	\ext@arrow 9999{\longleftrightarrowfill@}{#1}{#2}}
% \newcommand\longleftrightarrowfill@{%
% 	\arrowfill@\leftarrow\relbar\rightarrow}
% \makeatother

\newcommand{\customarrow}[1]{\xrightarrow{\makebox[#1]{}}}

\newtheorem{assumption}{Assumption}
\newtheorem{myclaim}{Claim}
\newtheorem{fact}{Fact}

\newtcolorbox{compactbox}[1][]{%
	colback=gray!5!white, % Background color
	colframe=black,		% Border color
	rounded corners,	% Rounded corners
	boxrule=0.5pt,		% Border thickness
	top=2pt,			% Top padding reduced
	bottom=2pt,		% Bottom padding reduced
	left=2pt,			% Left padding reduced
	right=2pt,		% Right padding reduced
	before skip=5pt,	% Space before the box reduced
	after skip=5pt,		% Space after the box reduced
	fonttitle=\bfseries,	% Title font bold
	title={\textbf{#1}},	% Use \textbf to make sure the title is bold
	attach title to upper, % Attach title inline
	% boxed title style={%
	%colback=gray!20!white, % Slightly darker background for title
	%colframe=black,	% Title frame color
	%top=-2pt,				% Title position adjustment
	%bottom=0pt,
	%left=4pt,
	%right=4pt,
	% },
	coltitle=black,		% Title text color
 % #1						% Additional user-defined options
}

\newtcolorbox{coloredalgo}[1][]{
	colback=gray!5!white, % Background color
	colframe=black, % Frame color
	boxrule=0.5pt, % Frame thickness
	#1 % Allows for further customizations passed at the environment call
}

\definecolor{mydarkblue}{HTML}{16468d}
\definecolor{mydarkgreen}{HTML}{226e4b}
\definecolor{myblue}{HTML}{bce2f5}
\definecolor{mygreen}{HTML}{bbead5}
\definecolor{myyellow}{HTML}{fff1dd}
\definecolor{mygrey}{HTML}{d9d9d9}

%\makeatletter
%\renewcommand\@cite[1]%{\leavevmode\unskip[\nobreakspace#1]\unskip\linebreak[0]}
%\makeatother

\newmdenv[
	backgroundcolor=gray!5,	% light gray background
	linecolor=white,		 % no border color
	innerleftmargin=0pt,	 % adjust inner margins
	innerrightmargin=0pt,
	innertopmargin=0pt,
	innerbottommargin=0pt,
	skipabove=5pt,			 % space above the box
	skipbelow=5pt,			% space below the box
        outermargin=0pt,
]{algorithmbackground}


\SetKwInput{KwPubIn}{Public input}
\SetKwInput{KwPriIn}{Private input}

\newlist{tightitemize}{itemize}{1}
\setlist[tightitemize]{label=\textbullet, leftmargin=*, topsep=0pt, partopsep=0pt, itemsep=0pt, parsep=0pt}

\begin{document}

\title{Mikros: Bitwidth-Aware SNARKs over Extensions of Small Fields}
% \subtitle{A}foldinverse
%Efficient zkSNARK for Matrix Expression Through Proof Composition 
% \title{Revisiting Bulletproofs: Surpassing Sumcheck in the zk-SNARK Campaign}
%Matilda: Matrix-Integrated Layered-Data Assurance/Authentication
%Miranda: Matrix-Integrated Recursive Non-interactive Atomic Derivation Assurance

%Sumcheck Replacement, Matrix Expressions, Commitment Reduction

%\titlerunning{}
\author{Paper \#01 \institute{\,} }
% \author{\, \institute{\,}}
% \author{Mingshu Cong
% \inst{\,\,*}
% \inst{1} \orcidID{0000-0001-5858-3411}
% \and
% Sherman S. M. Chow
% \inst{2}
% \orcidID{0000-0002-0629-6792}
% \and
% Tsz Hon Yuen
% \inst{3}
% \orcidID{0000-0002-0629-6792}
% \and
% Siu Ming Yiu
% \thanks{Siu Ming Yiu is the corresponding author.}
% \inst{1} \orcidID{0000-0002-3975-8500}}

% \authorrunning{M. Cong et al.}
% \institute{Department of Computer Science, The University of Hong Kong, Hong Kong SAR, China \\
% \email{\{mscong,smyiu\}@cs.hku.hk}
% \institute{Department of Computer Science, The Chinese University of Hong Kong, Hong Kong SAR, China \\
% \email{smchow@ie.cuhk.edu.hk}
% \institute{Faculty of Information Technology, Monash University, Melbourne, Australia\\
% \email{John.tszhonyuen@monash.edu}
% }
%
\maketitle


\begin{abstract}
Provers in succinct non-interactive arguments of knowledge (SNARKs)
typically perform arithmetic over a large field, even when
their low-bit-width witness entries can be represented in a small field.
Multiplying low-bit-width witness entries by full-bit-width random challenges produces
intermediate values in the large field, increasing
both prover time and memory usage.

We introduce \methodname, a framework for small-field SNARK protocols based
on \emph{vector polynomial packing}.
For inner- and Hadamard-product relations, our reduction transforms
a length-\(N=m\kappa\) vector claim over \(\Field\) into a
length-\(m\) vector claim over a degree-\(\kappa\) extension
\(\ExtensionField\).
The prover encodes the original relation as a low-degree
vector-valued polynomial identity and computes the corresponding
row-polynomial products over \(\Field\).
Evaluation at a random point in \(\ExtensionField\) yields the
shorter relation, which is then proved using an extension-field
protocol.
Assuming \(|\Field|\geq\kappa\), \(m\geq\kappa\), and precomputed
packing bases, both reductions require \(O(\kappa N)\) base-field
operations using schoolbook arithmetic and \(O(N)\) base-field
coordinates of prover memory, excluding commitments and openings.
For smaller fields and power-of-two \(\kappa\), tabulation over
an intermediate field gives the same online Hadamard-reduction
time with additional, witness-independent lookup tables.

We also provide a packed opening protocol for the underlying
polynomial commitment scheme (PCS).
Together with a permutation argument, these components yield
protocols for circuit satisfiability.
For \(N=2^{30}\) over an \(8\)-bit base field, our component
benchmarks show prover speedups of \(5.60\times\) and
\(4.31\times\) for inner and Hadamard products, respectively,
over the small-field-first baseline, with measured resident
memory usage reduced by a factor of \(4.48\) in both cases.
\end{abstract}


\section{Introduction}

A succinct non-interactive argument of knowledge (SNARK) allows a
prover to establish knowledge of a witness using a short proof that
can be verified much faster than the computation it
certifies~\cite{innovations/BitanskyCCT12}.
SNARKs enable efficient verification of computations performed by
untrusted parties; their zero-knowledge variants additionally
hide the witness (see, \eg,~\cite{tcc/BootleCG20,%
eurocrypt/BootleCL22,%
eurocrypt/BuenzFS20,%
eurocrypt/ChenBBZ23,%
eurocrypt/ChiesaHMMVW20,%
eurocrypt/ChiesaOS20,%
iacr/GabizonWC19,%
crypto/GolovnevLSTW23,%
eurocrypt/Groth16,%
crypto/Setty20,%
iacr/SettyL20}).
The prover's cost of producing these SNARK proofs, however, can substantially
exceed the cost of the computation itself.

The prover's use of a large working field is a major source
of this proving overhead.
Many efficient SNARKs combine a polynomial interactive oracle
proof (PIOP) with a polynomial commitment scheme
(PCS)~\cite{tcc/Ben-SassonCS16,eurocrypt/BunzFS20,%
eurocrypt/ChiesaHMMVW20,iacr/GabizonWC19,ccs/MallerBKM19}.
The circuit-satisfiability (Circ-SAT) relation is encoded as low-degree
polynomial identities.
In multilinear constructions, the PIOP uses random challenges
to reduce these identities to evaluation claims about the
multilinear extensions (MLEs) of the witness vectors, which
the PCS then authenticates.
To achieve the target soundness error, these protocols use
a challenge field sufficiently large relative to the polynomial
degrees, even when the witness entries fit in a much smaller field.
Arithmetic involving these challenges produces intermediate
values in the larger field, increasing both prover time
and memory usage.

Small fields and low-bit-width integers arise naturally in many modern
applications of SNARKs.
First, a growing body of zero-knowledge machine learning (zkML)
research generates SNARK proofs for neural-network (NN) computations
(\eg,~\cite{ccs/Balbas0VRS23,%
asiacrypt/CCYY25,%
nips/GhodsiGG17,%
iacr/GuoGQZ25,%
esorics/KeufferMC18,%
uss/Weng0XK021,%
ccs/LiuXZ21,%
iacr/FengQZDC21,%
corr/KangHSS22,%
ccs/GargGJMMPW23,%
tifs/WengWTYL23,%
tdsc/LeeKKO24,%
ccs/SunLZ24,%
ccs/AbbaszadehPK024,%
popets/WangH23,%
icml/WaiwitlikhitS0H24,%
sp/QuGYZ26,%
ccs/RiasiWBVH25,%
uss/GuoLHQTZ25,%
uss/0001SLLGCZ25}).
Modern production-grade NNs commonly use low-precision (\eg, 8-bit)
representations~\cite{corr/micikevicius22}.
Second, a recent line of work combines SNARKs with fully homomorphic
encryption (FHE) to achieve verifiable computation on encrypted
data (\eg, ~\cite{ccs/FioreGP14,%
pkc/FioreNP20,%
pkc/BoisCFK21,%
joc/GaneshNS23,%
cic/AtapoorBPS24,%
wahc/KnabenhansVMH24,%
crypto/GargGW24,%
asiacrypt/AranhaCGS25,%
crypto/CascudoCCFGS25,%
ccs/ThibaultW25,%
cic/GamaBKSV25,%
ccs/ZhangWLXDFL25,%
iacr/PeshawariaLFT25}).
FHE schemes often exploit small-field arithmetic, with plaintext
coordinates in fields such as
\(\mathbb{F}_{65537}\)~\cite{asiacrypt/AranhaCGS25}
and ciphertext coefficients represented via the Chinese remainder
theorem (CRT) over 30--60-bit prime fields.

For computations over small fields, we seek prover time and
memory usage that scale proportionally with the bit width of the base field
rather than that of the soundness field, a property we call
\emph{bitwidth awareness}.
In elliptic-curve-based PCSs such as KZG, Hyrax, and
Dory~\cite{asiacrypt/KateZG10,sp/WahbyTSTW18,tcc/Lee21},
short scalars reduce the number of group operations in the
multiscalar multiplications used to generate
commitments~\cite{eurocrypt/SettyTW24}.
Hash-based constructions exploit binary-field packing and
ring-switching to commit to small-field polynomials without
representing each input value as a full extension-field
element~\cite{eurocrypt/DiamondP25,eurocrypt/DiamondP26}.
Extending these savings to the PIOP layer is more difficult:
random challenges turn small-field inputs into extension-field
intermediate values.
Small-field-first algorithms, constraint packing, and related
sum-check optimizations substantially reduce this
overhead~\cite{iacr/BagadDT24,iacr/DaoT24,iacr/Gruen24,%
iacr/BagadDDT25,iacr/DaoDBDT26}.
After \(t\) optimized rounds, a conventional sum-check tail
operates on tables of length \(N/2^t\) and, for fixed product
degree \(d\geq2\), performs \(\Theta(N/2^t)\) extension-field
multiplications~\cite{crypto/Thaler13,iacr/BagadDT24}.
In a degree-\(\kappa\) extension represented in a polynomial basis,
schoolbook multiplication uses \(\Theta(\kappa^2)\) base-field
multiplications~\cite{iacr/BagadDT24}.
Increasing \(t\) reduces this tail at the expense of the work
in the optimized phase; \(t\) may depend on \(\kappa\).
The prover's total cost therefore depends on both phases.

Field agnosticism does not imply bitwidth awareness.
Field-agnostic constructions such as Brakedown and BaseFold
remove structural restrictions on the arithmetic field, but not
the field-size requirements imposed by
soundness~\cite{crypto/GolovnevLSTW23,crypto/ZeilbergerCF24}.
Their algebraic checks still rely on a sufficiently large
challenge field to achieve the target soundness error.
Supporting a small base field therefore does not by itself
ensure that challenge-dependent arithmetic and intermediate
storage retain the cost advantages of the low-bit-width witness~\cite{eurocrypt/DiamondP26,iacr/BagadDDT25}.
We therefore ask:
\begin{quote}
\emph{Can a SNARK achieve bitwidth-aware prover time and memory
usage while drawing its soundness challenges from a large field?}
\end{quote}

\smallskip\noindent\textbf{Our result.}
Let $\lambda$ be the security parameter, and consider an $N$-gate
arithmetic circuit over a small field of order $q$.
We choose a packing factor $\kappa\simeq\lambda/\log_2 q$,
subject to $\kappa\le q$; for fixed $\lambda$, $\kappa$ scales
inversely with the base-field bit width $\log_2 q$.
For a single instance, our prover uses
$O(\kappa N(\log_q N)^2)$ base-field operations and stores
$O(N\log_q N)$ base-field elements.
For a batch of $\kappa$ instances of the same circuit, the
amortized costs decrease to $O(\kappa N)$ base-field operations
and $O(N)$ base-field elements of memory per instance.
Under schoolbook arithmetic, multiplication in a degree-$\kappa$
extension costs $\Theta(\kappa^2)$ base-field operations, while
each extension-field element occupies $\kappa$ base-field
coordinates.
Thus, the amortized per-instance complexity corresponds to
$O(N/\kappa)$ extension-field multiplication equivalents and
the storage occupied by $O(N/\kappa)$ extension-field elements.

\smallskip\noindent{\bf Our core technique.}
Our \methodname framework is built on a novel technique called \emph{vector polynomial
packing}.
In the base-field instantiations, we partition a length-\(N=m\kappa\) vector over \(\Field\)
into \(\kappa\) blocks of length \(m\), and encode these blocks
as a length-\(m\) vector of polynomials using \(\kappa\)
linearly independent polynomials in \(\Field[X]\), each of
degree at most \(\kappa-1\).
This encoding expresses the original vector relation as a
low-degree polynomial identity involving length-\(m\) polynomial vectors.
Evaluating this identity at a uniformly random point
\(\ext{x}\in\ExtensionField\), where
\([\ExtensionField:\Field]=\kappa\), reduces it to a
length-\(m\) vector relation over \(\ExtensionField\).
For our inner- and Hadamard-product instantiations, an appropriate
choice of packing polynomials allows the prover to establish
the reduced claim using \(O(\kappa N)\) base-field operations.

We further show how to commit to a length-\(N\) vector over
\(\Field\) by packing its entries into a length-\(m\) vector
over \(\ExtensionField\) and committing to it using any
compatible PCS over \(\ExtensionField\).
Combining this packed PCS with our packed inner- and Hadamard-product protocol and
a bitwidth-aware permutation argument yields a bitwidth-aware
SNARK for the PLONK constraint system~\cite{iacr/GabizonWC19} of circuit satisfiability.


\smallskip\noindent{\bf Concrete performance.}
Our component benchmarks compare the packed reductions with a
small-field-first baseline. At \(N=2^{30}\) over an 8-bit base
field, the measured prover speedups are \(5.60\times\) for inner
products and \(4.31\times\) for Hadamard products, with a
\(4.48\times\) reduction in resident memory for both.
These measurements include packing and witness-dependent
preprocessing, and exclude commitments and PCS openings.
Section~\ref{sec:experiment} reports the full comparison across
8-, 16-, 32-, and 64-bit configurations and the measurement procedure.


\subsection{Related Work}
\label{sec:related-work}

\noindent{\bf PIOP--PCS framework for SNARKs.}
Mainstream SNARKs combine a polynomial interactive oracle
proof (PIOP) with a polynomial commitment scheme
(PCS)~\cite{eurocrypt/BunzFS20,eurocrypt/ChiesaHMMVW20}.
In this framework, circuit satisfiability is arithmetized as linear
constraints, such as rank-one constraint systems
(R1CS)~\cite{eurocrypt/GennaroGP013} or PLONK-style constraints~\cite{iacr/GabizonWC19,eurocrypt/ChenBBZ23},
which are then encoded as polynomial identities.
The PIOP uses random verifier challenges to reduce verification
of these identities to polynomial evaluation claims, and the
PCS authenticates the claimed evaluations against committed
polynomials~\cite{asiacrypt/KateZG10,tcc/PapamanthouST13}.
The Fiat--Shamir transform makes the resulting argument non-interactive in the
random-oracle model~\cite{crypto/FiatS86,tcc/Ben-SassonCS16}.

Early arithmetizations encoded constraints using univariate
polynomials~\cite{eurocrypt/GennaroGP013,cacm/ParnoHG016,iacr/GabizonWC19},
whereas subsequent works increasingly adopted multilinear extensions
(MLEs) over the Boolean hypercube~\cite{sp/WahbyTSTW18,crypto/XieZZPS19,crypto/Setty20,eurocrypt/ChenBBZ23}.
At the PIOP layer, polynomial identities arising from MLE-based
arithmetizations are typically reduced to sumcheck claims.
The sumcheck protocol was introduced by
Lund \etal~\cite{jacm/LundFKN92} and subsequently developed into a
fundamental building block for efficient proof
systems~\cite{jacm/GoldwasserKR15,crypto/Thaler13}.
A PCS then binds the prover to the oracle
polynomials and proves the evaluation claims produced by the PIOP.
Our work follows this PIOP--PCS framework and develops
bitwidth-aware techniques for the PIOP layer.

\medskip
\noindent{\bf Bitwidth-aware PCS constructions.}
Group-based PCSs~\cite{asiacrypt/KateZG10,tcc/PapamanthouST13,%
sp/BunzBBPWM18,sp/WahbyTSTW18,tcc/Lee21,asiacrypt/BunzMMTV21}
commit to polynomial coefficients or evaluation vectors using
multiscalar multiplications (MSMs).
MSM algorithms can exploit short integer scalars, making
commitment generation sensitive to input bit
width~\cite{eurocrypt/SettyTW24}.
Opening costs depend on the construction: for a size-$N$ polynomial, Dory requires
\(O(\sqrt{N})\) group operations and \(O(N)\) scalar-field
operations~\cite{tcc/Lee21,eurocrypt/SettyTW24}, where
multiplications between short witness entries and full-width
weights can benefit from bitwidth-aware arithmetic.

Hash-based PCSs~\cite{icalp/Ben-SassonBHR18,crypto/ArnonCFY24,%
eurocrypt/ArnonCFY25,iacr/KohritaT23,ccs/AmesHIYM17,%
focs/Ben-SassonCIKS20,crypto/GolovnevLSTW23,crypto/XieZS22,%
crypto/BlockFK24,iacr/LiLXYY26,crypto/ZeilbergerCF24,%
usenix/GuoLHQTZ25,eurocrypt/BrehmCFRRZ25,%
asiacrypt/NariST25,iacr/BunzFRW25}
commit to linear encodings of polynomial coefficient or
evaluation vectors using Merkle trees.
Their opening protocols combine algebraic consistency checks
with authenticated queries to the encoded vectors, using
code distance to detect inconsistencies.
When encoding is performed over a large field, short inputs
can produce full-width codeword symbols, so commitment
generation does not automatically benefit from low input
bit width.
Recent constructions avoid this large-field embedding
overhead for small-field polynomials.
Binius~\cite{eurocrypt/DiamondP25} adapts Brakedown to
multilinear commitments over binary towers without
coefficient-wise embedding overhead.
FRI-Binius~\cite{iacr/BenjaminDP24} combines the binary tower of Binius with binary-field FRI and BaseFold,
achieving polylogarithmic proof size and verifier complexity.
Ring-switching generalizes this approach by reducing small-field
multilinear openings to openings of packed polynomials over
an extension field.
By packing \(\kappa\) base-field coordinates into each element
of a degree-\(\kappa\) extension, it transforms an extension-field
multilinear PCS into a small-field PCS whose commitment cost
matches that of the underlying scheme on inputs of the same
bit size.

% Complete degree-2 Hadamard audit: a_i b_i - c_i = 0.
% N is the length of EACH of the three input vectors (total input Theta(N)).
% The public equality factor is included but is not a third small-field input.
%
% SFF: 2025/1117, Theorem 6.1/Section 6.2; 2026/587, Figure 2,
% Section 6 and Appendix C.5. With M=N/2^t, the quadratic grid has 3^t
% entries (projective nodes 0,1,infinity). M*3^t mixed products cost
% kappa*N*(3/2)^t; suffix weights and tail cost kappa^2*M;
% grid contractions cost kappa^2*3^t. Streamed binding is absorbed.
% Gruen: 2024/108, Sections 2.3.1 and 5.1--6.1; 2026/587 Section 7.
% The full bound ADDS FFT interpolation N*t, equality weighting and
% binding kappa*N, suffix weights/tail kappa^2*N/2^t, and first-round
% Lagrange/quotient processing kappa^2*2^t. It is not the abstract's
% constraint-evaluation-only bound. This count uses FFT/coset domains;
% the protocol permits general interpolation sets (Section 5.1, footnote 18).
% Flock: 2026/1329, Sections 4.2--4.4 (July 31, 2026 PDF).
% Byte chunks/subfield degree 8; 3<=t<=7, implementation t=6, kappa=128.
% Re-cost lookup LDEs, weighting, binding, output check, and conventional
% tail. Full time adds kappa^2*2^t; RAM adds kappa*2^t (fixed tables
% included). N>=4^t absorbs these terms. This is a derived bound.
% Dao--Thaler: 2024/1038, Sections 7--8, Algorithm 2, composed here with
% CONVENTIONAL sumcheck. This is not a claim about the best possible tail.
% Liu--Zhang: 2025/594, Section 4.2, Theorem 8/Protocol 5 (NOT Table 1).
% Set s=2, h=floor(log_2(kappa)/2), u_i=4^i for i<h.
% Full cost: sum_i N/2^i*(u_i^2+kappa*u_i) + kappa^2*N/2^h
% = O(kappa^(3/2)*N). Generate deterministic prefix weights using the
% O(kappa) special-generator multiplication of 2024/1038 Section 8.
% Store only random-suffix weights and the current folded witness tables:
% O(N + max_i N*2^i + kappa*N/2^h) = O(N*sqrt(kappa)), for N>=kappa.
% RMFE: 2021/694, Theorem 4/Figure 3/Appendix A.2, sparse Boolean R1CS:
% n=3N, m=N, nnz=Theta(N), constant-rate packing and security O(kappa).
% Packed sumcheck: 2025/719 (2025-10-03 version), Lemma 5 explicitly proves
% the Hadamard relation via Lemma 2 (d=2); Figure 3 includes the output c.
% Set r=k, N=k*N_0. Equality-weighted degree is 3, hence q>3(2r-1).
% Stream per-column forward differences; keep O(N+r^2) base coordinates.
% Our row adds the output-vector linear evaluation to Proposition
% proposition:hadamard-complexity, costing O(kappa*N) and O(N+kappa^2).
% Parameter-only additive terms and their absorption conditions are given
% in Appendix appx:hadamard-comparison.
\begin{table*}[!t]
\centering
\setlength{\tabcolsep}{1.5pt}
\renewcommand{\arraystretch}{1.15}
\begin{threeparttable}
\caption{Prover costs for length-\(N\) Hadamard-product proofs}
\label{tab:hadamard-comparison}
\begin{tabular*}{\textwidth}{@{\extracolsep{\fill}}
    >{\raggedright\arraybackslash}p{0.25\textwidth}
    >{\raggedright\arraybackslash}p{0.26\textwidth}
    >{\raggedright\arraybackslash}p{0.20\textwidth}
    >{\raggedright\arraybackslash}p{0.25\textwidth}@{}}
\toprule
Method & Online prover time & RAM
& Applicability \\
\midrule
Classical sumcheck
& \(O(\kappa^2N)\)
& \(O(\kappa N)\)
& General extensions \\
\addlinespace[3pt]
SSF~\cite{iacr/BagadDT24,iacr/BagadDDT25,iacr/DaoDBDT26}
& \(O(\kappa N(3/2)^t +\kappa^2N/2^t)\)
& \(O(N+\kappa N/2^t)\)
& General extensions \\
\addlinespace[3pt]
Dao--Thaler~\cite{iacr/DaoT24}
& \(O(\kappa^2N)\)
& \(O(\kappa N)\)
& Binary tower fields \\
\addlinespace[3pt]
Tower zerocheck~\cite{iacr/LiuZ25}
& \(O(\kappa^{3/2}N)\)
& \(O(N\sqrt{\kappa})\)
& Binary tower fields \\
\addlinespace[3pt]
RMFE--Aurora~\cite{iacr/CascudoG21}
& \(O(\kappa^2N+\kappa N\log N)\)
& \(O(\kappa N)\)
& Binary RMFE \\
\addlinespace[3pt]
Packed sumcheck~\cite{iacr/WeiWXZDWZ25}
& \(O(rN)\)
& \(O(N+r^2)\)
& \(|\Field|>3(2r-1)\) \\
\addlinespace[3pt]
Gruen~\cite{iacr/Gruen24}
& \(O(Nt+\kappa N +\kappa^2N/2^t)\)
& \(O(N+\kappa N/2^t)\)
& $|\Field| \ge 2^{t+1}-1$ \\
\addlinespace[3pt]
Flock~\cite{iacr/BunzRW26}
& \(O(N2^t+\kappa N+\kappa^2N/2^t)\)
& \(O(N+\kappa N/2^t)\)
& \(\Field=\mathbb F_2\) \\
\midrule
\methodname
& \(O(\kappa N)\)
& \(O(N)\)
& General extensions \\
\bottomrule
\end{tabular*}
\begin{tablenotes}[para,flushleft]
\footnotesize
\item Here \(\kappa=[\ExtensionField:\Field]\) and
\(t_{\mathbb K}=[\mathbb K:\Field]\) is chosen as in
Section~\ref{sec:hadamard}.
Arithmetic and RAM are measured in base-field operations
and elements, respectively.
\(t\) counts optimized
or skipped rounds, and \(r\) is the packed-sumcheck parameter.
Bounds use schoolbook arithmetic and coordinate-charged table accesses;
RAM includes inputs and read-only tables. Public preprocessing is separate.
Parameter-only additive terms are absorbed for sufficiently large \(N\):
\(N\geq6^t\) for SFF and \(N\geq4^t\) for univariate skip and Flock.
Both \methodname{} rows assume \(N\geq\kappa^2\).
Flock uses \(\mathbb F_{256}\subseteq\ExtensionField\) and
\(3\leq t\leq7\) (\(t=6\) in its implementation).
Complete bounds, implementation assumptions, and soundness-dependent
parameter choices are given in Appendix~\ref{appx:hadamard-comparison}.
\end{tablenotes}
\end{threeparttable}
\end{table*}

\medskip
\noindent{\bf Low-bit-width PIOP optimizations.}
In sumcheck-based PIOPs, random challenges produce
extension-field intermediate values even when the witnesses
lie in a small base field.
Existing approaches reduce this overhead in several ways.

Small-field-first methods precompute products over the base
field and postpone extension-field arithmetic, but switching
to conventional sumcheck after \(t\) optimized rounds still
leaves \(\Theta(N/2^t)\) extension-field multiplications for
fixed product degree \(d\geq 2\), while increasing \(t\)
requires more base-field
preprocessing~\cite{iacr/BagadDT24,iacr/BagadDDT25,iacr/DaoDBDT26}.
Gruen's univariate-skipping technique replaces several initial
rounds with one higher-degree check, reducing prover work
at the cost of higher-degree verifier
interpolation~\cite{iacr/Gruen24}.
Flock applies univariate skipping to batched Boolean computations
and computes binary-input low-degree extensions using chunked
lookup tables~\cite{iacr/BunzRW26}. It compresses these tables
using translation invariance, improving cache locality without
changing the number of vector XORs. Its zerocheck also uses
fixed equality weights and aggregates the linear output term
before interpolation.

Binary-tower constraint packing combines base-field constraints
using linearly independent extension-field
weights~\cite{eurocrypt/DiamondP25,iacr/DaoT24}.
Dao and Thaler~\cite{iacr/DaoT24} use inexpensive multiplication
by tower generators to reduce equality-table construction to
approximately \(N/\kappa\) general extension-field multiplications,
but their direct variable elimination requires individual
degree one and therefore does not yield the same bound for
proving general quadratic relations.
Liu and Zhang~\cite{iacr/LiuZ25} specialize sumcheck to Boolean
inputs over binary towers, achieving \(O(\kappa^{3/2}N)\)
prover time under quadratic multiplication costs, but this
bound relies on the binary tower structure over
\(\mathbb{F}_2\).


Reverse multiplication-friendly embeddings allow coordinatewise
Boolean products to be recovered from packed extension-field
products, but the resulting Boolean R1CS protocols also require
specialized checks of linear relations modulo a
subspace~\cite{iacr/CascudoG21}.

Packed sumcheck instead keeps computation in a small prime
field by combining parallel repetition with folding, but
the larger soundness error of individual small-field checks
requires additional work for soundness amplification and
for combining the repeated
claims~\cite{iacr/WeiWXZDWZ25}.


\medskip
\noindent{\bf Relation to our work.}
Our work combines packing at the relation and commitment layers.
For inner and Hadamard products, vector polynomial packing
encodes operand blocks as row polynomials and uses auxiliary
polynomials to account for cross-column products.
In the base-field instantiations, the prover computes row
products over \(\Field\), then evaluates at a random point in
\(\ExtensionField\) to reduce
the length-\(N\) claim to a length-\(N/\kappa\) inner-product
claim.
Our packed opening protocol authenticates these
challenge-dependent operands against a fixed-basis witness
commitment.
Together with a packed permutation argument, these components
support PLONK constraints while maintaining compact witness
representations.
Sections~\ref{sec:overview} and~\ref{sec:pcs} describe the
reductions and their commitment interface.
For \(|\Field|<\kappa\) and power-of-two \(\kappa\), our
univariate Hadamard construction
uses an intermediate field and interpolates only after weighted
aggregation. Combining blockwise linear-map tables, as used for
binary LDEs in Flock~\cite{iacr/BunzRW26}, with intermediate-field
multiplication tables yields \(O(N\kappa+\kappa^3)\) online
base-field operations. The construction and its table costs
are given in Section~\ref{sec:hadamard}.


\subsection{Our Contributions}

\noindent\textbf{Vector polynomial packing.}
    We introduce a technique that reduces length-\(N=m\kappa\)
    inner- and Hadamard-product relations over \(\Field\) to
    length-\(m\) inner-product relations over a degree-\(\kappa\)
    extension \(\ExtensionField\).
    In the base-field instantiations, the prover multiplies row polynomials over \(\Field\) before
    evaluating them at an extension-field challenge.
    This keeps the polynomial-product computation in the small field
    and confines the ensuing extension-field proof to vectors shorter
    by a factor of \(\kappa\).

\medskip \noindent \textbf{End-to-end bitwidth-aware construction.}
    We compose our packed reductions and permutation argument with
    a compatible bitwidth-aware PCS through ring-switching-based
    openings to obtain arguments for PLONK-style circuit satisfiability.
    We give constructions for both a single instance and a batch of
    instances of the same circuit.
    Packing across instances preserves the shared linear constraints,
    yielding lower amortized prover costs.
    Section~\ref{sec:together} gives soundness and efficiency bounds
    for both compositions, including their dependence on the
    underlying PCS.

\medskip\noindent \textbf{Independently useful Hadamard-product and permutation proofs.}
    Our bitwidth-aware protocols for coordinatewise multiplication
    and multiset equality apply beyond the circuit arithmetization
    considered here.
    For \(q=|\Field|<\kappa\), a tabulated intermediate-field
    Hadamard reduction uses \(O(N\kappa)\) online base-field
    operations and \(O(N)\) working base-field coordinates,
    assuming \(\kappa\) is a power of two and \(N\geq\kappa^2\).
    It retains the original witness commitments; its reusable
    tables and their preprocessing costs are accounted for
    separately in Section~\ref{sec:hadamard}.
    The permutation proof uses packed Hadamard checks to verify
    grand-product identities while retaining compact witness storage.
    It also provides a reusable component for permutation-based
    lookup arguments.

\section{Technical Overview}
\label{sec:overview}

Let \(\ExtensionField\) be a degree-\(\kappa\) extension of a
small field \(\Field\), large enough to meet the soundness
requirements at security parameter \(\lambda\).
Our central technique, \emph{algebraic vector packing}, reduces
length-\(N=m\kappa\) base-field relations to length-\(m\)
extension-field relations.
The challenge is to preserve the base-field relation.
We use fixed extension-field bases for linear relations and
polynomial bases for bilinear relations.
The two routes share a packed commitment interface,
so their proof protocols compose into an end-to-end bitwidth-aware SNAKR for Circ-SAT.

\subsection{Algebraic Packing of Small-Field Vectors via Two Routes}
\label{sec:dualPhase}

We partition \(\bm a,\bm b\in\Field^N\) into \(\kappa\) blocks
of length \(m\), and arrange them as columns:
\[
    \bm A=\concat{\bm a^{(1)},\ldots,\bm a^{(\kappa)}},
    \qquad
    \bm B=\concat{\bm b^{(1)},\ldots,\bm b^{(\kappa)}}
       \in\Field^{m\times\kappa}.
\]
There are two methods to combine these columns into length-\(m\) vectors.

\medskip\noindent\textbf{Extension-field packing for linear relations.}
Fix an \(\Field\)-basis
\(\ext{\bm\beta}=(\ext\beta_1,\ldots,\ext\beta_\kappa)^\top\)
of \(\ExtensionField\) and set
\(\ext{\bm a}=\bm A\ext{\bm\beta}\).
Any $\Field$-linear map
\(\mathcal L:\Field^m\to\Field^{m'}\), applied across the columns, can be represented by a matrix
\(\bm L\in\Field^{m'\times m}\).
For \(\bm c^{(j)}\in\Field^{m'}\), linear independence
of the basis give:
\[
    \mathcal L(\bm a^{(j)})=\bm c^{(j)}
       \quad\forall j\in[\kappa]
    \quad\Longleftrightarrow\quad
    \bm L\ext{\bm a}
       =\sum_{j=1}^{\kappa}\bm c^{(j)}\ext\beta_j.
\]
Thus one extension-field vector equation certifies all \(\kappa\)
base-field equations.
Here, the requirement that \(\bm L\) have base-field entries is
essential.

\medskip\noindent\textbf{Polynomial packing for bilinear relations.}
For a bilinear map \(\mathcal B\), applying
\(\mathcal B(\bm A\ext{\bm\beta},\bm B\ext{\bm\beta})\)
introduces cross-column products weighted by
\(\ext\beta_i\ext\beta_j\).
These do not, in general, encode the desired coordinatewise
relations in the original basis.

We instead retain a formal indeterminate \(X\).
Choose linearly independent polynomials
\(\phi_1,\ldots,\phi_\kappa\in\Field[X]\), tailored to the
relation, and write
\(\bm\phi=(\phi_1,\ldots,\phi_\kappa)^\top\).
The vector polynomial packings are then:
\begin{equation}
    \bm g_A(X)=\bm A\bm\phi(X) \in\Field[X]^m,
    \qquad
    \bm g_B(X)=\bm B\bm\phi(X)
       \in\Field[X]^m.
\label{eq:overview-packing}
\end{equation}
Unlike multiplication in a fixed extension-field basis,
polynomial multiplication does not reduce modulo an
irreducible polynomial.
We can therefore isolate the desired relation in a single coefficient,
or in the values at a set of interpolation points, while auxiliary polynomials absorb the remaining terms.

For the inner-product and Hadamard-product construction, this gives an identity
\(\mathcal F(\bm g_A(X),\bm g_B(X))=0\), where
\(\mathcal F\) incorporates the claimed output and auxiliary
polynomials.
With the witness commitments fixed, the prover sends the auxiliary
polynomials before the verifier samples \(\ext x\in\ExtensionField\).
Evaluation at \(\ext x\) reduces the identity to a bilinear
relation on length-\(m\) extension-field vectors, which is proved over \(\ExtensionField\).

\subsection{Instantiations for Concrete Relations}

\medskip\noindent\textbf{Packed inner-product proof.}
For \(c=\langle\bm a,\bm b\rangle
=\sum_{j=1}^{\kappa}\langle\bm a^{(j)},\bm b^{(j)}\rangle\),
choose:
\[
    \bm\phi_A(X)=(1,X,\ldots,X^{\kappa-1})^\top,
    \qquad
    \bm\phi_B(X)=(1,X^{-1},\ldots,X^{-(\kappa-1)})^\top.
\]
In \(\langle\bm g_A(X),\bm g_B(X)\rangle\), the matching
columns contribute to the constant coefficient, while
cross-column products have nonzero exponents.
The claim therefore holds if and only if there exist
\(h_+,h_-\in\Field[X]\), each of degree at most
\(\kappa-2\), such that:
\[
    \langle\bm g_A(X),\bm g_B(X)\rangle
       =c+Xh_+(X)+X^{-1}h_-(X^{-1}).
\]
Neither auxiliary term can alter the constant coefficient.
After \(h_+,h_-\) are fixed, a random
\(\ext x\sample\ExtensionField^\times\) reduces this identity
to one length-\(m\) inner-product claim.

\medskip\noindent\textbf{Packed Hadamard-product proof.}
For \(\bm a^{(j)}\circ\bm b^{(j)}=\bm c^{(j)}\), a random
tensor row weight \(\ext{\bm l}\) reduces each vector
equality to a scalar claim.
Writing
\(\ext d^{(j)}=\langle\ext{\bm l},\bm c^{(j)}\rangle\),
we must prove:
\[
    \ext d^{(j)}
       =\langle\ext{\bm l},
          \bm a^{(j)}\circ\bm b^{(j)}\rangle,
       \qquad j\in[\kappa].
\]
For \(|\Field|\geq\kappa\), choose distinct roots
\(\mathcal H=\{\gamma_1,\ldots,\gamma_\kappa\}\subseteq\Field\)
and their Lagrange basis \((\psi_1,\ldots,\psi_\kappa)\).
Use \(\bm\phi=(\psi_1,\ldots,\psi_\kappa)^\top\) in
Eq.~\eqref{eq:overview-packing}, and set
\(g_C(X)=\sum_{j=1}^{\kappa}\ext d^{(j)}\psi_j(X)\) and
\(\Psi(X)=\prod_{j=1}^{\kappa}(X-\gamma_j)\).
Since evaluation at \(\gamma_j\) recovers column \(j\),
all scalar claims hold exactly when:
\[
    \ext{\bm l}^{\top}
       \bigl(\bm g_A(X)\circ\bm g_B(X)\bigr)-g_C(X)
       =\Psi(X)h(X),
    \qquad \deg h\leq\kappa-2.
\]
Here \(h\in\ExtensionField[X]\).
The factor \(\Psi\) prevents the quotient from changing any
claimed value at the nodes.
Fixing \(g_C,h\) before sampling
\(\ext x\sample\ExtensionField\setminus\mathcal H\) leaves
a length-\(m\) weighted inner product over $\ExtensionField$.

In both reductions, the packed polynomial products in each row are computed over
\(\Field\) before aggregation with extension-field weights.
With precomputed base-polynomial evaluations, the prover uses
\(O(N\kappa+\kappa^3)\) base-field operations.
For \(m\geq\kappa\), this is \(O(N\kappa)\):
the final length-\(m\) proof costs \(O(m)\) extension-field
operations, or \(O(N\kappa)\) base-field operations under
schoolbook arithmetic.
Section~\ref{sec:hadamard} also gives a tabulated
intermediate-field variant for smaller fields, with its
preprocessing and table storage accounted for separately.

\medskip\noindent\textbf{Packed PCS.}
$\kappa$ base field vectors are committed to via a packed
\(\ext{\bm a}=\bm A\ext{\bm\beta}\).
However, the inner-product and Hadamard-product reductions
use challenge-dependent packing weights.
To connect these different packing weights, consider the column MLE claims
\(\ext{\bm\xi}^{\top}\bm A=\ext{\bm v}^{\top}\), where
\(\ext{\bm\xi}=\tensorspan{\ext{\bm x}}\) and
\(\ext{\bm x}\in\ExtensionField^{\log_2m}\).
We decompose
\(\ext{\bm\xi}=\bm\Xi\ext{\bm\beta}\) and
\(\ext{\bm v}^{\top}=\ext{\bm\beta}^{\top}\bm U\),
with \(\bm\Xi\in\Field^{m\times\kappa}\) and
\(\bm U\in\Field^{\kappa\times\kappa}\).
The extension-field claims thus become a base-field matrix identity,
to which the linear packing route applies:
\[
    \ext{\bm\xi}^{\top}\bm A=\ext{\bm v}^{\top}
    \quad\Longleftrightarrow\quad
    \bm\Xi^\top\bm A=\bm U
    \quad\Longleftrightarrow\quad
    \bm\Xi^\top\ext{\bm a}=\bm U\ext{\bm\beta}.
\]
After \(\ext{\bm v}\) is fixed, a random projection reduces
the last identity to a length-\(m\) inner product and one
extension-field PCS opening.
The verifier computes the MLE value of $\bm{\Xi}$ using the
tensor structure of \(\ext{\bm\xi}\), through a recursion on its
\(\kappa\times\kappa\) coordinate matrices.

\medskip\noindent\textbf{Packed permutation proof.}
A permutation relation is reduced to equality of random
grand-product fingerprints:
\[
    \prod_{i=1}^{N}(\ext z+a_i)
       =\prod_{i=1}^{N}(\ext z+b_i),
       \qquad \ext z\sample\ExtensionField.
\]
For powers-of-two \(m,\kappa\) with \(m\geq\kappa\), we
retain \(Z\) as a variable during the first \(\log_2\kappa\)
levels of the product tree.
Each level multiplies pairs of rows, so its coefficients
remain in \(\Field\).
At that depth, only \(m\) polynomial entries remain.
Evaluating them at \(\ext z\) gives a length-\(m\)
extension-field tree.

The proof traverses these levels in reverse, using the packed
Hadamard construction for each transition.
At each interpolation point, the prover forms the row-product
polynomials over the small field, aggregates their coefficients
with the row weights, and only then evaluates at \(\ext z\).

\medskip\noindent\textbf{Packed constraint system.}
These components support both single-instance Circ-SAT arguments and a
batched argument of multiple executions of the same circuit.
For a single execution, we reduce the PLONK gate constraints
to weighted Hadamard-product claims and encode witness values
together with their positions for the permutation proof.

For a batched argument, let \(N\) denote the number of gates per instance.
Denote the gate left inputs \(\bm{w}_L \in \Field^N\), right inputs
\(\bm{w}_R \in \Field^N\), and outputs \(\bm{w}_O \in \Field^N\), and set
\(\bm w_M\coloneqq \bm w_L\circ\bm w_R\).
The gate constraints of \(\kappa\) executions can be packed into
\(\ext{\bm w}_L,\ext{\bm w}_R,\ext{\bm w}_M,\ext{\bm w}_O\)
using \(\ext{\bm\beta}\).
Let \(\bm{q}_L, \bm{q}_R, \bm{q}_M, \bm{q}_O, \bm{q}_C\) be the shared selector
vectors that pick out the operands for each operation. All gate equations
hold if and only if:
\[
    \bm q_L\circ\ext{\bm w}_L
    +\bm q_R\circ\ext{\bm w}_R
    +\bm q_M\circ\ext{\bm w}_M
    +\bm q_O\circ\ext{\bm w}_O
    +\left(\textstyle\sum_{j=1}^{\kappa}\ext\beta_j\right)\bm q_C
       =\bm0_N.
\]
This is still a length-\(N\) vector identity over \(\ExtensionField\) and can be
proved using \(O(N)\) extension-field multiplications. Amortized over \(\kappa\)
instances, it reduces to an average of \(O(N\kappa)\) base-field multiplications
per instance.
\section{Preliminaries}


\subsection{Notation}
\label{sec:notation}
Let \(\Field\) be a small finite field of order \(q\), and let
\(\ExtensionField=\mathsf{GF}(q^\kappa)\) be a degree-\(\kappa\)
extension of \(\Field\). Thus \(q=|\Field|\) is a prime power
and \(\Field\subset\ExtensionField\).
For security parameter \(\lambda\), we choose \(\kappa\) to
satisfy the soundness bounds of our protocols, including
\(
    \lvert\ExtensionField\rvert
    =
    q^\kappa
    \geq
    2^\lambda.
\)
We write \(N\) for the witness length over \(\Field\) and set
\(m=N/\kappa\), assuming \(\kappa\mid N\).
For the main circuit instantiations, we require \(q>\kappa\).
This condition concerns the field order \(q\).
The intermediate-field Hadamard variant in
Section~\ref{sec:hadamard} also treats \(q<\kappa\); its tabulated
implementation reports online time, table construction, and total
storage separately.
Our main instantiations use odd prime fields;
Section~\ref{sec:binary-packing} treats binary fields using a small
extension as the working field.

We use lowercase letters for scalars, bold lowercase letters for
column vectors, and bold uppercase letters for matrices.
An overline denotes an element, vector, or matrix over
\(\ExtensionField\), as in \(\ext a\), \(\ext{\bm a}\), and
\(\ext{\bm A}\). Transposition is denoted by \(\top\).

Let
\(\overline{\beta}_1,\ldots,\overline{\beta}_\kappa\)
be a basis of \(\ExtensionField\) over \(\Field\), and let
\(\overline{\beta}_1^\dagger,\ldots,
\overline{\beta}_\kappa^\dagger\)
denote its trace-dual basis.
The polynomial rings over \(\Field\) and \(\ExtensionField\) are
denoted by \(\Field[X]\) and \(\ExtensionField[X]\), respectively.
We use \(X,Y,Z\) for formal indeterminates,
\(\overline{x},\overline{y},\overline{z}\) for extension-field
challenges, and \(i,j,\ell\) for indices.
% Sets are denoted using blackboard bold font,\eg,
% The set of natural numbers is $\mathbb{N}$,
% the set of powers of two is $2^\mathbb{N}$,
% and the set of real numbers is $\mathbb{R}$.
We write $[n]=\{1,\ldots,n\}$ and use $|\cdot|$ for cardinality.


We also use the following matrix and vector operations:
\begin{itemize}[noitemsep, topsep = 0pt]
\item inner product $c = \langle \bm a, \bm b \rangle$; addition ($\bm C = \bm A + \bm B$); transposition ($\bm C = \bm A^\top$);
\item multiplication
($\bm C = \bm A\bm B$); Kronecker (tensor) product ($\bm C = \bm A \otimes \bm B$);
\item Hadamard (element-wise) product $\bm{A} \circ \bm{B}$;
\item concatenation of vectors $\bm{c} = (\bm c_1 \| \bm c_2 \|\cdots \| \bm c_\kappa)$; and
\item horizontal packing of column vectors $\bm C = [\bm c_1 \mid \bm c_2 \mid \cdots \mid \bm c_\kappa]$.
\end{itemize}
    
We charge \(\bb\), \(\be\), and \(\ee\) for a base-field,
mixed-field, and extension-field multiplication, respectively,
and absorb additions into these costs. Schoolbook extension-field
arithmetic gives
\(
    \be=\kappa\bb
\) and
\(
    \ee=\kappa\be=\kappa^2\bb
\).
Tabulated implementations additionally allow random-access tables,
charging one unit per base-field coordinate read or written.
We distinguish online arithmetic, public table construction,
and total storage, including read-only tables.


We denote the prover and verifier by $\Prover$ and $\Verifier$.
We specify relations as
$\Relation=\{\text{public input}:\text{private input}\mid\text{constraint}\}$.

\subsection{SNARK Protocol in Vector Form}
\label{sec:defgeneral}
We express circuit satisfiability as vector identities and represent
multilinear-extension (MLE) evaluations as inner products with
tensor vectors. This formulation exposes the matrix structure used
by our packing reductions. We first recall the PLONK constraint
system and then state the evaluation and proof interfaces.

\begin{definition}[PLONK constraint system for circuit
satisfiability~\cite{iacr/GabizonWC19,iacr/AmbronaSTW22}]
\label{def:plonk}
An arithmetic circuit with \(N\) gates and \(n_{\mathsf{io}}\)
public inputs and outputs is specified by five selector vectors
\(
    \bm q_L,\bm q_R,\bm q_M,\bm q_O,\bm q_C
    \in\Field^N
\)
and an augmented wiring permutation
\(
    \sigma:
    [3N+n_{\mathsf{io}}]
    \to
    [3N+n_{\mathsf{io}}],
\)
encoding wire-copy constraints and connections to the public inputs
and outputs.
Given a public-I/O vector
\(\bm{\mathsf{io}}\in\Field^{n_{\mathsf{io}}}\), let
\(\bm w_L,\bm w_R,\bm w_O\in\Field^N\) denote the left, right,
and output wire assignments, respectively, and define the
augmented assignment
\(
    \bm w
    \coloneqq
    (\bm w_L\mathbin{\|}\bm w_R\mathbin{\|}\bm w_O
    \mathbin{\|}\bm{\mathsf{io}})
    \in\Field^{3N+n_{\mathsf{io}}}.
\)
The wire assignment satisfies the circuit when:
\begin{align}
    \bm q_L\circ\bm w_L
    +\bm q_R\circ\bm w_R
    +\bm q_M\circ(\bm w_L\circ\bm w_R)
    +\bm q_O\circ\bm w_O
    +\bm q_C
    &=\bm 0_N,
    \label{eq:plonk1}\\
    \sigma(\bm w)
    &=\bm w.
    \label{eq:plonk2}
\end{align}
Here,
\(\sigma(\bm w)\) denotes the coordinate permutation of
\(\bm w\) induced by \(\sigma\).
\end{definition}

\begin{definition}[R1CS constraint system for circuit satisfiability]
\label{def:r1cs}
An arithmetic circuit is specified by three sparse matrices
\(
    \bm S_L,\ \bm S_R,\ \bm S_O \in \Field^{N\times N},
\)
where \(N\) is the total number of multiplication gates and I/O gates
(the public inputs and outputs). Let
\(
    \bm w = (\bm w^{\mathsf{mul}} \,\|\, \bm{\mathsf{io}}) \in \Field^{N}
\)
be the assignment to the multiplication gates and the I/O gates.
The circuit is satisfied if and only if:
\[
    (\bm S_L \bm w) \circ (\bm S_R \bm w) = \bm S_O \bm w .
\]
\end{definition}

The following tensor representation identifies an MLE query with a
linear functional on its evaluation vector.

\begin{definition}[Tensor span]
\label{def:tensorspan}
Let \(n \in 2^\mathbb{N}\) and let
\(
    \bm x
    =
    (x_1,\ldots,x_{\log_2 n})
    \in\Field^{\log_2 n}.
\)
The \emph{tensor span} of \(\bm x\) is the length-\(n\) vector
\(
    \tensorspan{\bm x}
    \coloneqq
    \bigotimes_{i=1}^{\log_2 n} (1-x_i, x_i)^\top
    \in\Field^n.
\)
The same construction applies to
\(\ext{\bm x}\in\ExtensionField^{\log_2 n}\).
We define the sets of tensor spans over \(\Field\) and
\(\ExtensionField\), respectively, by:
\begin{align*}
    \TensorSpace(n)
    &\coloneqq
    \left\{
        \tensorspan{\bm x}
        \;\middle|\;
        \bm x\in\Field^{\log_2 n}
    \right\}
    \subseteq\Field^n,
    \qquad
    \ext{\TensorSpace}(n)
    \coloneqq
    \left\{
        \tensorspan{\ext{\bm x}}
        \;\middle|\;
        \ext{\bm x}\in\ExtensionField^{\log_2 n}
    \right\}
    \subseteq\ExtensionField^n.
\end{align*}
\end{definition}

\paragraph{Products of tensor weights.}
We also allow tensor queries with arbitrary two-entry factors.
Coordinatewise products preserve this representation:
\begin{equation}
    \tensorspan{\ext{\bm z}}\circ\tensorspan{\ext{\bm y}}
    =\bigotimes_{j=1}^{\log_2n}
    \begin{pmatrix}
        (1-\ext z_j)(1-\ext y_j)\\
        \ext z_j\ext y_j
    \end{pmatrix}.
    \label{eq:tensor-weight-product}
\end{equation}
When \((1-\ext z_j)(1-\ext y_j)\neq0\) for every \(j\),
factoring out these entries gives:
\[
    \tensorspan{\ext{\bm z}}\circ\tensorspan{\ext{\bm y}}
    =\left(\prod_j(1-\ext z_j)(1-\ext y_j)\right)
      \bigotimes_j
      \begin{pmatrix}
          1\\
          \dfrac{\ext z_j}{1-\ext z_j}
          \dfrac{\ext y_j}{1-\ext y_j}
      \end{pmatrix}.
\]
To express the result as a scaled MLE query, set
\(\ext d_j=(1-\ext z_j)(1-\ext y_j)+\ext z_j\ext y_j\).
If every \(\ext d_j\neq0\), then:
\[
    \tensorspan{\ext{\bm z}}\circ\tensorspan{\ext{\bm y}}
    =\left(\prod_j\ext d_j\right)
      \tensorspan{(\ext z_j\ext y_j/\ext d_j)_j}.
\]
Our protocols retain the two-entry factors in
Eq.~\eqref{eq:tensor-weight-product}, which also covers zero
entries and vanishing normalization factors without division.

\begin{lemma}[Schwartz-Zippel-DeMillo-Lipton Lemma]
\label{lemma:schwartz}
Let
\( \phi(\bm{x}) : \Field^\upsilon \to \Field \) be a nonzero \( \upsilon \)-variate polynomial such that the degree of \( \phi \) in each variable is at most \( d \). Let \( \bm{x} \ChooseRandom \mathbb{X}^{\upsilon} \) for some finite set \( \mathbb{X} \subseteq \Field \). Then:
\[
    \Pr_{\bm{x} \ChooseRandom \mathbb{X}^{\upsilon}}\!\left[ \phi(\bm{x}) = 0 \right]
    \le \frac{\upsilon \cdot d}{|\mathbb{X}|}.
\]
\end{lemma}

\begin{corollary}[Zero check in matrix form]
\label{corollary:schwartz}
Let \(n\) be a positive power of two and let
\(
    \bm a\in\Field^n\setminus\{\bm 0_n\}.
\)
Viewing \(\bm a\) as a vector over \(\ExtensionField\), sample
\(
    \ext{\bm x}
    \ChooseRandom
    \ExtensionField^{\log_2 n}.
\)
Then:
\[
    \Pr_{\ext{\bm x}\ChooseRandom
        \ExtensionField^{\log_2 n}}
    \!\left[
        \left\langle
            \bm a,
            \tensorspan{\ext{\bm x}}
        \right\rangle
        =0
    \right]
    \leq
    \frac{\log_2 n}{\lvert\ExtensionField\rvert}.
\]
\end{corollary}

\begin{definition}[Inner-product proof]
\label{def:ip}
Let \(m\) be a positive power of two. Suppose that the prover
\(\Prover\) is bound to two vectors
\(\ext{\bm a},\ext{\bm b}\in\ExtensionField^m\), and that the
verifier \(\Verifier\) holds the corresponding commitments.
Given a claimed value \(\ext{v}\in\ExtensionField\), an
\emph{inner-product proof} is an interactive protocol for the
relation:
\[
    \Language_{\mathsf{ip}}
    \coloneqq
    \left\{
        \ext{v}\in \ExtensionField,~      \ext{\bm a},\ext{\bm b}\in\ExtensionField^m
        \;\middle|\;
        \begin{array}{l}
            \ext{v}
            =\langle\ext{\bm a},\ext{\bm b}\rangle
        \end{array}
    \right\}.
\]
\end{definition}

Sum-check~\cite{crypto/Thaler13} and folding~\cite{asiacrypt/CCYY25}
give the following reduction to two MLE claims. We represent
sum-check round polynomials by their coefficients; the round check
\(g(0)+g(1)=c\) then applies in every characteristic.

\begin{fact}[\bf Inner-product reduction]
\label{fact:ip}
There is an inner-product reduction with prover cost
\(4m+O(\log m)\) multiplications in \(\ExtensionField\), communication
\(2\log_2 m+2\) $\ExtensionField$-elements, and verifier cost
\(O(\log m)\) $\ExtensionField$-multiplications.
Over \(\log_2 m\) rounds, the verifier samples challenges collected
in
\(
    \ext{\bm x}\in\ExtensionField^{\log_2 m},
\)
and the protocol reduces the inner-product relation to two
terminal MLE claims:
\[
    \ext{v}_1
    \EqOrNot
    \left\langle
        \ext{\bm a},
        \tensorspan{\ext{\bm x}}
    \right\rangle
    \quad\wedge\quad
    \ext{v}_2
    \EqOrNot
    \left\langle
        \ext{\bm b},
        \tensorspan{\ext{\bm x}}
    \right\rangle,
\]
where the prover supplies \(\ext{v}_1,\ext{v}_2\in\ExtensionField\).
The reduction has soundness error at most
\(
    \frac{2\log_2 m}{\lvert\ExtensionField\rvert}.
\)
\end{fact}

\begin{definition}[Multilinear polynomial commitment scheme in matrix form]
\label{def:pcs}
Let \(n \in 2^\mathbb{N}\). A multilinear PCS over \(\Field\)
with evaluations over \(\ExtensionField\) consists of algorithms
$(\Setup,\Com)$ and an opening protocol $\Open$:
\begin{itemize}[leftmargin=*,labelsep=0.5em]
    \item
    \(
        \pp\leftarrow\Setup(1^\lambda,n)
    \)
    takes the security parameter and the committed vector length
    \(n\), and outputs public parameters \(\pp\).

    \item
    \(
        C\leftarrow\Com(\pp,\bm a;\rho)
    \)
    takes a vector \(\bm a\in\Field^n\) and fresh randomness
    \(\rho\), and outputs a commitment \(C\). We write
    \(\Com(\bm a;\rho)\) when \(\pp\) is clear from context.

    \item The opening protocol
    \(
        \mathsf{bit}
        \leftarrow
        \Open\!\left(
            \pp,C,\ext{v},\ext{\bm x}
            :
            \bm a,\rho
        \right)
    \)
    is an interactive argument between \(\Prover\) and
    \(\Verifier\) for the relation:
    \[
        \Relation_{\mathsf{open}}
        \coloneqq
        \left\{
            \bigl(
                \pp,C,\ext{v},\ext{\bm x}
            \bigr)
                :
            \bigl(
                \bm a,\rho
            \bigr)
            \;\middle|\;
            \begin{aligned}
                C=\Com(\pp,\bm a;\rho) \wedge
                \ext{v}=\left\langle
            \bm a,
            \tensorspan{\ext{\bm x}}
        \right\rangle
            \end{aligned}
        \right\}.
    \]
    At the end of the protocol, \(\Verifier\) outputs
    \(\mathsf{bit}\in\{\TRUE,\FALSE\}\).
\end{itemize}
\end{definition}



\subsection{Extension-Field Arithmetic}

\begin{fact}[Base-Field Embedding]
\label{fact:injective}
Let\/ \(\ExtensionField/\Field\) be an extension.
The canonical inclusion
\(
    \iota:\Field\hookrightarrow\ExtensionField
\)
is injective and induces an injective inclusion
\(
    \Field[X]\hookrightarrow\ExtensionField[X].
\)
\end{fact}

The field trace maps extension-field elements back to \(\Field\)
and gives an explicit formula for their base-field coordinates.
\begin{definition}[Extension-field trace]
\label{def:trace}
The \emph{field trace} from \(\ExtensionField = \mathsf{GF}(q^\kappa)\) to \(\Field = \mathsf{GF}(q)\)
is the \(\Field\)-linear map:
\[
    \trace:\ExtensionField\to\Field,
    \qquad
    \ext{\alpha}
    \longmapsto
    \trace(\ext{\alpha})
    \coloneqq \ext{\alpha} + \ext{\alpha}^q + \ext{\alpha}^{q^2} + \cdots + \ext{\alpha}^{q^{\kappa-1}}.
\]
In particular, for all \(a,b\in\Field\) and
\(\ext{\alpha},\ext{\beta}\in\ExtensionField\):
\[
    \trace\!\left(
        a\ext{\alpha}+b\ext{\beta}
    \right)
    =
    a\trace(\ext{\alpha})
    +
    b\trace(\ext{\beta}).
\]
\end{definition}

\begin{definition}[Basis and trace-dual basis]
\label{def:basis}
An ordered tuple
\(
    \mathcal B
    =
    \bigl(
        \ext{\beta}_1,\ldots,\ext{\beta}_\kappa
    \bigr)
    \in\ExtensionField^\kappa
\)
is a \emph{basis} of \(\ExtensionField\) over \(\Field\) if every
\(\ext{\alpha}\in\ExtensionField\) admits a unique representation:
\[
    \ext{\alpha}
    =
    \sum_{i=1}^{\kappa}a_i\ext{\beta}_i,
    \qquad
    a_1,\ldots,a_\kappa\in\Field.
\]
The \emph{trace-dual basis} of \(\mathcal B\) is the unique basis
\(
    \mathcal B^\dagger
    =
    \bigl(
        \ext{\beta}_1^\dagger,\ldots,
        \ext{\beta}_\kappa^\dagger
    \bigr)
\)
satisfying:
\[
    \trace\!\left(
        \ext{\beta}_i\ext{\beta}_j^\dagger
    \right)
    =
    \delta_{ij}
    \qquad
    \text{for all }i,j\in[\kappa],
\]
where \(\delta_{ij}\) is the Kronecker delta. Consequently, the
coordinates of \(\ext{\alpha}\) with respect to \(\mathcal B\)
can be recovered as:
\[
    a_i
    =
    \trace\!\left(
        \ext{\alpha}\ext{\beta}_i^\dagger
    \right),
    \qquad
    i\in[\kappa].
\]
\end{definition}

Our base-field instantiations use interpolation nodes in \(\Field\), so that the Lagrange
basis has base-field coefficients even when the interpolated
polynomial has extension-field values. The same definition over
an intermediate field \(\mathbb K\) is used for the small-field
variant in Section~\ref{sec:hadamard}.


\begin{definition}[Lagrange basis over the extension field]
\label{def:lagrange}
Assume that \(q\geq\kappa\), and fix an ordered set of distinct
nodes \(\mathcal H=\{\alpha_1,\ldots,\alpha_\kappa\}\subseteq\Field\).
The indices \(j\in[\kappa]\) label the chosen field elements.
The vanishing polynomial of \(\mathcal H\) is:
\[
    \Psi(X)
    \coloneqq
    \prod_{j=1}^{\kappa}(X-\alpha_j)
    \in\Field[X].
\]
For each \(j\in[\kappa]\), define the Lagrange basis polynomial
\(
    \psi_j(X)
    \coloneqq
    \prod_{\substack{\ell\in[\kappa]\\ \ell\neq j}}
    \frac{X-\alpha_\ell}{\alpha_j-\alpha_\ell}
    \in\Field[X].
\)
These polynomials satisfy:
\[
    \psi_j(\alpha_i)=\delta_{ij}
    \qquad
    \text{for all }i,j\in[\kappa].
\]
Moreover, \(\{\psi_j(X)\}_{j=1}^{\kappa}\) forms an
\(\ExtensionField\)-basis for the space of polynomials in
\(\ExtensionField[X]\) of degree less than \(\kappa\).
Consequently, every \(g(X)\in\ExtensionField[X]\) with
\(\deg g<\kappa\) admits the unique representation:
\[
    g(X)
    =
    \sum_{j=1}^{\kappa}
    \ext{a}_j\psi_j(X),
    \qquad
    \ext{a}_j=g(\alpha_j)\in\ExtensionField.
\]
\end{definition}

\section{Building Blocks}
\label{sec:blocks}

SNARK constructions typically reduce circuit satisfiability, expressed through constraint systems such as R1CS or PLONK, to matrix or vector relations involving inner products, Hadamard products, and permutations.
These relations can be encoded as polynomial identities and proved using sumcheck, or proved using other specialized protocols, such as folding-based inner-product arguments~\cite{sp/BunzBBPWM18}.
In this section, we assume prover protocols for inner and Hadamard products over \(\ExtensionField\) whose prover costs are linear in the vector length, measured in extension-field operations.
We develop a unified vector polynomial packing technique that reduces base-field relations to shorter extension-field relations, yielding bitwidth awareness by applying these protocols to the packed instances.
In Section~\ref{sec:together}, we combine these building blocks into a SNARK protocol for PLONK constraints.

\subsection{Packed Inner-Product Proof}
\label{sec:ip}

We introduce vector polynomial packing by the packed inner-product proof that reduces a
length-\(N\) inner-product claim over \(\Field\) to a
length-\(m\) claim over \(\ExtensionField\), where \(N=m\kappa\).
Consider the claim:
\[
    c=\langle\bm a,\bm b\rangle,
    ~
    \bm a,\bm b\in\Field^N.
\]
We partition each vector into \(\kappa\) blocks of length \(m\), and
arrange these blocks as:
\begin{equation}
\label{eq:matAB}
    \bm A
    \coloneqq
    \concat{
        \bm a^{(1)},\ldots,\bm a^{(\kappa)}
    },
    \qquad
    \bm B
    \coloneqq
    \concat{
        \bm b^{(1)},\ldots,\bm b^{(\kappa)}
    }
    \in\Field^{m\times\kappa}.
\end{equation}
We use the monomial basis
\(
    \bm{\phi}_A(X)
    =
    (1,X,\ldots,X^{\kappa-1})
\)
and the reciprocal monomial basis
\(
    \bm{\phi}_B(X)
    =
    (1,X^{-1},\ldots,X^{-(\kappa-1)})
\)
to pack the columns into a polynomial vector and a Laurent
polynomial vector, respectively:
\[
    \bm g_A(X)
    \coloneqq
    \sum_{j=1}^{\kappa}\bm a^{(j)}X^{j-1}
    \in\Field[X]^m,
    \qquad
    \bm g_B(X)
    \coloneqq
    \sum_{j=1}^{\kappa}\bm b^{(j)}X^{-(j-1)}
    \in\Field[X^{-1}]^m.
\]
Their inner product defines the Laurent polynomial:
\begin{equation}
    f(X)
    \coloneqq
    \left\langle
        \bm g_A(X),\bm g_B(X)
    \right\rangle.
\label{eq:ip-laurent-polynomial}
\end{equation}
Its constant coefficient is
\(\sum_{j=1}^{\kappa}
\langle\bm a^{(j)},\bm b^{(j)}\rangle
=\langle\bm a,\bm b\rangle\),
while all cross-column terms have nonzero exponents.
Thus, \(c=\langle\bm a,\bm b\rangle\) if and only if there
exist \(h_+,h_-\in\Field[X]\), each of degree at most
\(\kappa-2\), such that:
\begin{equation}
    f(X)
    =
    c
    +Xh_+(X)
    +X^{-1}h_-(X^{-1}).
\label{eq:ip-correction-identity}
\end{equation}

With \(c\) and the commitments to \(\bm A\) and \(\bm B\)
fixed, the prover sends the coefficients of \(h_+\) and \(h_-\).
The verifier then samples
\(\ext{x}\sample\ExtensionField^\times\).
The prover evaluates both encodings at \(\ext{x}\) to obtain:
\[
    \ext{\bm a}
    \coloneqq
    \bm g_A(\ext{x}) \in \ExtensionField^m,
    \qquad
    \ext{\bm b}
    \coloneqq
    \bm g_B(\ext{x})
    \in\ExtensionField^m.
\]
Evaluating Eq.~\eqref{eq:ip-correction-identity} at \(\ext{x}\)
reduces the polynomial identity to the length-\(m\)
inner-product claim:
\begin{equation}
    \left\langle
        \ext{\bm a},\ext{\bm b}
    \right\rangle
    =
    c
    +\ext{x}h_+(\ext{x})
    +\ext{x}^{-1}h_-(\ext{x}^{-1}).
\label{eq:packed-inner-product-check}
\end{equation}
The verifier computes the RHS from \(c\),
the auxiliary polynomials, and \(\ext{x}\).
The parties then apply the extension-field inner-product
protocol of Fact~\ref{fact:ip}, reducing this claim to two
MLE claims on the length-\(m\) packed vectors.
These terminal claims are authenticated against the
original witness commitments via a packed PCS.
Algorithm~\ref{alg:ip} in Appendix~\ref{appx:pseudocode}
gives the packed inner-product protocol.

\begin{proposition}[Packed inner-product proof]
\label{proposition:ip-complexity}
Algorithm~\ref{alg:ip} is a proof protocol for the base-field inner-product relation
with soundness error at most:
\[
    \frac{2\kappa-2}{|\ExtensionField|-1}
    +
    \frac{2\log_2 m}{|\ExtensionField|}.
\]
Excluding the terminal MLE openings, its prover time cost is
\(
    O(N \kappa + \kappa^3)\bb
\)
and its peak prover memory is \(O(N+\kappa^2) \Field\).
\end{proposition}

We prove Proposition~\ref{proposition:ip-complexity}
in Appendix~\ref{appx:proof-ip}.


\subsection{Packed Hadamard-Product Proof}
\label{sec:hadamard}

We next apply vector polynomial packing, using a different choice of
basis polynomials, to a weighted
Hadamard-product relation:
\begin{equation}
    \ext{c}
    =
    \left\langle
        \ext{\bm\xi},
        \bm a\circ\bm b
    \right\rangle \in \ExtensionField,
    \qquad
    \ext{\bm\xi}
    =
    \ext{\bm r}\kronecker\ext{\bm l} \in \ExtensionField^N,
\label{eq:hadamard-claim}
\end{equation}
where
\(
    \bm a,\bm b\in\Field^N,~
    \ext{\bm l}\in\ExtensionField^m,~
    \ext{\bm r}\in\ExtensionField^\kappa,
\)
and
\(
    N=m\kappa.
\)
When expressed using MLE polynomials, this weighted
Hadamard relation corresponds to the degree-two
zerocheck setting considered by Dao and Thaler~\cite{iacr/DaoT24}
and Gruen~\cite{iacr/Gruen24}.

The tensor-structured weight vector factors as
\(\ext{\bm\xi}=\ext{\bm r}\kronecker\ext{\bm l}\).
Using the block matrices
\(\bm A,\bm B\in\Field^{m\times\kappa}\)
defined in Eq.~\eqref{eq:matAB}, we can therefore rewrite
Eq.~\eqref{eq:hadamard-claim} in matrix form as:
\[
    \ext{c}
    =
    \ext{\bm l}^{\top}
    (\bm A\circ\bm B)
    \ext{\bm r}
    \in\ExtensionField.
\]

To prove this identity, the prover sends
\(
    \ext{\bm c}^{\top}
    \coloneqq
    \ext{\bm l}^{\top}(\bm A\circ\bm B) \in \ExtensionField^\kappa.
\)
The verifier checks
\(\ext{c}\EqOrNot
\langle\ext{\bm c},\ext{\bm r}\rangle\).
We then encode these \(\kappa\) column claims at distinct
interpolation nodes to verify them jointly,
where the $\kappa$ entries of $\ext{c}$ is denoted by $\ext{c}_1, \cdots, \ext{c}_\kappa \in \ExtensionField$.

We choose both \(\bm{\phi}_A(X)\) and \(\bm{\phi}_B(X)\)
to be the Lagrange basis
\((\psi_1(X),\ldots,\psi_\kappa(X))\)
from Definition~\ref{def:lagrange}, and let \(\Psi(X)\)
denote the corresponding vanishing polynomial.
We define:
\[
\begin{aligned}
    \bm g_A(X)
    &\coloneqq
    \sum_{j=1}^{\kappa}
    \bm a^{(j)}\psi_j(X),
    \quad
    \bm g_B(X)
    \coloneqq
    \sum_{j=1}^{\kappa}
    \bm b^{(j)}\psi_j(X),
    \quad
    g_C(X)
    \coloneqq
    \sum_{j=1}^{\kappa}
    \ext{c}_j\psi_j(X).
\end{aligned}
\]
When \(|\Field|\ge 2\kappa-1\), we choose \(\kappa\) distinct
interpolation nodes in \(\Field\), so the Lagrange basis
polynomials have coefficients in \(\Field\).
Consequently:
\[
    \bm g_A(X),\bm g_B(X)\in\Field[X]^m,
    \qquad
    g_C(X)\in\ExtensionField[X].
\]

We define the residual polynomial:
\[
    f(X)
    \coloneqq
    \left\langle
        \ext{\bm l},
        \bm g_A(X)\circ\bm g_B(X)
    \right\rangle
    -
    g_C(X)
    \in\ExtensionField[X].
\]
At each interpolation root \(\gamma_j\in\mathcal H\),
the property of the Lagrange basis gives:
\[
    f(\gamma_j)
    =
    \left\langle
        \ext{\bm l},
        \bm a^{(j)}\circ\bm b^{(j)}
    \right\rangle
    -
    \ext{c}_j.
\]
Thus, all \(\kappa\) column claims hold if and only if
\(f\) vanishes on \(\mathcal H\), or equivalently,
\(\Psi(X)\) divides \(f(X)\).
Since \(\deg f\le 2\kappa-2\) and \(\deg\Psi=\kappa\),
this is equivalent to the existence of
\(h(X)\in\ExtensionField[X]\) with \(\deg h\le\kappa-2\)
such that
\begin{equation}
    f(X)=\Psi(X)h(X).
\label{eq:hadamard-quotient}
\end{equation}

For valid claims, the values of \(f\) on \(\mathcal H\)
are known to be zero.
When \(|\Field|\ge 2\kappa-1\), the prover can therefore
recover \(f\) by evaluating it at \(\kappa-1\) distinct
points in \(\Field\setminus\mathcal H\) and interpolating
these values together with the known zeros.
Dividing the resulting polynomial by \(\Psi\) then gives \(h\).

After the prover sends \(\ext{\bm c}\) and \(h\), the verifier
samples
\(
    \ext{x}
    \sample
    \ExtensionField\setminus\mathcal H.
\)
The Lagrange evaluations give the packing weights:
\[
    \ext{\bm\rho}(\ext{x})
    \coloneqq
    \bigl(
        \psi_1(\ext{x}),\ldots,\psi_\kappa(\ext{x})
    \bigr)^\top \in \ExtensionField^\kappa.
\]
The prover computes the packed vectors:
\[
\begin{aligned}
    \ext{\bm a}
    &\coloneqq
    \bm A\ext{\bm\rho}(\ext{x})
    =
    \bm g_A(\ext{x}) \in \ExtensionField^m,
    \qquad
    \ext{\bm b}
    \coloneqq
    \bm B\ext{\bm\rho}(\ext{x})
    =
    \bm g_B(\ext{x})  \in \ExtensionField^m,
\end{aligned}
\]
and the verifier directly computes a degree-$(\kappa-1)$ polynomial value:
\[
    \ext{c}_x
    \coloneqq
    g_C(\ext{x})
    =
    \left\langle
        \ext{\bm c},
        \ext{\bm\rho}(\ext{x})
    \right\rangle \in \ExtensionField.
\]
Evaluating Eq.~\eqref{eq:hadamard-quotient} at \(\ext{x}\)
reduces the Hadamard-product claim to:
\begin{equation*}
    \left\langle
        \ext{\bm a},
        \ext{\bm b}\circ\ext{\bm l}
    \right\rangle
    =
    \ext{c}_x
    +
    \Psi(\ext{x})h(\ext{x}).
\end{equation*}
This is again a length-\(m\) inner-product relation over
\(\ExtensionField\).
Applying the inner-product protocol of Fact~\ref{fact:ip}
reduces it to two MLE claims on
\(\ext{\bm a}\) and \(\ext{\bm b}\).

Algorithm~\ref{alg:hadamard} in Appendix~\ref{appx:pseudocode}
specifies the protocol.

\iffalse
Fact~\ref{fact:ip} reduces this length-\(m\) inner product to two
MLE claims. Write its challenge point and terminal values as:
\[
    \ext{\bm y}\in\ExtensionField^{\log_2m},
    \qquad
    \ext{v}_1,\ext{v}_2\in\ExtensionField.
\]
The terminal claims on the original vectors are:
\begin{align}
\label{eq:hadamard-terminal}
    \ext{v}_1
    &\EqOrNot
    \left\langle
        \bm a,
        \ext{\bm\rho}(\ext{x})
        \kronecker
        \tensorspan{\ext{\bm y}}
    \right\rangle,
    \qquad
    \ext{v}_2
    \EqOrNot
    \left\langle
        \bm b,
        \ext{\bm\rho}(\ext{x})
        \kronecker
        \left(
            \ext{\bm l}
            \circ
            \tensorspan{\ext{\bm y}}
        \right)
    \right\rangle.
\end{align}

The second row weight is represented by multiplying the corresponding
two-entry factors, as in Eq.~\eqref{eq:tensor-weight-product};
it need not be a normalized MLE weight.
Both claims fit the packed-opening interface in
Eq.~\eqref{eq:pcs-opening-relation}.
Algorithm~\ref{alg:hadamard} in Appendix~\ref{appx:pseudocode}
additionally aggregates the column claims at a fresh MLE point,
retaining these row factors.
\fi

\smallskip\noindent\textbf{Small fields via an intermediate field.}
When \(q=|\Field|<\kappa\), we choose the interpolation roots
in an intermediate field.
Starting from the desired extension degree \(\kappa\),
choose \(t=2\lceil\log_q\kappa\rceil-1\), increasing \(t\)
if necessary until
\(
    q^t\ge  2 t\left\lceil\frac{\kappa}{t}\right\rceil - 1.
\)
Then round \(\kappa\) up to the nearest multiple of \(t\).
We reuse \(\kappa\) for the adjusted degree, which satisfies
\(t\mid\kappa\), \(q^t\ge 2\kappa-1\), and
\(t=O(\log_q\kappa)\).
This gives the field tower:
\[
    \Field\subseteq\mathbb K=\mathbb F_{q^t}
    \subseteq\ExtensionField,
    \qquad
    [\ExtensionField:\mathbb K]=\kappa/t.
\]

Choose \(\mathcal H\subseteq\mathbb K\) of size \(\kappa\)
and define \(\bm\psi(X),\Psi(X),\bm g_A(X),\bm g_B(X),g_C(X)\) as above.
After zero padding if needed, arrange the witnesses as
\(\bm A,\bm B\in\Field^{m\times\kappa}\), with \(N=m\kappa\).
The witness entries remain over \(\Field\), while:
\[
    \bm g_A (X),\bm g_B (X)\in\mathbb K[X]^m,
    \qquad
    g_C(X)\in\ExtensionField[X].
\]

We implement the interpolation procedure using
precomputed addition and multiplication tables over
\(\mathbb K\), which also handle inputs in
\(\Field\times\mathbb K\).
Together, these tables require \(O(t^2q^{2t})\)
base-field preprocessing operations and \(O(tq^{2t})\)
base-field coordinates of storage.
Caching the evaluation and interpolation data requires
an additional \(O(\kappa^2t)\) coordinates.
We charge each table lookup as $O(1)$.

Under this convention, evaluating the witness encodings
and aggregating their row products costs \(O(N\kappa)\).
Evaluating \(g_C\), interpolating \(f\), and computing
the quotient cost \(O(\kappa^3/t)\).
Including the remaining protocol steps, but excluding PCS
commitments and openings, the online prover time and
total storage are
\[
    T_{\Prover}=O(N\kappa+\kappa^3),
    \qquad
    S_{\Prover}=O(N+\kappa^2t+tq^{2t}),
\]
measured in base-field operations and coordinates.
The storage bound includes all precomputed data and is
achieved by processing sample points and witness rows
sequentially.
For \(m\ge\kappa\), the online prover time simplifies
to \(O(N\kappa)\).

The packed polynomial identity, challenge sampling over
\(\ExtensionField\), and terminal extension-field inner-product
proof carry over to the adjusted parameters.
The same soundness and complexity bounds for the
remaining steps therefore apply.

\begin{proposition}[Packed Hadamard-product proof]
\label{proposition:hadamard-complexity}
Let \(q=|\Field|\), \(N=m\kappa\), and
\(t=[\mathbb K:\Field]\), where \(\mathbb K\) is the
interpolation field and may equal \(\Field\).
Algorithm~\ref{alg:hadamard} is complete for the weighted
Hadamard-product relation.
With the terminal MLE claims verified, its soundness error
is at most:
\[
    \frac{2\kappa-2}{|\ExtensionField|-\kappa}
    +
    \frac{2\log_2N}{|\ExtensionField|}.
\]
With the precomputed data described above, its online
prover cost, excluding terminal MLE openings, is
\(
    O(N\kappa+\kappa^3)\,\bb.
\)
Its peak prover memory, including all precomputed data, is
\(
    O\bigl(N+\kappa^2t+(t-1)q^{2t}\bigr) \Field.
\)
\end{proposition}

We prove Proposition~\ref{proposition:hadamard-complexity} in
Appendix~\ref{appx:proof-hadamard}.



\subsection{Packed Permutation Proof}
\label{sec:permutation}

We next consider the permutation relation:
\[
    \bm a\stackrel{\mathsf{perm}}{=}\bm b,
    \qquad
    \bm a,\bm b\in\Field^N.
\]
It asserts that the two vectors contain the same multiset of
entries. Such proofs verify the PLONK copy constraints in
Eq.~\eqref{eq:plonk2} and are used in lookup
constructions~\cite{eurocrypt/SettyTW24}.
After the witness commitments are fixed, the verifier samples
\(\ext z\sample\ExtensionField\). The prover claims a common
fingerprint \(\ext\Theta\) and proves:
\begin{equation}
    \ext\Theta
    =\mathsf{GrandProd}(\ext z\bm1+\bm a)
    =\mathsf{GrandProd}(\ext z\bm1+\bm b).
\label{eq:packed-permutation-fingerprint}
\end{equation}
The two fingerprint polynomials are identical if and only if
the permutation relation holds. It therefore suffices to prove
the two grand-product claims. We describe the proof for \(\bm a\).

\smallskip\noindent\textbf{Grand products from Hadamard products.}
Arrange \(\bm a\) as \(\bm A\in\Field^{m\times\kappa}\), as in
Eq.~\eqref{eq:matAB}, with \(N=m\kappa\).
Let \(q=|\Field|\), and assume that \(m,\kappa\) are powers
of two with \(m\geq\kappa\geq2\).
We first construct a polynomial product tree over \(\Field\),
keeping \(\kappa\) columns and halving the number of rows at each
level:
\begin{equation}
\begin{aligned}
    \bm C_0(Z)&\coloneqq Z \cdot \bm1_{m\times\kappa}+\bm A,\\
    \bm C_{\tau+1}(Z)
      &\coloneqq\bm C_{\tau,U}(Z)\circ\bm C_{\tau,D}(Z),
      \qquad \tau \in [0,\log_2\kappa-1].
\end{aligned}
\label{eq:packed-permutation-product-tree}
\end{equation}
where \(\bm C_{\tau,U}(Z)\) and \(\bm C_{\tau,D}(Z)\) are the upper
and lower row halves of \(\bm C_\tau(Z)\). Thus:
\[
    \bm C_\tau(Z)\in
       \Field[Z]^{(m/2^\tau)\times\kappa},
    \qquad
    \deg_Z \bm C_\tau(Z)=2^\tau.
\]
The prover computes these polynomial coefficients over \(\Field\), then
evaluates \(\bm C_{\log_2\kappa}(Z)\) using precomputed powers of
\(\ext z\) and $\Field-\ExtensionField$ products.
Its \(m\) entries form the vector
\(\bm c_{\log_2\kappa}(\ext z)\in\ExtensionField^m\).
The remaining levels are computed over
\(\ExtensionField\):
\[
    \bm c_{\tau+1}(\ext z)
       =\bm c_{\tau,U}(\ext z)\circ\bm c_{\tau,D}(\ext z),
    \qquad
    \tau \in [\log_2\kappa, \log_2 N-1].
\]
where \(U,D\) now denote vector halves. The final entry is
\(\bm c_{\log_2N}(\ext z)=\ext\Theta\).

The proof traverses these trees in reverse. Starting from
\(\ext\Theta\), weighted Hadamard proofs over \(\ExtensionField\)
reduce the last \(\log_2m\) levels to an MLE claim on
\(\bm c_{\log_2\kappa}(\ext z)\).
Factoring its tensor weight into column and row weights gives:
\[
    \ext v_{\log_2\kappa}
      =\ext{\bm l}_{\log_2\kappa}^{\top}
         \bm C_{\log_2\kappa}(\ext z)
         \ext{\bm r}_{\log_2\kappa}.
\]
At each subsequent level, the recursive claim has the same form:
\[
    \ext v_{\tau+1}
       =\ext{\bm l}_{\tau+1}^{\top}
          \bm C_{\tau+1}(\ext z)\ext{\bm r}_{\tau+1},
    \qquad
    \ext{\bm l}_{\tau+1}\in\ExtensionField^{m/2^{\tau+1}},
    \quad
    \ext{\bm r}_{\tau+1}\in\ExtensionField^\kappa,
\]
with a tensor-structured row weight and fixed column weights.

\smallskip\noindent\textbf{Base-field interpolation: \(q\geq2\kappa-1\).}
In this case, the interpolation points lie in \(\Field\).
We use the Lagrange basis from Section~\ref{sec:hadamard}:
\[
    \bm\psi(X)\coloneqq
       (\psi_1(X),\ldots,\psi_\kappa(X))^\top,
    \qquad
    \bm g_\tau(Z,X)\coloneqq\bm C_\tau(Z)\bm\psi(X).
\]
Choose \(\kappa-1\) distinct points
\(\gamma'_1,\ldots,\gamma'_{\kappa-1}\in
\Field\setminus\mathcal H\).
For each level \(\tau\in [\log_2\kappa]\) and point
\(\gamma'_j\), fix \(X=\gamma'_j\) and retain \(Z\) as a
variable. Define the polynomial vector:
\[
\begin{aligned}
    \bm\eta_{\tau,j}(Z)
      &\coloneqq
      \bm g_{\tau-1,U}(Z,\gamma'_j)
      \circ\bm g_{\tau-1,D}(Z,\gamma'_j)
      =\sum_{\nu=0}^{2^\tau}\bm p_{\tau,j,\nu}Z^\nu,
    \quad
    \bm p_{\tau,j,\nu}\in\Field^{m/2^\tau}.
\end{aligned}
\]
Each entry has degree at most \(2^\tau\) in \(Z\).
In particular, since \(\sum_j\psi_j(X)=1\):
\[
    \bm g_0(Z,\gamma'_j)
       =Z \cdot \bm1_m+\bm A\bm\psi(\gamma'_j).
\]
The prover computes the coefficients
\(\bm p_{\tau,j,\nu}\) over \(\Field\). The weighted
polynomial is:
\[
    F_{\tau,j}(Z)
       \coloneqq\left\langle\ext{\bm l}_\tau,
                             \bm\eta_{\tau,j}(Z)\right\rangle
       =\sum_{\nu=0}^{2^\tau}
          \left\langle\ext{\bm l}_\tau,
                       \bm p_{\tau,j,\nu}\right\rangle Z^\nu
       \in\ExtensionField[Z].
\]
The prover first aggregates the coefficient vectors
\(\bm p_{\tau,j,\nu}\) against \(\ext{\bm l}_\tau\), obtaining
the extension-field coefficients of \(F_{\tau,j}(Z)\).
It then evaluates this scalar polynomial at \(\ext z\).
The powers \(1,\ext z,\ldots,\ext z^\kappa\) are shared across
all levels.

For the transition from \(\bm C_{\tau+1}\) to \(\bm C_\tau\),
where \(\tau \in [0,\log_2\kappa-1]\), the prover applies
Section~\ref{sec:hadamard} by sending the column claims:
\[
    \ext{\bm u}_{\tau+1}^{\top}
       \coloneqq
    \ext{\bm l}_{\tau+1}^{\top}
        \bigl(\bm C_{\tau,U}(\ext z)
               \circ\bm C_{\tau,D}(\ext z)\bigr).
\]
For each column, the prover first aggregates the coefficients
of \(\bm C_{\tau+1}(Z)\) against \(\ext{\bm l}_{\tau+1}\),
then evaluates the resulting scalar polynomial at \(\ext z\).
The verifier checks
\(\ext v_{\tau+1}=\ext{\bm u}_{\tau+1}^{\top}
\ext{\bm r}_{\tau+1}\).
As in Section~\ref{sec:hadamard}, define:
\[
\begin{aligned}
    g_C(X)&\coloneqq\ext{\bm u}_{\tau+1}^{\top}\bm\psi(X),\\
    f(X)&\coloneqq
       \left\langle\ext{\bm l}_{\tau+1},
          \bm g_{\tau,U}(\ext z,X)
             \circ\bm g_{\tau,D}(\ext z,X)\right\rangle
       -g_C(X).
\end{aligned}
\]
For valid column claims, \(f(X)\) vanishes on \(\mathcal H\).
Its off-domain samples are:
\[
    f(\gamma'_j)
       =F_{\tau+1,j}(\ext z)-g_C(\gamma'_j).
\]
Interpolating them together with the zeros on \(\mathcal H\)
gives \(f(X)\), and division by \(\Psi(X)\)) gives
\(h(X)\in\ExtensionField[X]\) with \(\deg h (X)\leq\kappa-2\).
The prover fixes \(\ext{\bm u}_{\tau+1}\) and \(h(X)\) before the
verifier samples a fresh
\(\ext x\sample\ExtensionField\setminus\mathcal H\).
The resulting check is:
\[
    \left\langle\ext{\bm l}_{\tau+1},
       \bm g_{\tau,U}(\ext z,\ext x)
         \circ\bm g_{\tau,D}(\ext z,\ext x)\right\rangle
       =g_C(\ext x)+\Psi(\ext x)h(\ext x).
\]
This is a weighted Hadamard claim of length \(m/2^{\tau+1}\)
over \(\ExtensionField\).

\smallskip\noindent\textbf{Recursive queries.}
We implement each extension-field weighted Hadamard proof by the degree-2 zero-check via
sumcheck over $\ExtensionField$. The checked polynomial has individual degree at most
three, and the two child claims use the same row MLE point
\(\ext{\bm y}\).
After their values \(\ext v_{\tau,U},\ext v_{\tau,D}\) are fixed,
the verifier samples a fresh \(\ext y_0\sample\ExtensionField\).
The next claim is obtained by setting:
\[
\begin{aligned}
    \ext v_\tau
      &\coloneqq(1-\ext y_0)\ext v_{\tau,U}
                    +\ext y_0\ext v_{\tau,D},
                    \qquad
    \ext{\bm l}_\tau
      \coloneqq\tensorspan{\ext y_0\mathbin\|\ext{\bm y}},
    \qquad
    \ext{\bm r}_\tau
      \coloneqq\ext{\bm\rho}(\ext x).
\end{aligned}
\]
It has the desired form
\(\ext v_\tau=\ext{\bm l}_\tau^\top
\bm C_\tau(\ext z)\ext{\bm r}_\tau\).
The same packed Hadamard-product proof is applied.
At level zero, both weight vectors sum to one, leaving the
claim:
\[
    \ext v_0-\ext z
       =\ext{\bm l}_0^\top\bm A\ext{\bm r}_0.
\]

\smallskip\noindent\textbf{Cost over \(\Field\).}
At level \(\tau\), forming the \(\kappa-1\) encodings from
\(\bm C_{\tau-1}\) costs \(O(N\kappa)\,\bb\).
Their row products cost \(O(m\kappa2^\tau)\,\bb\), totaling
\(O(N\kappa)\,\bb\) across all levels.
Coefficient aggregation costs \(O(N)\,\be\) per level,
including the column claims. The aggregated polynomials have
degree at most \(2^\tau\), so evaluating them at \(\ext z\)
costs \(O(\kappa2^\tau)\,\ee\) per level and
\(O(\kappa^2)\,\ee\) in total.
To form the terminal Hadamard vectors, the prover first
contracts the coefficients of \(\bm C_\tau(Z)\) with
\(\ext{\bm\rho}(\ext x)\), then evaluates the resulting row
polynomials at \(\ext z\), at a cost of
\(O(N)\,\be+O(m)\,\ee\) per level.
Using shared powers of \(\ext x\), interpolation and public
query generation cost \(O(\kappa^3)\,\bb\) per level.
The short Hadamard proofs and extension-field tree have
total length \(O(m)\).
Thus, under schoolbook arithmetic, the prover cost for one
grand-product proof is
\(
    O(N\kappa\log_2\kappa+\kappa^4)\bb.
\)


\smallskip\noindent\textbf{Intermediate-field tables over smaller fields: \(q<2\kappa-1\).}
Choose \(t\) and $\kappa$ such that $q^t \ge 2\kappa-1$ and $t \mid \kappa$, and we can use the tower:
\[
    \Field\subseteq\mathbb K=\mathbb F_{q^t}
       \subseteq\ExtensionField,
    \qquad [\ExtensionField:\mathbb K]=\kappa/t.
\]
We choose \(\mathcal H\) and the additional interpolation nodes
in \(\mathbb K\). The coefficients of \(\bm C_\tau(X)\) remain
in \(\Field\), while those of \(\bm g_\tau(X)\) and
\(\bm\eta_{\tau,j}(X)\) lie in \(\mathbb K\).
The polynomial identities and recursive queries are unchanged.

We precompute a multiplication table over \(\mathbb K\)
and count each table lookup as an \(O(1)\)-time operation.
Representing an extension-field element by \(\kappa/t\)
coordinates over \(\mathbb K\), a mixed
\(\mathbb K\times\ExtensionField\) product costs \(O(\kappa)\).
Using this multiplication table and precomputed Lagrange basis
evaluations, the grand-product reduction over \(\mathbb K\)
similarly achieves an online prover cost of
\(
    O\!\left(N\kappa\log_2\kappa+\frac{\kappa^4}{t}\right)\,\bb.
\)

Algorithm~\ref{alg:permutation} in Appendix~\ref{appx:pseudocode}
specifies the protocol for both cases.

\begin{proposition}[Packed permutation proof]
\label{proposition:permutation-complexity}
Let \(q=|\Field|\) be fixed and let
\(N=m\kappa\), where \(m,\kappa\) are powers of two and
\(m\geq\kappa\geq2\).
Algorithm~\ref{alg:permutation} is a proof protocol for the grand-product relation
with the soundness error:
\[
\begin{aligned}
    \epsilon_{\mathsf{GP}}\leq{}&
    \frac{\frac32\log_2m(\log_2m-1)+\log_2m}
         {|\ExtensionField|}+\sum_{\tau=0}^{\log_2\kappa-1}
      \left(
        \frac{2\kappa-2}{|\ExtensionField|-\kappa}
        +\frac{3\log_2(m/2^{\tau+1})+1}{|\ExtensionField|}
      \right).
\end{aligned}
\]
Consequently, when both input vectors are committed, the
permutation argument has soundness error at most:
\[
    \epsilon_{\mathsf{Perm}}\leq{} \min\left\{1,\frac{N-1}{|\ExtensionField|}\right\}
       +2\epsilon_{\mathsf{GP}}+\epsilon_{\mathsf{PCS}},
\]
where \(\epsilon_{\mathsf{PCS}}\) is the joint soundness error
of authenticating the two terminal packed-opening claims.
With the intermediate-field multiplication table and block
The online prover cost for one
grand product is
\(
    O\!\left(N\kappa\log_2\kappa+\kappa^4\right)\,\bb.
\)
The peak prover memory, including all precomputed tables, is:
\(
    O(N+\kappa^3+tq^{2t})\Field
\)
\end{proposition}

Proof of this proposition is provided in Appendix~\ref{appx:proof-permutation}.

\subsection{Packed Polynomial Commitment Scheme}
\label{sec:pcs}

We complete the packed reductions by authenticating their
terminal claims against a fixed-basis commitment.
Let \(N=m\kappa\) and \(\bm A\in\Field^{m\times\kappa}\).
The preceding protocols produce claims of the form:
\[
    \ext v
    =\langle\bm A\ext{\bm\rho},\ext{\bm\xi}\rangle
    =\ext{\bm\xi}^{\top}\bm A\ext{\bm\rho},
    \qquad
    \ext{\bm\rho}\in\ExtensionField^\kappa,
    \quad
    \ext{\bm\xi}\in\ExtensionField^m,
\]
where \(\ext{\bm\xi}\) is a tensor-structured row weight.
The query weights \(\ext{\bm\rho}\) may differ from the fixed
\(\Field\)-basis \(\ext{\bm\beta}\in\ExtensionField^\kappa\)
used to commit to the packed vector:
\[
    \ext{\bm a}=\bm A\ext{\bm\beta}\in\ExtensionField^m.
\]

To relate the two packings, the prover sends the column evaluations
\(\ext{\bm u}^{\top}\coloneqq\ext{\bm\xi}^{\top}\bm A\).
The verifier checks
\(\ext v\EqOrNot\langle\ext{\bm u},\ext{\bm\rho}\rangle\).
Write the query and the column claims in base-field coordinates:
\[
    \ext{\bm\xi}=\bm\Xi\ext{\bm\beta},
    \qquad
    \ext{\bm u}^{\top}=\ext{\bm\beta}^{\top}\bm U,
    \qquad
    \bm\Xi\in\Field^{m\times\kappa},
    \quad
    \bm U\in\Field^{\kappa\times\kappa}.
\]
Since \(\bm A\) has base-field entries, the column claims are
equivalent to \(\bm U=\bm\Xi^\top\bm A\).
Packing this matrix identity with the fixed basis gives:
\[
    \bm U\ext{\bm\beta}
    =\bm\Xi^\top\bm A\ext{\bm\beta}
    =\bm\Xi^\top\ext{\bm a}\in\ExtensionField^\kappa.
\]
After \(\ext{\bm u}\) is fixed, a fresh random projection reduces
this vector identity to a length-\(m\) inner-product claim involving
the committed vector \(\ext{\bm a}\).
Fact~\ref{fact:ip} then leaves one MLE opening of
\(\ext{\bm a}\) and one public query evaluation.
The verifier computes the latter directly from the tensor factors
of \(\ext{\bm\xi}\).

This is the ring-switching technique~\cite{eurocrypt/DiamondP26}
expressed in our vector and matrix notation.
Appendix~\ref{appx:pcs} gives the construction and its analysis;
Algorithm~\ref{alg:packed-pcs} specifies the protocol.
Under schoolbook arithmetic and \(m\geq\kappa\), the packed
opening reduction uses \(O(\kappa N)\,\bb\) and
\(O(N)\Field\) memory,
given an underlying extension-field PCS that commits and opens in \(O(m)\ee\) and uses \(O(m)\ExtensionField\) of memory.

\section{Putting Everything Together}
\label{sec:together}

To construct an end-to-end bitwidth-aware SNARK for circuit
satisfiability (Circ-SAT), we combine the building blocks
of Section~\ref{sec:blocks} to prove the PLONK constraints
of Definition~\ref{def:plonk}.
The same approach can be adapted to other constraint systems,
such as R1CS.
In addition to witness values, the constraint system encodes
wire positions to certify the wiring of the circuit.
Bitwidth awareness must therefore account for both the
\(\lceil\log_2 q\rceil\)-bit witness entries and the
\(\lceil\log_2 N\rceil\)-bit position labels.
We seek prover time and memory usage that scale linearly
with this \emph{effective bitwidth} \(\lceil\log_2 q\rceil + \lceil\log_2 N\rceil\).
For large circuits, processing position labels can introduce
an additional \(\log N\) factor in the corresponding costs.
When packing multiple executions of the same circuit,
the shared wiring allows its encoding and the associated
proof work to be amortized across executions.

\subsection{Proof for the Copy Constraint}
\label{sec:plonk-copy}

We reduce the PLONK copy constraint
\(\sigma(\bm w)=\bm w\) in Eq.~\eqref{eq:plonk2} to a
permutation relation over an intermediate field.
Let \(\bm w\in\Field^L\), where
\(L=3N+n_{\mathsf{io}}\), and let \(q=|\Field|\).
The reduction encodes each witness value together with its
position.

\smallskip\noindent\textbf{Joint encoding of values and positions.}
Choose \(\underline q=2^s\leq q\), with
\(s=\Theta(\log_2q)\), and elements
\(\delta_1,\ldots,\delta_s\in\Field\) such that the map
\(\sum_{t=1}^s2^{t-1}b_t\mapsto\sum_{t=1}^s\delta_tb_t\)
is injective on \(\{0,1\}^s\).
For prime \(q\), take \(s=\lfloor\log_2q\rfloor\) and
\(\delta_t=2^{t-1}\).
Appendix~\ref{appx:mleiota} gives a choice for every finite field.
We reserve \(\lceil\log_{\underline q} L\rceil\) base-field
coordinates for the position and one for the witness value.
Choose \(d\) as the least power of two at least
\(\lceil\log_{\underline q}L\rceil+1\), and choose a
power-of-two extension degree \(\kappa\) with \(2d\leq\kappa\).
This gives the tower:
\[
    \Field\subset\mathbb K\subset\ExtensionField,
    \quad
    [\mathbb K:\Field]=d,
    \quad
    d\mid\kappa,
    \quad
    [\ExtensionField:\mathbb K]=\kappa_{\mathbb K}
       =\kappa/d.
\]
We fix an \(\Field\)-basis
\(\bm\theta=(\theta_1,\ldots,\theta_d)^\top\in\mathbb K^d\)
of \(\mathbb K\).

Then, we can represent \(\sigma\) by
\(\bm\sigma=(\sigma(1),\ldots,\sigma(L))^\top\), and let
\(\bm i=(1,\ldots,L)^\top\).
Write their zero-based positions in radix \(\underline{q}\):
\[
\begin{aligned}
    \bm\sigma-\bm1_L
       &=\sum_{j=1}^{\lceil\log_{\underline{q}} L\rceil}\underline q^{j-1}\bm\sigma_{(j)},
       \quad
    \bm i-\bm1_L
       =\sum_{j=1}^{\lceil\log_{\underline{q}} L\rceil}\underline q^{j-1}\bm i_{(j)},
    \quad
    \bm\sigma_{(j)},\bm i_{(j)}
       \in [0,\underline{q}-1]^L.
\end{aligned}
\]
These are integer decompositions. In field expressions, we
interpret each digit through the injection above.
Consequently,
we encode the permutation images and identity positions as:
\[
    \bm\varsigma
       \coloneqq\sum_{j=1}^{\lceil\log_{\underline q} L\rceil}\theta_j\bm\sigma_{(j)} \in \mathbb{K}^L,
    \qquad
    \bm\iota
       \coloneqq\sum_{j=1}^{\lceil\log_{\underline q} L\rceil}\theta_j\bm i_{(j)}
       \in\mathbb K^L.
\]
The labels \(\iota_1,\ldots,\iota_L\) are distinct, and
\(\varsigma_i=\iota_{\sigma(i)}\).
The remaining basis coordinate binds each label to a witness
value. We define:
\[
    \bm a\coloneqq\bm\varsigma+\theta_d \bm w  \in \mathbb{K}^L,
    \qquad
    \bm b\coloneqq\bm\iota+\theta_d\bm w
       \in\mathbb K^L.
\]

\begin{lemma}[Copy constraints from a permutation relation]
\label{eq:permutation}
For the encodings above and every \(\bm w\in\Field^L\):
\[
    \sigma(\bm w)=\bm w
    \quad\Longleftrightarrow\quad
    \bm a\stackrel{\mathsf{perm}}{=}\bm b.
\]
\end{lemma}

The basis coordinates make the joint encoding injective.
Consequently, the label in \(a_i\) can match only the label in
\(b_{\sigma(i)}\), and equality of these entries enforces
\(w_i=w_{\sigma(i)}\).
Appendix~\ref{appx:proof-copy-constraint} gives the proof.

\smallskip\noindent\textbf{Applying the packed permutation proof.}
With the witness commitments fixed, both operands are determined
by \(\bm w\) and the public permutation.
We treat \(\mathbb K\) as the base field and apply
Algorithm~\ref{alg:permutation} over \(\ExtensionField\)
with packing width \(\kappa_{\mathbb K}\), chosen so that
\(|\mathbb K|\geq2\kappa_{\mathbb K}-1\).
All required interpolation points can therefore be chosen
in \(\mathbb K\).
The permutation proof leaves weighted evaluations with arbitrary
column weights and tensor row weights on \(\bm a\) and \(\bm b\).
By linearity, these reduce to evaluations of
\(\bm\varsigma,\bm\iota\in\mathbb K^L\) and
\(\bm w\in\Field^L\).
The verifier evaluates \(\bm\iota\) directly using
\(O(\kappa_{\mathbb K}+\log L)\) extension-field operations,
as shown in Appendix~\ref{appx:proof-position-encoding}; ordinary
MLE queries need only \(O(\log L)\) operations.
The remaining claims are verified using the packed PCS.

\subsection{Packing for a Single Circ-SAT Instance}
\label{sec:plonk-packing}

With the copy constraints handled by the permutation proof,
it remains to prove the PLONK gate constraints in Eq.~\eqref{eq:plonk1}.
Let \(N=m\kappa\), with \(m,\kappa\) powers of two, and
reuse the polynomial basis of Section~\ref{sec:pcs}, with
\(\ext\beta_1=1\), for the wire commitments.
The packed Hadamard proof of Section~\ref{sec:hadamard}
handles the length-\(N\) gate equation.
For the copy proof, pad both operands with zeros to length
\(4N\) and sample its fingerprint challenge from
\(\ExtensionField\setminus\{0\}\).
The common padding factors then cancel.

We reduce the gate constraints to weighted Hadamard claims by
evaluating the residual's multilinear extension. The verifier samples
\(\ext{\bm x}\sample\ExtensionField^{\log_2N}\) and sets
\(\ext{\bm\xi}=\tensorspan{\ext{\bm x}}\).
It checks one random projection of the gate residual:
\[
    \left\langle \ext{\bm\xi},
    \bm q_L\circ\bm w_L+\bm q_R\circ\bm w_R
    +\bm q_M\circ(\bm w_L\circ\bm w_R)
    +\bm q_O\circ\bm w_O+\bm q_C\right\rangle=0.
\]
Define the virtual vector \(\bm w_M=\bm w_L\circ\bm w_R\).
The prover sends
\(\ext c_L,\ext c_R,\ext c_M,\ext c_O,\ext c_C\);
the verifier checks their sum and the following claims:
\begin{align}
    &\ext c_L+\ext c_R+\ext c_M+\ext c_O+\ext c_C=0,
    \nonumber\\
    &\ext c_L=\langle\ext{\bm\xi},\bm q_L\circ\bm w_L\rangle,
    \quad
    \ext c_R=\langle\ext{\bm\xi},\bm q_R\circ\bm w_R\rangle,
    \quad
    \ext c_O=\langle\ext{\bm\xi},\bm q_O\circ\bm w_O\rangle,
    \label{eq:plonkWeightedHadamard1}\\
    &\ext c_M=\langle\ext{\bm\xi},\bm q_M\circ\bm w_M\rangle,
    \quad
    \ext c_C=\langle\ext{\bm\xi},\bm q_C\rangle.
    \label{eq:plonkWeightedHadamard2}
\end{align}
Each of the three products in
Eq.~\eqref{eq:plonkWeightedHadamard1} has the form of
Eq.~\eqref{eq:hadamard-claim}. Algorithm~\ref{alg:hadamard}
reduces each to a selector evaluation and an evaluation of a committed
wire vector.

For \(\ext c_M\), the same reduction leaves an evaluation of
\(\bm q_M\) and an evaluation of \(\bm w_M\). The latter is
itself a weighted Hadamard claim on \(\bm w_L,\bm w_R\), which
we prove with a second invocation of Algorithm~\ref{alg:hadamard}.

We batch the terminal claims into a witness group and a public
group, using the packed opening reduction of Section~\ref{sec:pcs}.
The prover commits to the concatenation of the three packed wire
vectors before any challenge; the selector and position-table
commitment is preprocessed, with a total packed length of \(m_{\mathsf{pub}}=
(5+3\lceil\log_{\underline q}(3N+n_{\mathsf{io}})\rceil)m\).

\begin{theorem}[Single-instance packing]
\label{theorem:single-instance-packing}
Assume the parameters above, \(m\geq\max\{\kappa,q,n_{\mathsf{io}}\}\),
and the PCS assumptions of Proposition~\ref{proposition:packed-pcs}.
The single-instance protocol is complete, with soundness error
at most:
\begin{equation}
    \epsilon_\mathsf{single}
       =O\left(\frac{N}{|\ExtensionField|}\right) + \epsilon_{\mathsf{PCS}},
\label{eq:together-soundness}
\end{equation}
where \(\epsilon_{\mathsf{PCS}}\) bounds the joint binding and
opening error of the underlying PCS calls.
With circuit preprocessing, the efficiency bounds are:
\begin{equation}
\begin{aligned}
    \text{Prover time}\quad
       &O\left(N\kappa\log_q N\log_2\kappa\right)\,\bb
          +O(\kappa^4)
          +\mathcal{TP}_{\mathsf{pcs}}(3m+m_{\mathsf{pub}}),\\
    \text{Prover memory}\quad
       &O(N\log_q N)\,\Field+O(\kappa^4+n_{\mathsf{io}})
          +\mathcal{SP}_{\mathsf{pcs}}(3m+m_{\mathsf{pub}}),\\
    \text{Verifier time}\quad
       &O\left(
          \kappa\log_q N\log_2N+(\log_2N)^2
          +n_{\mathsf{io}}
          \right)\,\ee
          +\mathcal{TV}_{\mathsf{pcs}}(3m+m_{\mathsf{pub}}),\\
    \text{Proof size}\quad
       &O\left(\kappa\log_q N
          +\kappa\log_2\kappa
          +(\log_2N)^2\right)\,\ExtensionField
          +\mathcal{S}_{\mathsf{pcs}}(3m+m_{\mathsf{pub}}),
\end{aligned}
\label{eq:together-single-efficiency}
\end{equation}
where \(\mathcal{TP}_{\mathsf{pcs}}\),
\(\mathcal{SP}_{\mathsf{pcs}}\),
\(\mathcal{TV}_{\mathsf{pcs}}\), and
\(\mathcal{S}_{\mathsf{pcs}}\) denote the PCS
prover time, prover memory, verifier time, and proof size for
total logical length \(3m+m_\mathsf{pub}\), respectively,
including domain padding. The unlabeled remainders count
base-field operations and base-field elements.
\end{theorem}

The proof is deferred to Appendix~\ref{appx:proof-single-instance}.
Algorithm~\ref{alg:single-instance-packing} gives the protocol.

\subsection{Packing Multiple Circ-SAT Instances of the Same Circuit}
\label{sec:multi-instance-packing}

The $\log_q N$ term in
Theorem~\ref{theorem:single-instance-packing} introduced by the
position-encoding overhead can be amortized by packing executions of the same circuit.
Since all instances share the same selectors and wiring
permutation, the corresponding linear operations commute
with packing.
This yields amortized per-instance prover time
\(O(N\kappa)\,\bb\) and memory \(O(N)\,\Field\), excluding
the underlying PCS.

Consider \(\kappa\) executions of the same \(N\)-gate circuit,
where \(N,\kappa\) are powers of two,
\(\kappa\geq2\), and \(N\geq\max\{\kappa,n_{\mathsf{io}}\}\).
Let \(L=3N+n_{\mathsf{io}}\).
For each \(j\in[\kappa]\), let
\(\bm w^{(j)}\in\Field^{3N+n_{\mathsf{io}}}\)
be the witness vector from Definition~\ref{def:plonk}.
Using the polynomial basis \(\ext{\bm\beta}\) of
Section~\ref{sec:pcs}, with \(\ext\beta_1=1\), define:
\[
    \bm W
       =\concat{\bm w^{(1)},\ldots,\bm w^{(\kappa)}}
       \in\Field^{L\times\kappa},
    \qquad
    \ext{\bm w}=\bm W\ext{\bm\beta}\in\ExtensionField^L.
\]
For \(\ell\in\{L,R,O\}\), arrange the corresponding wire
vectors as \(\bm W_\ell\in\Field^{N\times\kappa}\).
Their packings and the virtual product vector are:
\begin{equation}
    \ext{\bm w}_\ell=\bm W_\ell\ext{\bm\beta}
       \in\ExtensionField^N,
    \qquad
    \ext{\bm w}_M=(\bm W_L\circ\bm W_R)\ext{\bm\beta}.
\label{eq:together-multiple-wires}
\end{equation}
The prover commits to the concatenation of
\(\ext{\bm w}_L,\ext{\bm w}_R,\ext{\bm w}_O\), of logical
length \(3N\), before any challenge.
The product vector remains uncommitted.

\begin{lemma}[Packed PLONK constraints]
\label{lemma:packedPlonk}
The \(\kappa\) instances satisfy the PLONK constraints if and only if:
\begin{align}
    \bm q_L\circ\ext{\bm w}_L
    +\bm q_R\circ\ext{\bm w}_R
    +\bm q_M\circ\ext{\bm w}_M
    +\bm q_O\circ\ext{\bm w}_O
    +\left(\sum_{j=1}^{\kappa}\ext\beta_j\right)\bm q_C
       &=\bm0_N,
       \label{eq:together-multiple-gates}\\
    \sigma(\ext{\bm w})&=\ext{\bm w}.
       \label{eq:packed-plonk2}
\end{align}
\end{lemma}

The shared selectors and permutation are \(\Field\)-linear,
and the basis packing is injective.
Thus each packed equality holds if and only if it holds in
every column of \(\bm W\).

\smallskip\noindent\textbf{Gate constraints.}
The verifier samples
\(\ext{\bm x}\sample\ExtensionField^{\log_2N}\)
and uses sumcheck to verify that the
\(\tensorspan{\ext{\bm x}}\)-weighted residual sum in
Eq.~\eqref{eq:together-multiple-gates} is zero.
The terminal check requires five selector MLE claims
and four packed-wire MLE evaluations at a common point
\(\ext{\bm y} \in \ExtensionField^{\log_2 N}\).
These claims are authenticated by the batched openings below.

To authenticate the packed product MLE value \(\ext v_M\),
the prover proves:
\begin{equation}
    \ext v_M
    =
    \tensorspan{\ext{\bm y}}^\top
    (\bm W_L\circ\bm W_R)\ext{\bm\beta}.
\label{eq:together-multiple-product-opening}
\end{equation}
The packed Hadamard-product proof in Algorithm~\ref{alg:hadamard}
reduces this to two weighted evaluation claims on
\(\bm W_L\) and \(\bm W_R\).
The column-projection step of Section~\ref{sec:pcs} converts
them to inner-product claims on the committed wire packings,
which are batched with the ordinary wire queries below.

\smallskip\noindent\textbf{Copy constraints.}
We apply the position encoding of Section~\ref{sec:plonk-copy}
in the basis \(\ext{\bm\beta}\) to obtain
\(\bm\iota,\bm\varsigma\in\ExtensionField^L\).
Appendix~\ref{appx:mleiota} evaluates the identity-label MLE
in \(O(\log L)\) extension-field operations.
The labels \(\iota_i\) are distinct and
\(\varsigma_i=\iota_{\sigma(i)}\).
For a fresh challenge \(\ext y\sample\ExtensionField\),
the copy constraint in Eq.~\eqref{eq:packed-plonk2} reduces
to a permutation check over $\ExtensionField$:
\[
    \bm\iota+\ext y\,\ext{\bm w}
       \stackrel{\mathsf{perm}}{=}
    \bm\varsigma+\ext y\,\ext{\bm w} \in \ExtensionField^L.
\]
The verifier samples an independent
\(\ext z\sample\ExtensionField\), and the prover uses two
extension-field product trees to establish:
\[
    \prod_{i=1}^{L}(\ext z+\iota_i+\ext y\,\ext w_i)
       =
    \prod_{i=1}^{L}(\ext z+\varsigma_i+\ext y\,\ext w_i).
\]
Each tree is padded with ones to length \(4N\).
The two reductions leave MLE claims on \(\bm\varsigma\),
\(\bm\iota\), and \(\ext{\bm w}\).
The verifier evaluates \(\bm\iota\) directly and computes the
public-I/O and padding contributions without commitments.
All remaining claims are batched as in
Section~\ref{sec:plonk-packing}, separately for the witness
and public vectors, with one underlying PCS opening per group.
When \(5+3\lceil\log_{\underline q}L\rceil\leq\kappa\),
the five selectors and the position-digit columns for the three
wire segments fit into one \(N\times\kappa\) base-field matrix.
Packing its rows gives one preprocessed public commitment of
length \(N\), for a total logical PCS length of \(4N\).

\begin{theorem}[Multiple-instance packing]
\label{theorem:multiple-instance-packing}
Assume the parameters above, \(q\geq2\kappa-1\),
\(5+3\lceil\log_{\underline q}L\rceil\leq\kappa\), and the PCS
assumptions of Proposition~\ref{proposition:packed-pcs}.
The protocol for \(\kappa\) instances of the same circuit
is complete, with soundness error at most:
\begin{equation}
    \epsilon_{\mathsf{multi}}
       =O\left(\frac{N}{|\ExtensionField|}\right)
          +\epsilon_{\mathsf{PCS}},
\label{eq:together-multiple-soundness}
\end{equation}
where \(\epsilon_{\mathsf{PCS}}\) bounds the joint binding and
opening error of the underlying PCS calls.
With circuit preprocessing, the total costs for the batch are:
\begin{equation}
\begin{aligned}
    \text{Prover time}\quad
       &O(N\kappa^2)\,\bb+\mathcal{TP}_{\mathsf{pcs}}(4N),\\
    \text{Prover memory}\quad
       &O(N\kappa)\,\Field+\mathcal{SP}_{\mathsf{pcs}}(4N),\\
    \text{Verifier time}\quad
       &O\left(\kappa\log_2N+(\log_2N)^2
          +n_{\mathsf{io}}\right)\,\ee+\mathcal{TV}_{\mathsf{pcs}}(4N),\\
    \text{Proof size}\quad
       &O\left(\kappa+(\log_2N)^2\right)\,\ExtensionField+\mathcal{S}_{\mathsf{pcs}}(4N),
\end{aligned}
\label{eq:together-multiple-efficiency}
\end{equation}
where \(\mathcal{TP}_{\mathsf{pcs}}\),
\(\mathcal{SP}_{\mathsf{pcs}}\),
\(\mathcal{TV}_{\mathsf{pcs}}\), and
\(\mathcal{S}_{\mathsf{pcs}}\) denote the aggregate PCS prover
time, prover memory, verifier time, and proof size for total
logical length \(4N\), respectively.

Excluding the underlying PCS, dividing the prover time and peak
memory by \(\kappa\)
gives amortized costs of \(O(N\kappa)\bb\) and
\(O(N)\Field\) per instance.
\end{theorem}

The proof is deferred to
Appendix~\ref{appx:proof-multiple-instance}.
Algorithm~\ref{alg:multiple-instance-packing} gives the protocol.

A PCS whose prover time and memory are linear in its
extension-field input length therefore gives an end-to-end
bitwidth-aware construction.

\iffalse
\subsection{Binary-Field Instantiation}
\label{sec:binary-packing}

For binary circuits, we use a small binary extension as the working
field of the entire PIOP. By contrast, the intermediate-field
Hadamard variant in Section~\ref{sec:hadamard} retains the
original base-field witness and its fixed-basis commitment,
using the intermediate field only for certificate construction.
For compatible \(t\mid\kappa\) and \(2^t\geq2\kappa-1\),
its tabulated implementation has \(O(N\kappa)\) online
base-field cost when \(N\geq\kappa^2\). The storage includes
\(O(\kappa^2t+t2^{2t})\) precomputed base-field coordinates.
The working-field instantiation here supplies interpolation nodes
and represents the witness in byte-field coordinates. Choose:
\[
    \mathbb F_2\subset\mathbb K=\mathbb F_{2^8}
       \subset\ExtensionField,
    \qquad [\mathbb K:\mathbb F_2]=8,
    \qquad [\ExtensionField:\mathbb K]=\kappa_{\mathbb K}.
\]
Arithmetize the circuit over \(\mathbb K\), including constraints
\(w^2-w=0\) for the binary wires, and define \(m,N\) for this
arithmetization. Instantiate the building blocks, packed PCS, and
multiple-instance composition with
\((\Field,q,\kappa)=(\mathbb K,256,\kappa_{\mathbb K})\).
Here \(N=m\kappa_{\mathbb K}\) counts working-field entries.

\smallskip\noindent\textbf{Packing and openings.}
Assume \(256\geq2\kappa_{\mathbb K}-1\), so \(\mathbb K\)
contains the interpolation nodes and off-domain samples needed
by Section~\ref{sec:hadamard}. Represent \(\ExtensionField\) in a
\(\mathbb K\)-basis. All coordinate matrices in
Section~\ref{sec:pcs} use this basis. The verifier checks public
query contractions using Eq.~\eqref{eq:pcs-query-recurrence},
while the gate and copy challenges lie in \(\ExtensionField\).

\smallskip\noindent\textbf{Position encoding.}
Set \(s=8\) and \(\underline q=256\), and assume
\(\lceil\log_{256}L\rceil\leq\kappa_{\mathbb K}\).
Choose an \(\mathbb F_2\)-basis
\((\delta_1,\ldots,\delta_8)\) of \(\mathbb K\) and encode each
eight-bit position block as \(\sum_{t=1}^8\delta_tb_t\).
These digits give injective labels in \(\ExtensionField\) through
Eq.~\eqref{eq:together-multiple-labels}.
The proof in Appendix~\ref{appx:proof-position-encoding} applies
to these coefficients and gives the same MLE identity,
Eq.~\eqref{eq:position-encoding-evaluation}.
Thus the randomized copy argument operates in the stated tower.
The role of \(\mathbb K\) here is to support the PIOP arithmetic.
Using \(\mathbb F_{2^8}\) instead for the joint value-and-position
encoding of Section~\ref{sec:plonk-packing} over \(\mathbb F_2\)
would impose \(1+\lceil\log_2L\rceil\leq8\), or \(L\leq2^7\).

\smallskip\noindent\textbf{Using a binary linear code.}
Appendix~\ref{appx:pcs} gives the code-switching
instantiation of the PCS. A binary generator matrix extends
first to \(\mathbb K\) and then to \(\ExtensionField\) without
changing its entries. The source encoding can therefore be
computed on the eight bit-coordinate matrices over \(\mathbb F_2\),
while the packed opening reduction uses the
\(\ExtensionField/\mathbb K\) basis.
Code switching supplies the terminal extension-field opening
through a characteristic-two-compatible BaseFold
instantiation~\cite{eurocrypt/DiamondP26}.

\smallskip\noindent\textbf{Efficiency and soundness.}
The proof of Theorem~\ref{theorem:multiple-instance-packing} uses
field identities, injective labels, and polynomial degree bounds,
which apply in characteristic two. With the packed PCS of
Section~\ref{sec:pcs}, its efficiency and soundness bounds hold with
\(q=256\) and \(\kappa=\kappa_{\mathbb K}\), subject to the
stated field-size and packing conditions.
Each packed opening has error at most
\(\epsilon_{\mathsf{PCS}}+(1+2\log_2m)/|\ExtensionField|\).
Arithmetic costs are measured in \(\mathbb K\)-operations,
and communication and space in \(\mathbb K\)-coordinates.
For the code-based instantiation, substitute the underlying PCS
error and costs of the chosen Blaze construction as specified in
Appendix~\ref{appx:pcs}. The packed-opening, gate,
Hadamard, and copy-reduction bounds apply with the same parameters.
\fi

\section{Experimental Evaluation}
\label{sec:experiment}

We compare the inner- and Hadamard-product protocols based on
vector polynomial packing with the small-field-first (SFF)
baseline at \(N=2^{30}\).
Table~\ref{table:component} reports prover time, verifier time,
message payload, and resident set size (RSS).
The four configurations use base fields stored in 8-, 16-, 32-,
and 64-bit words, with packing factors \(\kappa=16,8,4,2\),
respectively.
These configurations satisfy \(|\Field|\geq\kappa\) and evaluate
the base-field implementations, not the intermediate-field
tabulated variant of Section~\ref{sec:hadamard}.

\paragraph{Measurement scope.}
The benchmarks measure the PIOP components, including packing and
witness-dependent preprocessing. Commitments, PCS openings, and
deferred source-oracle evaluations are outside the measured scope.
All runs use an Intel Xeon Gold 6342 with AVX-512 and 64 Rayon
workers. RSS is measured in isolated runs separately from timing
and includes the inputs resident during the prover window.
The table notes specify the run counts, warmup, and execution-order
controls for each configuration. In particular, the 8- and 16-bit
timings and the updated 32-bit Hadamard timings are single-run
measurements.

\paragraph{Prover time and memory.}
Packing gives its largest measured benefit in the 8-bit
configuration: the inner-product prover is \(5.60\times\) faster
and the Hadamard-product prover is \(4.31\times\) faster than SFF.
Both use \(4.02\) GiB of measured resident memory, a reduction by a
factor of approximately \(4.48\).
The prover speedups decrease to \(3.54\times\) and \(2.82\times\)
at 16 bits, and to \(2.07\times\) and \(1.82\times\) at 32 bits.
The updated 32-bit Hadamard implementation reduces RSS from
\(23.99\) GiB to \(16.00\) GiB; RSS has not been remeasured for
the updated 32-bit inner-product comparison.
At 64 bits, packed prover time is slightly higher for both relations
and RSS is essentially unchanged.

\paragraph{Verifier time and communication.}
The packed Hadamard verifier is faster in all four configurations,
with measured speedups from \(1.57\times\) to \(2.41\times\).
For inner products, verification improves at 8 and 16 bits and is
slightly slower at 32 and 64 bits.
Inner-product payloads decrease in every configuration.
Hadamard packing trades additional auxiliary messages for a shorter
inner-product proof: its 8-bit payload increases from \(1{,}472\)
to \(1{,}872\) bytes, while the other three configurations have
smaller payloads.
These measurements quantify the reduction-layer tradeoffs;
the cost of a composed circuit proof also depends on the PCS.

% Sources: frozen consolidated report, with current 32-bit IP and Hadamard updates.
% https://github.com/mirandaprivate/mikros/blob/eccf9d72dbe481f9e8707d12aadc45ebab60971a/mikros/compare_small_field.md
% https://github.com/mirandaprivate/mikros/blob/eccf9d72dbe481f9e8707d12aadc45ebab60971a/mikros/IP_STABILITY_N30.md
% Selected scale: N = 2^30 for all fields and both relations.
% Historical values: unrounded four-field-comparison-v1/all_results.csv.
% Historical CSV SHA-256: adbd61cc7b165f50c9a3443ee53bd4ded92fafe53ea021858b97525a224c5b76
% 32-bit IP: ip-stability-n30-v3/report/all_methods.csv, native vs joint (new).
% Updated CSV SHA-256: d5bd74c07b0965f02ec4e72ccb7daffd1982310905ee2e5faedf36f76a4481c9
% 32-bit Hadamard: mikros/M31_HADAMARD_RSS_DIAGNOSIS.md, public WIDE_V1.
% This report was locally staged, not yet on remote main, when incorporated.
% Report SHA-256: 04e519d70ed4a9723fe0f5935d627c1136dcae09091a213a161c5a4e68290515
% Data: m31-hadamard-diagnostic-v1/report/all_results.csv (separate timing/RSS runs).
% CSV SHA-256: f65cae711beee61fbe6c9510a67fa33318a91bdfb6f80cff8658754cca9e701a
% Repeated-run timing ratios divide unrounded per-method medians, not paired ratios.
% No new 32-bit IP RSS; 8/16/64-bit results retain historical data.
\begin{table*}[!t]
\centering
\setlength{\tabcolsep}{1.5pt}
\begin{threeparttable}
\caption{Packed and small-field-first PIOP components at \(N=2^{30}\).}
\label{table:component}
\small
\scriptsize
\begin{tabular*}{\textwidth}{@{\extracolsep{\fill}}lrrrrrrrrrrrr@{}}
\toprule
& \multicolumn{3}{c}{8-bit}
& \multicolumn{3}{c}{16-bit}
& \multicolumn{3}{c}{32-bit}
& \multicolumn{3}{c}{64-bit} \\
\cmidrule(lr){2-4}\cmidrule(lr){5-7}
\cmidrule(lr){8-10}\cmidrule(l){11-13}
& \multicolumn{1}{c}{SFF}
& \multicolumn{1}{c}{Packed}
& \multicolumn{1}{c}{$\times$}
& \multicolumn{1}{c}{SFF}
& \multicolumn{1}{c}{Packed}
& \multicolumn{1}{c}{$\times$}
& \multicolumn{1}{c}{SFF}
& \multicolumn{1}{c}{Packed}
& \multicolumn{1}{c}{$\times$}
& \multicolumn{1}{c}{SFF}
& \multicolumn{1}{c}{Packed}
& \multicolumn{1}{c}{$\times$} \\
\midrule
\multicolumn{13}{l}{\textbf{Inner product (IP)}} \\
Prover (s)
  & $3.036$ & $0.542$ & \textcolor{black!80}{\textit{5.60}\(\times\)}
  & $2.829$ & $0.799$ & \textcolor{black!80}{\textit{3.54}\(\times\)}
  & $4.198$ & $2.030$ & \textcolor{black!80}{\textit{2.07}\(\times\)}
  & $3.003$ & $3.145$ & \textcolor{black!80}{\textit{0.96}\(\times\)} \\
Verifier ($\mu$s)
  & $137.64$ & $108.68$ & \textcolor{black!80}{\textit{1.27}\(\times\)}
  & $76.80$ & $51.04$ & \textcolor{black!80}{\textit{1.50}\(\times\)}
  & $21.63$ & $22.49$ & \textcolor{black!80}{\textit{0.96}\(\times\)}
  & $30.15$ & $31.94$ & \textcolor{black!80}{\textit{0.94}\(\times\)} \\
Payload (B)
  & $992$ & $894$ & \textcolor{black!80}{\textit{1.11}\(\times\)}
  & $992$ & $924$ & \textcolor{black!80}{\textit{1.07}\(\times\)}
  & $992$ & $952$ & \textcolor{black!80}{\textit{1.04}\(\times\)}
  & $992$ & $976$ & \textcolor{black!80}{\textit{1.02}\(\times\)} \\
RSS (GiB)
  & $18.02$ & $4.02$ & \textcolor{black!80}{\textit{4.48}\(\times\)}
  & $20.01$ & $8.01$ & \textcolor{black!80}{\textit{2.50}\(\times\)}
  & \textit{--} & \textit{--} & \textcolor{black!80}{\textit{--}}
  & $31.99$ & $31.99$ & \textcolor{black!80}{\textit{1.00}\(\times\)} \\
\midrule
\multicolumn{13}{l}{\textbf{Hadamard product}} \\
Prover (s)
  & $3.635$ & $0.843$ & \textcolor{black!80}{\textit{4.31}\(\times\)}
  & $3.108$ & $1.100$ & \textcolor{black!80}{\textit{2.82}\(\times\)}
  & $5.100$ & $2.798$ & \textcolor{black!80}{\textit{1.82}\(\times\)}
  & $3.289$ & $3.385$ & \textcolor{black!80}{\textit{0.97}\(\times\)} \\
Verifier ($\mu$s)
  & $650.75$ & $413.53$ & \textcolor{black!80}{\textit{1.57}\(\times\)}
  & $268.62$ & $111.51$ & \textcolor{black!80}{\textit{2.41}\(\times\)}
  & $86.18$ & $53.02$ & \textcolor{black!80}{\textit{1.63}\(\times\)}
  & $81.08$ & $34.89$ & \textcolor{black!80}{\textit{2.32}\(\times\)} \\
Payload (B)
  & $1{,}472$ & $1{,}872$ & \textcolor{black!80}{\textit{0.79}\(\times\)}
  & $1{,}472$ & $1{,}392$ & \textcolor{black!80}{\textit{1.06}\(\times\)}
  & $1{,}472$ & $1{,}168$ & \textcolor{black!80}{\textit{1.26}\(\times\)}
  & $1{,}472$ & $1{,}072$ & \textcolor{black!80}{\textit{1.37}\(\times\)} \\
RSS (GiB)
  & $18.00$ & $4.02$ & \textcolor{black!80}{\textit{4.48}\(\times\)}
  & $19.99$ & $8.03$ & \textcolor{black!80}{\textit{2.49}\(\times\)}
  & $23.99$ & $16.00$ & \textcolor{black!80}{\textit{1.50}\(\times\)}
  & $31.99$ & $32.00$ & \textcolor{black!80}{\textit{1.00}\(\times\)} \\
\bottomrule
\end{tabular*}
\small
\begin{tablenotes}[para,flushleft]
\item SFF denotes small-field-first; each ratio is SFF/Packed, computed
before rounding, so values above \(1\) favor Packed. For repeated runs,
timing ratios divide the per-method medians.
Columns denote base-field storage widths: the primes are \(251\),
\(65521\), \(2^{31}-1\) (Mersenne31 in a 32-bit word), and
\(2^{64}-2^{32}+1\), with packing factors \(\kappa=16,8,4,2\),
respectively.
All runs use an Intel Xeon Gold 6342, AVX-512, and 64 Rayon workers.
The 8/16-bit timings are single runs without warmup, with SFF first.
The 32-bit IP results use the joint-packing variant and medians of six
runs per method without warmup in one process, balanced over execution
order, with the controller pinned to CPU80 and NUMA interleaving.
The 32-bit Hadamard results use the public \texttt{WIDE\_V1} path:
one timing run per method without warmup in separate processes,
with NUMA interleaving and a shared worker CPU set, without the IP
experiment's per-worker pinning or CPU80 controller binding.
The 64-bit timings are historical medians of eight alternating-order
runs after two warmups.
Prover time includes packing and witness-dependent preprocessing.
Verifier time excludes deferred source-oracle evaluation.
Payload counts fixed-width field-coordinate messages; commitments
and PCS are excluded.
RSS is the observed prover-window process memory, including inputs,
measured separately from timing in isolated runs; it is not the
fixture-inclusive process peak. The 32-bit Hadamard RSS is remeasured
for \texttt{WIDE\_V1}; \textit{--} denotes RSS not remeasured for the
updated 32-bit IP comparison.
\end{tablenotes}
\end{threeparttable}
\end{table*}

\input{tab/tab-e2e}
\input{tab/tab-binary}

\section{Conclusion}
\label{sec:conclusion}

\methodname uses vector polynomial packing to reduce the dimension of
small-field relations before proving them over a large challenge
field. For inner and Hadamard products, the prover computes
row-polynomial products over \(\Field\) and reduces the original
length-\(N\) claim to a length-\(N/\kappa\) inner product over
\(\ExtensionField\).
Under the assumptions of Section~\ref{sec:highLevelEfficiency},
this gives \(O(\kappa N)\) base-field operations and \(O(N)\)
base-field coordinates of prover memory for the relation
reductions, excluding commitments and openings.
For \(|\Field|<\kappa\) and power-of-two \(\kappa\),
intermediate-field tabulation attains the same online Hadamard
bound with additional reusable tables. Section~\ref{sec:hadamard}
gives their preprocessing and storage costs.

The packed opening protocol authenticates challenge-dependent
evaluations against a fixed-basis witness commitment.
Together with the position-aware permutation argument, it yields
protocols for a single circuit execution and for multiple
executions of the same circuit.
The construction extends to binary arithmetizations by using
\(\mathbb F_2\subset\mathbb F_{2^8}\subset\ExtensionField\),
with byte-field arithmetic in the working-field PIOP.
Our component benchmarks show that the arithmetic and memory
savings are largest for the smallest witness fields tested.

Further work includes evaluating concrete PCS instantiations in
end-to-end proofs of quantized inference and adding compatible
zero-knowledge masking to the subprotocol transcripts.
The underlying reduction is modular: vector polynomial packing exploits
the witness field, and the terminal proof and commitment scheme
authenticate the resulting shorter relation.


\clearpage

\section*{Acknowledgement}
The authors used a generative AI tool to assist with language editing
and the organization of the exposition.

\bibliographystyle{splncs04}
\bibliography{ref}

\appendix
\section{Security Definitions}
\label{appx:def}

\subsection{Completeness and Soundness}
\label{appx:soundness}
\begin{definition}[Proof/Argument/Argument of Knowledge]
\label{def:proofArgumentAoK}
A protocol consists of algorithms $(\Setup,\Prover,\Verifier)$.
We distinguish proofs, arguments, and arguments of knowledge by
requiring soundness, computational soundness, and knowledge
soundness, respectively. Throughout, we require perfect completeness.
\begin{itemize}
	\item $\Setup$:
	On input $1^\lambda$, output public parameters $\pp$.
	\item Prover $\Prover$:
	On input $\pp$ and a statement $u$, produce a transcript $\pi$
    for membership in $\Language$. For a relation $\Relation$,
    the prover additionally receives a witness $w$.
	\item Verifier $\Verifier$:
	On input $\pp$, a statement $u$, and a transcript $\pi$,
    output {\sf true} or {\sf false}.
\end{itemize}
Let $\PVinteract$ denote the interaction between $\Prover$ and $\Verifier$.
\end{definition}

\begin{definition}[Perfect Completeness]
	\label{def:perfect-completeness}
The protocol $(\Setup,\Prover,\Verifier)$ for a relation \(\Relation\)
has perfect completeness if, for every \((u,w)\in\Relation\):
\[
    \Pr\left[
        \Verifier(\pp,u,\pi)=\TRUE
        \;\middle|\;
        \pp\leftarrow\Setup(1^\lambda),\;
        \pi\leftarrow\Prover(\pp,u,w)
    \right]
    =1,
\]
where the probability is over the randomness of
\(\Setup\), \(\Prover\), and \(\Verifier\).
\end{definition} 

\begin{definition}[Soundness~\cite{iacr/CanettiCHLRR18}]
	\label{def:soundness}
	A protocol $(\Setup,\Prover,\Verifier)$ for a language
    $\Language$ or relation $\mathcal R$ has soundness error
    $\kappa^\mathsf{s}(u)\colon\mathbb N\mapsto[0,1]$ if every
    computationally unbounded interactive prover $\Prover^*$ satisfies:
	\begin{equation*}
	% \label{eq:soundnessExpectedTime}
	\Pr
		\left(
			\left.
			\begin{aligned}
				& \Verifier(\pp, u, \pi) = \TRUE \\
			\end{aligned}
			\right.
			\left|
			\begin{aligned}
				& u \notin \Language~/~ \Relation;~ \pi \leftarrow (\Prover^* \interact \Verifier) (\pp, u) \\
			\end{aligned}
			\right.
		\right)
	\le \kappa^\mathsf{s}(u).
	\end{equation*}
    The protocol is sound if $\kappa^\mathsf{s}(u)$ is negligible.
\end{definition}

Restricting the adversarial prover to probabilistic polynomial time
gives computational soundness.
\begin{definition}[Computational Soundness]
	\label{def:ComputationalSoundness}
	A protocol $(\Setup,\Prover,\Verifier)$ for a language
    $\Language$ or relation $\mathcal R$ has computational soundness
    error $\kappa^\mathsf{cs}(u)\colon\mathbb N\mapsto[0,1]$ if
    every PPT prover $\Prover^*$ satisfies:
	\begin{equation*}
	% \label{eq:soundnessExpectedTime}
	\Pr
		\left(
			\left.
			\begin{aligned}
				& \Verifier(\pp, u, \pi) = \TRUE \\
			\end{aligned}
			\right.
			\left|
			\begin{aligned}
				& u \notin \Language~/~\Relation;~ \pi \leftarrow (\Prover^* \interact \Verifier)(\pp, u) \\
			\end{aligned}
			\right.
		\right)
	\le \kappa^\mathsf{cs}(u).
	\end{equation*}
    The protocol is computationally sound if
    $\kappa^\mathsf{cs}(u)$ is negligible.
\end{definition}
By definition, $\kappa^\mathsf{cs}(u)\leq\kappa^\mathsf{s}(u)$.

An argument of knowledge additionally requires extraction of a
witness~\cite{crypto/AttemaCK21}.
\begin{definition}[Knowledge Soundness]
	\label{def:knowledgeSoundness}
	A protocol $(\Setup,\Prover,\Verifier)$ for a relation $\Relation$
    has knowledge-soundness error
    $\kappa^\mathsf{ks}(u)\colon\mathbb N\mapsto[0,1]$ if there
    exists a knowledge extractor $\Extractor$ with the following
    guarantee. Given $u$ and black-box access to a prover $\Prover^*$,
    the extractor runs in expected time polynomial in $|u|$, counting
    each oracle query as one step, and satisfies:
	\begin{equation*}
	% \label{eq:soundnessExpectedTime}
	\Pr((u, \Extractor(\Prover^*, u) ) \in \Relation )
	\ge
	\epsilon(\Prover^*, u) - \kappa^\mathsf{ks}(u),
	\end{equation*}
	where $\epsilon(\Prover^*,u)$ is the acceptance probability of
    $\Prover^*$. The protocol is knowledge sound if
    $\kappa^\mathsf{ks}(u)$ is negligible.
\end{definition}

If a protocol is knowledge-sound with knowledge-soundness error
$\kappa^\mathsf{ks}(u)$, then it is also computationally sound with
error $\kappa^\mathsf{cs}(u) \le \kappa^\mathsf{ks}(u)$.
Otherwise, some false statement $u$ and PPT prover $\Prover^*$
would satisfy:
\[
	\Pr
	\left(
		\left.
		\begin{aligned}
			& \Verifier(\pp, u, \pi) = \TRUE \\
		\end{aligned}
		\right.
		\left|
		\begin{aligned}
			& u \notin \mathcal{R};~ \pi \leftarrow (\Prover^* \interact \Verifier)(\pp, u) \\
		\end{aligned}
		\right.
	\right)
	> \kappa^\mathsf{ks}(u).
\]
The extraction guarantee would then yield a witness $w$ with
$(u,w)\in\Relation$ with positive probability, contradicting the
falsity of $u$.


\subsection{Round-by-Round Soundness}
\label{appx:rbyDef}

\begin{definition}[Round-by-Round Soundness~\cite{iacr/CanettiCHLRR18}]
	\label{def:roundByRoundSoundness}
	Let $\Pi=(\Setup,\Prover,\Verifier)$ be a $(2\nu+1)$-message
    public-coin interactive proof for a relation $\Relation$.
    For a statement $u\in\{0,1\}^*$, write
    $\pi_j=(\alpha_0,x_1,\alpha_1,\ldots,\alpha_{j-1},x_j,\alpha_j)$
    for a transcript prefix containing the first $2j+1$ messages,
    where $j\in[0,\nu]$. Let $\mathsf H(u,\pi_j)$ denote the
    distribution of the verifier's next message.
    The protocol has round-by-round soundness error
    $\{\kappa_j^\mathsf{rbr}\}_{j\in[0,\nu]}$ if there exist
    deterministic, not necessarily efficient, functions
    $\{\mathsf{State}_j\}_{j\in[0,\nu]}$ mapping a statement and
    transcript prefix to $\mathsf{accept}$ or $\mathsf{reject}$,
    with the following properties:
		\begin{itemize}
			\item If $u \notin \Language$, then $\mathsf{State}_0(u, \emptyset) = \mathsf{reject}$, where $\emptyset$ denotes the empty transcript.
			\item If $\mathsf{State}_{j-1}(u,\pi_{j-1})=\mathsf{reject}$,
            then the next verifier challenge and prover response satisfy:
			\[
				\Pr\!\!
				\left(
					\mathsf{State}_j(u, \pi_j) = \mathsf{accept~}\!\!
					\left|~\!
						\begin{aligned}
							&\!\mathsf{State}_{j - 1}(u, \pi_{j - 1}) = \mathsf{reject}; \\
							& x_j \leftarrow \mathsf{H}(u, \pi_{j - 1}); \pi_j \leftarrow (\pi_{j - 1} \| x_j \| \alpha_j)\!
						\!\end{aligned}
					\right.
				\right)\!
				\!\le\!\kappa_j^\mathsf{rbr}.
			\]
			\item For any full transcript $\pi$ for which the verifier halts, if $\mathsf{State}_\nu(u, \pi) = \mathsf{reject}$ then $\Verifier(u, \pi) = \FALSE$.
		\end{itemize}
\end{definition}

The following lemma bounds the overall error by the sum of the
round-by-round errors.
\begin{lemma}[Round-by-Round Soundness Implies Soundness~\hspace{-1.58pt}\textup{\cite[Proposition~5.4]{iacr/CanettiCHLRR18}}]
\label{lemma:rbr}
	Let $\Pi$ be a $(2 \nu + 1)$-move interactive proof with round-by-round soundness error $\kappa_j^\mathsf{rbr}, j \in [\nu]$.
	Then $\Pi$ has soundness error at most
    $\sum_{j=1}^\nu\kappa_j^\mathsf{rbr}$.
\end{lemma}

\clearpage
\section{Deferred Technical Details}
\label{appx:tech}

\subsection{Pseudocode}
\label{appx:pseudocode}

We give the protocols in order of composition:
Algorithms~\ref{alg:ip}, \ref{alg:hadamard}, and
\ref{alg:permutation} implement the packed reductions of
Section~\ref{sec:blocks}; Algorithm~\ref{alg:packed-pcs}
authenticates their terminal evaluations as in Section~\ref{sec:pcs};
and Algorithms~\ref{alg:single-instance-packing} and
\ref{alg:multiple-instance-packing} compose these components
for the two settings of Section~\ref{sec:together}.
For the intermediate-field variant of Section~\ref{sec:hadamard},
Algorithm~\ref{alg:hadamard} uses the Lagrange basis over
\(\mathbb K\) and computes its quotient via
Eq.~\eqref{eq:hadamard-small-field-aggregation}.
Tabulation changes this computation, not the messages or checks.

% !TEX root = ../main.tex
\begin{algorithm}[!t]
\caption{Packed inner-product proof}
\label{alg:ip}
\KwIn{
\(N=m\kappa\), where \(m, \kappa \in 2^\mathbb{N}\);
\(c\in\Field\);
\(\bm a,\bm b\in\Field^N\)
}
\KwOut{
Terminal values \(\ext{v}_1,\ext{v}_2\) and MLE points
\(\ext{\bm y}_a, \ext{\bm y}_b \in\ExtensionField^{\log_2 N}\);
acceptance bit \(\mathsf{ok}\)
}

Arrange \(\bm a,\bm b\) columnwise as
\(\bm A\leftarrow\concat{\bm a^{(1)},\ldots,\bm a^{(\kappa)}}\)
and \(\bm B\leftarrow\concat{\bm b^{(1)},\ldots,\bm b^{(\kappa)}}
\in\Field^{m\times\kappa}\)\;

\tcp{Compute the auxiliary coefficients over the base field}
$\Prover$ sets:\\
\Indp
\(
    \bm g_A(X)
    \leftarrow
    \sum_{j=1}^{\kappa}
    \bm a^{(j)}X^{j-1}
\);~
\(
    \bm g_B(X)
    \leftarrow
    \sum_{j=1}^{\kappa}
    \bm b^{(j)}X^{j-1}
\)\;

\ForEach{
\(\tau\in
\{-(\kappa-1),\ldots,-1,1,\ldots,\kappa-1\}\)
}{
    Set \(
        f_\tau
        \leftarrow
        \sum_{\substack{j,j'\in[\kappa]\\j-j'=\tau}}
        \left\langle
            \bm a^{(j)},\bm b^{(j')}
        \right\rangle
    \)\;
}

Set
\(
    h_+(X)
    \leftarrow
    \sum_{\tau=1}^{\kappa-1}
    f_\tau X^{\tau-1}
\);~
\(
    h_-(X)
    \leftarrow
    \sum_{\tau=1}^{\kappa-1}
    f_{-\tau}X^{\tau-1}
\)\;
\Indm

\(\PtoV\):
the coefficients of \(h_+(X)\) and \(h_-(X)\)\;

\(\VtoP\):
\(\ext{x}\ChooseRandom\ExtensionField^\times\)\;

\tcp{Reduce to a length-\(m\) extension-field inner product}
$\Prover$ sets
\(
    \ext{\bm a}
    \leftarrow
    \bm g_A(\ext{x})
    \in\ExtensionField^m
\),
\(
    \ext{\bm b}
    \leftarrow
    \bm g_B(\ext{x}^{-1})
    \in\ExtensionField^m
\)\;

$\Verifier$ sets
\(
    \ext{c}
    \leftarrow
    c
    +\ext{x}h_+(\ext{x})
    +\ext{x}^{-1}h_-(\ext{x}^{-1})
\)\;

$\Prover$ and $\Verifier$ apply Fact~\ref{fact:ip} to
\(
    \ext{c}
    \EqOrNot
    \left\langle
        \ext{\bm a},\ext{\bm b}
    \right\rangle,
\)
and obtain terminal values
\(
    \ext{v}_1,\ext{v}_2
\),
the challenge point
\(
    \ext{\bm y}'\in
    \ExtensionField^{\log_2m}
\),
and the acceptance bit \(\mathsf{ok}\)\;

\(\Prover\) and \(\Verifier\) set:\\
\Indp
\(
    \ext{\bm y}_a
    \leftarrow
    \bigl(
        \ext{x}^{2^{\log_2\kappa-1}},
        \ext{x}^{2^{\log_2\kappa-2}},
        \ldots,
        \ext{x}
    \bigr)
    \mathbin{\|}
    \ext{\bm y}' \in \ExtensionField^{\log_2 N}
\)\;
\(
    \ext{\bm y}_b
    \leftarrow
    \bigl(
        \ext{x}^{-2^{\log_2\kappa-1}},
        \ext{x}^{-2^{\log_2\kappa-2}},
        \ldots,
        \ext{x}^{-1}
    \bigr)
    \mathbin{\|}
    \ext{\bm y}' \in \ExtensionField^{\log_2 N}
\)\;
\Indm

Output the terminal claims
\(
    \ext{v}_1
    \EqOrNot
    \left\langle
        \bm a,
        \tensorspan{\ext{\bm y}_a}
    \right\rangle
\) and
\(
    \ext{v}_2
    \EqOrNot
    \left\langle
        \bm b,
        \tensorspan{\ext{\bm y}_b}
    \right\rangle
\)\;

\Return
\(
    (\mathsf{ok},\ext{v}_1,\ext{v}_2,\ext{\bm y}_a, \ext{\bm y}_b)
\)\;
\end{algorithm}

% !TEX root = ../main.tex
\begin{algorithm}[!t]
\caption{Packed Hadamard-product proof}
\label{alg:hadamard}
\KwIn{
\(N=m\kappa\), \(m\in2^{\mathbb N}\), \(\kappa\geq2\);
a fixed claim \(\ext c\in\ExtensionField\) and matrices
\(\bm A,\bm B\in\Field^{m\times\kappa}\) fixed by the input commitments;
column weights \(\ext{\bm r}\in\ExtensionField^\kappa\);
row weights
\(\ext{\bm l}=\bigotimes_{j=1}^{\log_2m}
(\ext l_{j,0},\ext l_{j,1})^\top\);
nodes \(\mathcal H=\{\gamma_1,\ldots,\gamma_\kappa\}\subseteq\mathbb K\),
with \(\Field\subseteq\mathbb K\subseteq\ExtensionField\) and
\(|\mathbb K|\geq2\kappa-1\)
}
\KwOut{
Terminal values \(\ext v_1,\ext v_2\);
column weights \(\ext{\bm\rho}\in\ExtensionField^\kappa\);
tensor row weights \(\ext{\bm\xi}_1,\ext{\bm\xi}_2\in\ExtensionField^m\);
acceptance bit \(\mathsf{ok}\)
}

\tcp{Encode the weighted column claims}
\(\Prover\) computes
\(\ext{\bm c}^{\top}\leftarrow
\ext{\bm l}^{\top}(\bm A\circ\bm B)\)\;
In the Lagrange basis on \(\mathcal H\), define
\(\bm g_A(X)\leftarrow\sum_j\bm a^{(j)}\psi_j(X)\),
\(\bm g_B(X)\leftarrow\sum_j\bm b^{(j)}\psi_j(X)\), and
\(g_C(X)\leftarrow\sum_j\ext c_j\psi_j(X)\)\;
\(\Prover\) computes the residual
\(f(X)\leftarrow
\langle\ext{\bm l},\bm g_A(X)\circ\bm g_B(X)\rangle-g_C(X)\)
at \(\kappa-1\) fixed nodes in \(\mathbb K\setminus\mathcal H\),
interpolates with the zeros on \(\mathcal H\), and sets
\(h(X)\leftarrow f(X)/\Psi(X)\)\;
\(\PtoV\): \(\ext{\bm c}\) and the \(\kappa-1\) coefficients of \(h\)\;
\(\Verifier\) sets
\(\mathsf{ok}_0\leftarrow
(\ext c\EqOrNot\langle\ext{\bm c},\ext{\bm r}\rangle)\)\;
\(\VtoP\): \(\ext x\ChooseRandom\ExtensionField\setminus\mathcal H\)\;

\tcp{Reduce to a length-\(m\) extension-field inner product}
\(\Prover\) and \(\Verifier\) compute
\(\ext{\bm\rho}\leftarrow
(\psi_1(\ext x),\ldots,\psi_\kappa(\ext x))^\top\)\;
\(\Prover\) computes
\(\ext{\bm a}\leftarrow\bm A\ext{\bm\rho}\) and
\(\ext{\bm b}\leftarrow\bm B\ext{\bm\rho}\)\;
\(\Verifier\) computes
\(\ext c_x\leftarrow
\langle\ext{\bm c},\ext{\bm\rho}\rangle=g_C(\ext x)\)\;
\(\Prover\) and \(\Verifier\) apply Fact~\ref{fact:ip} to
\(\langle\ext{\bm a},\ext{\bm b}\circ\ext{\bm l}\rangle
\EqOrNot\ext c_x+\Psi(\ext x)h(\ext x)\),
obtaining \(\mathsf{ok}_{\mathsf{IP}}\),
\(\ext v_1,\ext v_2\), and
\(\ext{\bm y}\in\ExtensionField^{\log_2m}\)\;

\tcp{Output the column weights and tensor row factors}
Set
\(\ext{\bm\xi}_1\leftarrow\tensorspan{\ext{\bm y}}\) and
\(\ext{\bm\xi}_2\leftarrow
\ext{\bm l}\circ\tensorspan{\ext{\bm y}}\)\;
Represent the \(j\)-th factor of \(\ext{\bm\xi}_2\) by
\((\ext l_{j,0}(1-\ext y_j),\ext l_{j,1}\ext y_j)^\top\)\;
Output
\(\ext v_1\EqOrNot\ext{\bm\xi}_1^\top\bm A\ext{\bm\rho}\)
and
\(\ext v_2\EqOrNot\ext{\bm\xi}_2^\top\bm B\ext{\bm\rho}\)\;
Set
\(\mathsf{ok}\leftarrow\mathsf{ok}_0\wedge\mathsf{ok}_{\mathsf{IP}}\)\;
\Return{\((\mathsf{ok},\ext v_1,\ext v_2,\ext{\bm\rho},
\ext{\bm\xi}_1,\ext{\bm\xi}_2)\)}\;
\end{algorithm}

% !TEX root = ../main.tex
\begin{algorithm}[!t]
\linespread{0.97}\selectfont
\caption{Packed permutation proof}
\label{alg:permutation}
\KwIn{
\(N=m\kappa\), where \(m,\kappa\in2^{\mathbb N}\);
\(\bm a,\bm b\in\Field^N\)
}
\KwOut{
Terminal values \(\ext{v}_a,\ext{v}_b\) and MLE points
\(\ext{\bm y}_a,\ext{\bm y}_b
\in\ExtensionField^{\log_2N}\);
acceptance bit \(\mathsf{ok}\)
}

\tcp{Reduce multiset equality to equality of grand products}
Fix the commitments to \(\bm a\) and \(\bm b\), then sample
\(\VtoP\):
\(\ext z\ChooseRandom\ExtensionField\)\;

\ForEach{\(\bm d\in\{\bm a,\bm b\}\)}{
    Arrange \(\bm d\) columnwise as
    \(
        \bm D^{(\bm d)}
        \leftarrow
        \concat{
            \bm d^{(1)},\ldots,\bm d^{(\kappa)}
        }
        \in\Field^{m\times\kappa}.
    \)

    \(\Prover\) initializes the column-product tree with
    \(
        \bm C^{(\bm d)}_0(Z)
        \leftarrow
        Z\bm 1_{m\times\kappa}
        +
        \bm D^{(\bm d)}
    \)\;

    \For{\(t=0,\ldots,\log_2\kappa-1\)}{
        \(
            \bm C^{(\bm d)}_{t+1}(Z)
            \leftarrow
            \bm C^{(\bm d)}_{t,L}(Z)
            \circ
            \bm C^{(\bm d)}_{t,R}(Z).
        \)
    }

    \(\Prover\) evaluates the last column layer:
    \(
        \ext{\bm c}^{(\bm d)}_0
        \leftarrow
        \bm C^{(\bm d)}_{\log_2\kappa}(\ext z)
        \in\ExtensionField^m.
    \)

    \For{\(i=0,\ldots,\log_2m-1\)}{
        \(
            \ext{\bm c}^{(\bm d)}_{i+1}
            \leftarrow
            \ext{\bm c}^{(\bm d)}_{i,L}
            \circ
            \ext{\bm c}^{(\bm d)}_{i,R}.
        \)
    }

    Set
    \(
        \Theta_{\bm d}
        \leftarrow
        \ext c^{(\bm d)}_{\log_2m}.
    \)
}

\(\PtoV\):
\(\Theta_{\bm a},\Theta_{\bm b}\)\;

\(\Verifier\) checks
\(
    \mathsf{ok}_{\mathsf{perm}}
    \leftarrow
    \left(
        \Theta_{\bm a}
        \EqOrNot
        \Theta_{\bm b}
    \right).
\)

\ForEach{\(\bm d\in\{\bm a,\bm b\}\)}{
    Set
    \(
        \ext{s}^{(\bm d)}_{\log_2m}
        \leftarrow
        \Theta_{\bm d}
    \)
    and
    \(
        \ext{\bm\xi}^{(\bm d)}_{\log_2m}
        \leftarrow(1).
    \)

    \tcp{Unfold the row-product tree}
    \For{\(i=\log_2m-1,\ldots,0\)}{
        Write the current claim as
        \(
            \ext{s}^{(\bm d)}_{i+1}
            \EqOrNot
            \left\langle
                \ext{\bm c}^{(\bm d)}_{i,L},
                \ext{\bm c}^{(\bm d)}_{i,R}
                \circ
                \ext{\bm\xi}^{(\bm d)}_{i+1}
            \right\rangle.
        \)

        Apply Fact~\ref{fact:ip} and randomly combine the two
        terminal claims to obtain
        \(
            \ext{s}^{(\bm d)}_i
            \EqOrNot
            \left\langle
                \ext{\bm\xi}^{(\bm d)}_i,
                \ext{\bm c}^{(\bm d)}_i
            \right\rangle
        \)
        and the acceptance bit
        \(\mathsf{ok}^{(\bm d)}_{\mathsf{row},i}\), retaining
        the tensor factors of \(\ext{\bm\xi}^{(\bm d)}_i\)\;
    }

    Set
    \(
        \ext{q}^{(\bm d)}_{\log_2\kappa}
        \leftarrow
        \ext{s}^{(\bm d)}_0
    \)
    and
    \(
        \ext{\bm\zeta}^{(\bm d)}_{\log_2\kappa}
        \leftarrow
        \ext{\bm\xi}^{(\bm d)}_0.
    \)

    \tcp{Unfold the column-product tree}
    \For{\(t=\log_2\kappa-1,\ldots,0\)}{
        Write the current claim as
        \(
            \ext{q}^{(\bm d)}_{t+1}
            \EqOrNot
            \left\langle
                \ext{\bm\zeta}^{(\bm d)}_{t+1},
                \operatorname{vec}
                \bigl(
                    \bm C^{(\bm d)}_{t,L}(\ext z)
                \bigr)
                \circ
                \operatorname{vec}
                \bigl(
                    \bm C^{(\bm d)}_{t,R}(\ext z)
                \bigr)
            \right\rangle.
        \)

        Apply the coefficient-packing reduction of
        Algorithm~\ref{alg:hadamard} and use Fact~\ref{fact:ip}
        for its evaluated inner product to obtain
        \(
            \ext{q}^{(\bm d)}_t
            \EqOrNot
            \left\langle
                \ext{\bm\zeta}^{(\bm d)}_t,
                \operatorname{vec}
                \bigl(
                    \bm C^{(\bm d)}_t(\ext z)
                \bigr)
            \right\rangle
        \)
        and the acceptance bit
        \(\mathsf{ok}^{(\bm d)}_{\mathsf{col},t}\), retaining
        the tensor factors of \(\ext{\bm\zeta}^{(\bm d)}_t\)\;
    }

    Write
    \(
        \ext{\bm\zeta}^{(\bm d)}_0
        =
        \tensorspan{\ext{\bm y}_{\bm d}}
    \)
    for
    \(
        \ext{\bm y}_{\bm d}
        \in\ExtensionField^{\log_2N}.
    \)

    \(\Prover\) and \(\Verifier\) remove the public shift:
    \(
        \ext{v}_{\bm d}
        \leftarrow
        \ext{q}^{(\bm d)}_0
        -
        \ext z
        \left\langle
            \ext{\bm\zeta}^{(\bm d)}_0,
            \bm 1_N
        \right\rangle.
    \)

    Output the terminal claim
    \(
        \ext{v}_{\bm d}
        \EqOrNot
        \left\langle
            \bm d,
            \tensorspan{\ext{\bm y}_{\bm d}}
        \right\rangle.
    \)
}

Set
\(
    \ext{v}_a\leftarrow\ext{v}_{\bm a}
\),
\(
    \ext{v}_b\leftarrow\ext{v}_{\bm b}
\),
\(
    \ext{\bm y}_a\leftarrow\ext{\bm y}_{\bm a}
\),
and
\(
    \ext{\bm y}_b\leftarrow\ext{\bm y}_{\bm b}.
\)

Set
\(
    \mathsf{ok}
    \leftarrow
    \mathsf{ok}_{\mathsf{perm}}
    \wedge
    \bigwedge_{\bm d\in\{\bm a,\bm b\}}
    \left(
        \bigwedge_{i=0}^{\log_2m-1}
        \mathsf{ok}^{(\bm d)}_{\mathsf{row},i}
    \right)
    \wedge
    \bigwedge_{\bm d\in\{\bm a,\bm b\}}
    \left(
        \bigwedge_{t=0}^{\log_2\kappa-1}
        \mathsf{ok}^{(\bm d)}_{\mathsf{col},t}
    \right).
\)

\Return
\(
    (
        \mathsf{ok},
        \ext{v}_a,
        \ext{v}_b,
        \ext{\bm y}_a,
        \ext{\bm y}_b
    )
\)\;
\end{algorithm}

% !TEX root = ../main.tex
\begin{algorithm}[!t]
\caption{Packed opening protocol}
\label{alg:packed-pcs}
\KwPubIn{
\(\mathsf{pp}_{\pc},\mathsf{cm}\);
\(\Field\subset\ExtensionField\) with basis \(\ext{\bm\beta}\)
and \(\ext\beta_1=1\);
\(N=m\kappa\), \(\ell=\log_2m\);
\(\ext v\in\ExtensionField\),
\(\ext{\bm\rho}\in\ExtensionField^\kappa\), and
\(\ext{\bm\xi}=\bigotimes_{j=1}^{\ell}
\ext{\bm\xi}^{(j)}\), given by factors
\(\ext{\bm\xi}^{(j)}=
(\ext\xi_0^{(j)},\ext\xi_1^{(j)})^\top\in\ExtensionField^2\)
}
\KwPriIn{
\(\Prover\): \(\bm A\in\Field^{m\times\kappa}\) and
commitment randomness \(\rho_{\pc}\)
}
\KwOut{Acceptance bit \(\mathsf{ok}\)}

\(\Verifier\) returns \(0\) on a malformed message or a failed
check, reduction, or opening\;
\tcp{Authenticate the column evaluations by an extension-field projection}
\(\Prover\) sets
\(\ext{\bm a}\leftarrow\bm A\ext{\bm\beta}\) and decodes
\(\bm\Xi\in\Field^{m\times\kappa}\) from
\(\ext{\bm\xi}=\bm\Xi\ext{\bm\beta}\)\;
\(\Prover\) computes \(\bm U\leftarrow\bm\Xi^\top\bm A\)\;
\(\PtoV\): \(\ext{\bm u}\leftarrow\bm U^\top\ext{\bm\beta}
\in\ExtensionField^\kappa\)\;
\(\Verifier\) checks
\(\ext v\EqOrNot\ext{\bm u}^{\top}\ext{\bm\rho}\)\;
\(\Verifier\) recovers \(\bm U\in\Field^{\kappa\times\kappa}\)
from \(\ext{\bm u}^{\top}=\ext{\bm\beta}^{\top}\bm U\)\;
\(\VtoP\): \(\ext{\bm r}\ChooseRandom\ExtensionField^\kappa\)\;
\(\Verifier\) computes
\(\ext c\leftarrow\ext{\bm r}^{\top}\bm U\ext{\bm\beta}\)\;
\(\Prover\) computes the operand
\(\bm\Xi\ext{\bm r}\in\ExtensionField^m\)\;
\(\Prover,\Verifier\) apply Fact~\ref{fact:ip} over
\(\ExtensionField\) to
\(\ext c\EqOrNot\langle\ext{\bm a},\bm\Xi\ext{\bm r}\rangle\),
and obtain terminal values \(\ext w_{\mathsf{com}},\ext w\)
and point \(\ext{\bm y}\in\ExtensionField^\ell\)\;
\tcp{The verifier contracts the public queries directly}
Use the multiplication matrices \(\bm M_{\ext z}\) from
Eq.~\eqref{eq:pcs-multiplication-matrix}, applied by extension-field
scaling of rows and columns without materializing \(\bm M_{\ext z}\)\;
\(\Verifier\) initializes
\(\bm W\leftarrow\bm e_1\bm e_1^\top\in\Field^{\kappa\times\kappa}\),
where \(\bm e_1=(1,0,\ldots,0)^\top\)\;
\For{\(j=\ell,\ldots,1\)}{
    \(\Verifier\) updates
    \(\bm W\leftarrow
       \bm W\bm M_{\ext\xi_0^{(j)}}^\top
       +\bm M_{\ext y_j}\bm W
        \bm M_{\ext\xi_1^{(j)}-\ext\xi_0^{(j)}}^\top\)\;
    \tcp{All occurrences on the right use the previous \(\bm W\), as in Eq.~\eqref{eq:pcs-query-recurrence}}
}
\(\Verifier\) checks
\(\ext w\EqOrNot\ext{\bm\beta}^{\top}\bm W\ext{\bm r}\)\;

\(\mathsf{ok}\leftarrow
\pc.\Open\!\left(
    \mathsf{pp}_{\pc},\mathsf{cm},\ext w_{\mathsf{com}},
    \ext{\bm y}:\ext{\bm a},\rho_{\pc}
\right)\)\;
\Return \(\mathsf{ok}\)\;
\end{algorithm}

% !TEX root = ../main.tex
\begin{algorithm}[!t]
\footnotesize
\caption{PLONK proof with single-instance packing}
\label{alg:single-instance-packing}
\KwPubIn{
Parameters of Theorem~\ref{theorem:single-instance-packing};
\(\bm{\mathsf{io}}\); the preprocessed public commitment of
logical length \(m_{\mathsf{pub}}\)
}
\KwPriIn{
\(\Prover\): \(\bm w_L,\bm w_R,\bm w_O\in\Field^N\);
circuit tables and their opening data
}
\KwOut{Acceptance bit \(\mathsf{ok}\)}

\(\Verifier\) returns \(0\) on a malformed message or a failed
check, reduction, or opening\;
\tcp{Commit before sampling challenges}
Arrange \(\bm w_\ell\) as \(\bm W_\ell\in\Field^{m\times\kappa}\)
for \(\ell\in\{L,R,O\}\)\;
\(\PtoV\): one commitment to the concatenation of
\(\bm W_L\ext{\bm\beta},\bm W_R\ext{\bm\beta},
\bm W_O\ext{\bm\beta}\), of logical length \(3m\)\;
Set the virtual vector \(\bm w_M\leftarrow\bm w_L\circ\bm w_R\)\;

\tcp{Reduce the gate claims; defer all terminal openings}
\(\VtoP\): \(\ext{\bm x}\ChooseRandom\ExtensionField^{\log_2N}\);
set \(\ext{\bm\xi}\leftarrow\tensorspan{\ext{\bm x}}\)\;
\(\PtoV\): \(\ext c_L,\ext c_R,\ext c_M,\ext c_O,\ext c_C\)
as in Eq.~\eqref{eq:plonkWeightedHadamard1} and
Eq.~\eqref{eq:plonkWeightedHadamard2}\;
\(\Verifier\) checks
\(\ext c_L+\ext c_R+\ext c_M+\ext c_O+\ext c_C\EqOrNot0\)\;
\ForEach{\(\ell\in\{L,R,O\}\)}{
    Apply Algorithm~\ref{alg:hadamard} to
    \(\ext c_\ell\EqOrNot
    \langle\ext{\bm\xi},\bm q_\ell\circ\bm w_\ell\rangle\);
    retain both terminal queries\;
}
Apply Algorithm~\ref{alg:hadamard} to
\(\ext c_M\EqOrNot
\langle\ext{\bm\xi},\bm q_M\circ\bm w_M\rangle\);
retain its selector query\;
Express its product-vector query as
\(\ext v\EqOrNot
\ext{\bm l}^{\top}(\bm W_L\circ\bm W_R)\ext{\bm r}\).
Apply Algorithm~\ref{alg:hadamard} again and retain both terminals\;
Retain \(\ext c_C\EqOrNot\langle\ext{\bm\xi},\bm q_C\rangle\)\;

\tcp{Reduce the copy claims over \(\mathbb K\)}
Form \(\bm w\in\Field^L\), including the public I/O, and
\(\bm a\leftarrow\bm\varsigma+\theta_d\bm w\),
\(\bm b\leftarrow\bm\iota+\theta_d\bm w\)\;
Pad both operands with zeros to length \(4N\).
Apply Algorithm~\ref{alg:permutation} with base field \(\mathbb K\),
extension \(\ExtensionField\), and packing width \(\kappa_{\mathbb K}\),
sampling its fresh challenge
\(\ext z\ChooseRandom\ExtensionField\setminus\{0\}\)\;
Expand its two weighted terminals as in
Appendix~\ref{appx:proof-single-instance}.
Restrict to the wire segments and regroup the tensor factors
to packing width \(\kappa\)\;
\(\PtoV\): six wire values and
\(3\lceil\log_{\underline q}L\rceil\) permuted-digit values;
retain their weighted queries\;
\(\Verifier\) evaluates the identity-label, public-I/O, and padding
contributions directly and checks both recombinations\;

\tcp{Batch the witness and public openings}
For every retained query, execute the column-message and scalar-check
steps of Section~\ref{sec:pcs}.
Fix all column messages before sampling independent projection vectors\;
Project each column identity, extend its query by zeros outside
the committed block, and sum the inner-product claims in each group\;
For each group, apply Fact~\ref{fact:ip} and check its public
query-side evaluation as in Appendix~\ref{appx:proof-single-instance}\;
Authenticate the two terminal values by underlying PCS openings
at logical lengths \(3m\) and \(m_{\mathsf{pub}}\)\;
\Return \(1\)\;
\end{algorithm}

% !TEX root = ../main.tex
\begin{algorithm}[!t]
\footnotesize
\linespread{0.97}\selectfont
\caption{PLONK proof with packing across instances}
\label{alg:multiple-instance-packing}
\KwPubIn{
Parameters of Theorem~\ref{theorem:multiple-instance-packing};
the \(\kappa\) public-I/O vectors;
the preprocessed commitment to \(\bm Q\ext{\bm\beta}\in\ExtensionField^N\)
}
\KwPriIn{
\(\Prover\): \(\bm W_L,\bm W_R,\bm W_O\in\Field^{N\times\kappa}\);
circuit tables and their opening data
}
\KwOut{Acceptance bit \(\mathsf{ok}\)}

\(\Verifier\) returns \(0\) on a malformed message or a failed
check, reduction, or opening\;
\tcp{Commit before sampling challenges}
Set \(\ext{\bm w}_\ell\leftarrow\bm W_\ell\ext{\bm\beta}\)
for \(\ell\in\{L,R,O\}\), and the virtual vector
\(\ext{\bm w}_M\leftarrow(\bm W_L\circ\bm W_R)\ext{\bm\beta}\)\;
\(\PtoV\): one commitment to the concatenation of
\(\ext{\bm w}_L,\ext{\bm w}_R,\ext{\bm w}_O\),
of logical length \(3N\)\;

\tcp{Reduce the gate claims; defer all terminal openings}
\(\VtoP\): \(\ext{\bm x}\ChooseRandom\ExtensionField^{\log_2N}\)\;
Run sumcheck with initial claim \(0\) and individual degree at most
three for the \(\tensorspan{\ext{\bm x}}\)-weighted residual in
Eq.~\eqref{eq:together-multiple-gates};
obtain \(\ext{\bm y},\ext c\)\;
\(\PtoV\): selector values \(\ext v_\ell\) for
\(\ell\in\{L,R,M,O,C\}\) and wire values \(\ext u_\ell\) for
\(\ell\in\{L,R,M,O\}\), at \(\ext{\bm y}\)\;
\(\Verifier\) checks
\(\ext c\EqOrNot
\langle\tensorspan{\ext{\bm x}},\tensorspan{\ext{\bm y}}\rangle
[\sum_{\ell\in\{L,R,M,O\}}\ext v_\ell\ext u_\ell+
(\sum_{j=1}^{\kappa}\ext\beta_j)\ext v_C]\)\;
Retain the five selector queries and three ordinary wire queries\;
Apply Algorithm~\ref{alg:hadamard} to
\(\ext u_M\EqOrNot
\tensorspan{\ext{\bm y}}^\top
(\bm W_L\circ\bm W_R)\ext{\bm\beta}\);
retain its two weighted terminal queries\;

\tcp{Reduce the copy claims by two product trees}
Form \(\ext{\bm w}\in\ExtensionField^L\), including the packed public I/O\;
\(\VtoP\): fresh independent
\(\ext y,\ext z\ChooseRandom\ExtensionField\)\;
Construct product trees on
\(\ext z\bm1_L+\bm\iota+\ext y\ext{\bm w}\) and
\(\ext z\bm1_L+\bm\varsigma+\ext y\ext{\bm w}\),
padding each input with ones to length \(4N\)\;
\(\PtoV\): a common root \(\ext\Theta\)\;
For each tree, run the degree-three layer sumchecks of
Appendix~\ref{appx:proof-multiple-instance}, starting at \(\ext\Theta\).
After checking the sumcheck terminal identity and fixing both child values,
sample a fresh merge
challenge and continue to the next layer; retain the terminal input claim\;
Expand both terminals, retaining the prefix weights of the three
wire segments. \(\PtoV\): six wire values and three permuted-label values\;
\(\Verifier\) evaluates the identity-label, public-I/O, and padding
contributions directly and checks both recombinations\;

\tcp{Batch the witness and public openings}
Express the five selector and three permuted-label queries as weighted
queries to \(\bm Q\), as in Appendix~\ref{appx:proof-multiple-instance}\;
Execute the column-message and scalar-check steps of Section~\ref{sec:pcs}
for these eight queries and the two Hadamard terminals\;
Fix all column messages and the nine ordinary wire values.
Sample independent projection vectors for the packed queries and
independent scalar weights for the ordinary queries\;
Project the column identities and combine all inner-product claims
within the witness and public groups, extending queries by zeros
outside their committed blocks\;
For each group, apply Fact~\ref{fact:ip} and verify the combined query-side
evaluation, retaining all witness block-prefix weights\;
Authenticate the two terminal values by underlying PCS openings
at logical lengths \(3N\) and \(N\)\;
\Return \(1\)\;
\end{algorithm}


\clearpage

\subsection{Packed Polynomial Commitment Scheme}
\label{appx:pcs}

% !TEX root = main.tex
We express the ring-switching construction of Diamond and
Posen~\cite{eurocrypt/DiamondP26} in our vector and matrix notation
to support the weighted tensor queries of Section~\ref{sec:blocks}.

Let \(N=m\kappa\) and \([\ExtensionField:\Field]=\kappa\).
We commit to the fixed-basis packing \(\bm A\ext{\bm\beta}\)
of \(\bm A\in\Field^{m\times\kappa}\) and reduce each query
to one MLE opening. The verifier evaluates the resulting public
query by a recursion on small base-field matrices.
We then instantiate the underlying PCS from a base-field linear
code via Blaze's code-switching construction.

\subsubsection{Given a PCS over the Extension Field}

\medskip\noindent\textbf{Commitment and opening relation.}

Let \(m,\kappa\) be powers of two, with \(\kappa\geq2\), and write
\(\ell=\log_2m\).
Fix an \(\Field\)-basis
\(\ext{\bm\beta}=(\ext\beta_1,\ldots,\ext\beta_\kappa)^\top\)
of \(\ExtensionField\), with \(\ext\beta_1=1\), and let \(\pc\)
be a multilinear PCS for vectors in \(\ExtensionField^m\).
Extension-field elements are stored in this basis, so packing
and coordinate extraction only reorganize their base-field coordinates.
The schoolbook cost bounds below use a polynomial basis.

\paragraph{Setup and commitment.}
Run \(\pc.\Setup(1^\lambda,m)\) to obtain
\(\mathsf{pp}_{\pc}\), and include the fields, basis, and dimensions
\(N,m,\kappa\) in the public parameters.
For the witness matrix
\(\bm A=\concat{\bm a^{(1)},\ldots,\bm a^{(\kappa)}}
\in\Field^{m\times\kappa}\), pack each row into one
extension-field element, obtaining
\(\ext{\bm a}=\bm A\ext{\bm\beta}\in\ExtensionField^m\).
Compute the commitment
\(\mathsf{cm}=\pc.\Com(\mathsf{pp}_{\pc},\ext{\bm a};
\rho_{\pc})\), where \(\rho_{\pc}\) is the commitment
randomness, and retain \(\bm A\) and \(\rho_{\pc}\).
The fixed basis makes binding to \(\ext{\bm a}\) equivalent
to binding to \(\bm A\).

\paragraph{Opening relation.}
A query specifies column weights
\(\ext{\bm\rho}\in\ExtensionField^\kappa\), a tensor vector
\(\ext{\bm\xi}\in\ExtensionField^m\), and a claimed value
\(\ext v\in\ExtensionField\). The opening relation is:
\begin{equation}
    \ext v
    =\langle\bm A\ext{\bm\rho},\ext{\bm\xi}\rangle
    =\ext{\bm\xi}^{\top}\bm A\ext{\bm\rho}.
\label{eq:pcs-opening-relation}
\end{equation}
The row query is supplied as
\(\ext{\bm\xi}=\bigotimes_{j=1}^{\ell}\ext{\bm\xi}^{(j)}\),
where each factor is
\(\ext{\bm\xi}^{(j)}=(\ext\xi^{(j)}_0,
\ext\xi^{(j)}_1)^\top\in\ExtensionField^2\).
Ordinary MLE queries use the tensor spans of the corresponding
parts of the evaluation point for both \(\ext{\bm\rho}\) and
\(\ext{\bm\xi}\). Allowing arbitrary two-entry factors also
supports the coordinatewise products of tensor queries arising
in the weighted Hadamard reduction.

\medskip\noindent\textbf{Packed opening protocol.}

The commitment uses the fixed weights \(\ext{\bm\beta}\), whereas
the query uses \(\ext{\bm\rho}\). The prover first sends the
column evaluations. Coordinate decomposition and a random
projection then reduce their correctness to an inner-product
claim about the committed vector.

\paragraph{Reduce the column claims.}
Write the public query and the claimed column values in base-field
coordinates:
\[
    \ext{\bm\xi}=\bm\Xi\ext{\bm\beta},
    \qquad
    \ext{\bm u}^{\top}=\ext{\bm\beta}^{\top}\bm U,
    \qquad
    \bm\Xi\in\Field^{m\times\kappa},
    \quad
    \bm U\in\Field^{\kappa\times\kappa}.
\]
The honest prover computes \(\bm U=\bm\Xi^\top\bm A\) and sends
\(\ext{\bm u}=\bm U^\top\ext{\bm\beta}\).
The verifier extracts \(\bm U\) from this message and checks
\(\ext v=\ext{\bm u}^{\top}\ext{\bm\rho}\).
The matrix \(\bm\Xi\) is defined implicitly by the public query.

Since \(\bm A\) has base-field entries, the remaining claim
\(\ext{\bm u}^{\top}=\ext{\bm\xi}^{\top}\bm A\) is equivalent
to \(\bm U=\bm\Xi^\top\bm A\). Packing the rows of this
identity introduces the committed vector \(\ext{\bm a}\):
\begin{equation}
    \bm U=\bm\Xi^\top\bm A
    \quad\Longleftrightarrow\quad
    \bm U\ext{\bm\beta}
    =\bm\Xi^\top\ext{\bm a}.
\label{eq:pcs-encoded-matrix}
\end{equation}
After receiving \(\ext{\bm u}\), the verifier samples
\(\ext{\bm r}\sample\ExtensionField^\kappa\), independently of
the original query weights, and projects
Eq.~\eqref{eq:pcs-encoded-matrix} to one inner-product claim:
\begin{equation}
    \ext c
    =\ext{\bm r}^{\top}\bm U\ext{\bm\beta}
    =\langle\ext{\bm a},\bm\Xi\ext{\bm r}\rangle.
\label{eq:pcs-random-opening}
\end{equation}
The verifier computes \(\ext c\) by packing the rows of
\(\bm U\) and taking their inner product with \(\ext{\bm r}\).

Apply Fact~\ref{fact:ip} over \(\ExtensionField\) to
Eq.~\eqref{eq:pcs-random-opening}.
Let \(\ext{\bm y}\in\ExtensionField^\ell\) be its terminal point.
The prover's terminal values
\(\ext w_{\mathsf{com}},\ext w\in\ExtensionField\) give two claims:
\begin{equation}
    \ext w_{\mathsf{com}}
    =\langle\ext{\bm a},\tensorspan{\ext{\bm y}}\rangle,
    \qquad
    \ext w=\tensorspan{\ext{\bm y}}^{\top}\bm\Xi\ext{\bm r}.
\label{eq:pcs-terminal-claims}
\end{equation}
The first is an MLE opening of the committed vector.
The verifier checks the second by the following matrix recursion.

\paragraph{Evaluate the public query.}
Write \(\tensorspan{\ext{\bm y}}=\bm Y\ext{\bm\beta}\), where
\(\bm Y\in\Field^{m\times\kappa}\).
It suffices to compute the small matrix:
\begin{equation}
    \bm W=\bm Y^\top\bm\Xi\in\Field^{\kappa\times\kappa},
    \qquad
    \ext w=\ext{\bm\beta}^{\top}\bm W\ext{\bm r}.
\label{eq:pcs-query-matrix}
\end{equation}
The public tensor factors determine \(\bm Y\) and \(\bm\Xi\),
so the verifier can compute their product without forming either
length-\(m\) coordinate table.

For \(\ext z\in\ExtensionField\), let
\(\bm M_{\ext z}\in\Field^{\kappa\times\kappa}\) be the coordinate
matrix of multiplication by \(\ext z\), characterized by:
\begin{equation}
    \ext{\bm\beta}^{\top}\bm M_{\ext z}
    =\ext z\,\ext{\bm\beta}^{\top}.
\label{eq:pcs-multiplication-matrix}
\end{equation}
Thus left multiplication by \(\bm M_{\ext z}\) scales each
coordinate column by \(\ext z\) in \(\ExtensionField\), and right
multiplication by \(\bm M_{\ext z}^\top\) scales each coordinate row.
In particular, \(\bm M_{1-\ext z}=\bm I-\bm M_{\ext z}\).

Let \(\bm Y_j,\bm\Xi_j\) be the coordinate matrices for the
suffix factors \(j,\ldots,\ell\):
\[
    \bm Y_j\ext{\bm\beta}
    =\bigotimes_{i=j}^{\ell}
        \begin{pmatrix}1-\ext y_i\\ \ext y_i\end{pmatrix},
    \qquad
    \bm\Xi_j\ext{\bm\beta}
    =\bigotimes_{i=j}^{\ell}\ext{\bm\xi}^{(i)}.
\]
Here \(\bm Y_1=\bm Y\), \(\bm\Xi_1=\bm\Xi\), and
\(\bm Y_{\ell+1}=\bm\Xi_{\ell+1}=\bm e_1^\top\), where
\(\bm e_1=(1,0,\ldots,0)^\top\in\Field^\kappa\).
The two possible values of the first suffix bit give the block
decompositions:
\begin{equation}
\begin{aligned}
    \bm Y_j
    &=
    \begin{bmatrix}
        \bm Y_{j+1}\bm M_{1-\ext y_j}^\top\\
        \bm Y_{j+1}\bm M_{\ext y_j}^\top
    \end{bmatrix},\\
    \bm\Xi_j
    &=
    \begin{bmatrix}
        \bm\Xi_{j+1}\bm M_{\ext\xi^{(j)}_0}^\top\\
        \bm\Xi_{j+1}\bm M_{\ext\xi^{(j)}_1}^\top
    \end{bmatrix}.
\end{aligned}
\label{eq:pcs-query-blocks}
\end{equation}
Consequently, the verifier computes
\(\bm W_j=\bm Y_j^\top\bm\Xi_j\) by initializing
\(\bm W_{\ell+1}=\bm e_1\bm e_1^\top\) and processing
\(j=\ell,\ldots,1\):
\begin{equation}
\begin{aligned}
    \bm W_j
    &=
    \bm M_{1-\ext y_j}\bm W_{j+1}
        \bm M_{\ext\xi^{(j)}_0}^\top
    +\bm M_{\ext y_j}\bm W_{j+1}
        \bm M_{\ext\xi^{(j)}_1}^\top\\
    &=
    \bm W_{j+1}\bm M_{\ext\xi^{(j)}_0}^\top
    +\bm M_{\ext y_j}\bm W_{j+1}
        \bm M_{\ext\xi^{(j)}_1-\ext\xi^{(j)}_0}^\top.
\end{aligned}
\label{eq:pcs-query-recurrence}
\end{equation}
This recursion yields \(\bm W=\bm W_1\) without forming
the suffix coordinate matrices.
Each step of the second line uses at most \(3\kappa\)
extension-field multiplications, implemented by row and column
scaling without materializing \(\bm M_{\ext z}\).
The recursion takes \(O(\kappa\ell)\) extension-field
multiplications and \(O(\kappa^2\ell)\) base-field additions and
coordinate moves, using \(O(\kappa^2)\) base-field coordinates
of working memory. It applies in every characteristic.

The verifier packs the columns of \(\bm W\), computes
\(\ext{\bm\beta}^{\top}\bm W\ext{\bm r}\), and checks the second
identity in Eq.~\eqref{eq:pcs-query-matrix}.
This computation is exact and requires no additional prover message.

\paragraph{Terminal opening.}
Finally, the parties authenticate \(\ext w_{\mathsf{com}}\)
at \(\ext{\bm y}\) by invoking
\(\pc.\Open(\mathsf{pp}_{\pc},\mathsf{cm},
\ext w_{\mathsf{com}},\ext{\bm y}
:\ext{\bm a},\rho_{\pc})\).
The verifier accepts if the scalar checks, the inner-product
reduction, and the underlying opening accept.
Algorithm~\ref{alg:packed-pcs} in Appendix~\ref{appx:pseudocode}
summarizes the protocol.

\medskip\noindent\textbf{Security and complexity.}

\begin{proposition}[Packed polynomial commitment scheme]
\label{proposition:packed-pcs}
Let \(\pc\) bind the prover to a vector in \(\ExtensionField^m\),
with evaluation-soundness error \(\epsilon_{\mathsf{PCS}}\) for
adaptively chosen MLE claims about that vector.
The packed commitment is binding and the opening protocol is
complete, with evaluation-soundness error at most:
\begin{equation}
    \epsilon_{\mathsf{PCS}}
    +\frac{1+2\ell}{|\ExtensionField|}.
\label{eq:pcs-soundness}
\end{equation}
The opening reduction requires no additional lower bound on
\(|\Field|\).

Excluding the underlying PCS opening, the prover's arithmetic cost is:
\begin{equation}
    O(\kappa N)\,\bb
    +O(N)\,\be
    +O(m+\kappa)\,\ee.
\label{eq:pcs-prover-cost}
\end{equation}
The verifier's arithmetic cost is:
\begin{equation}
    O\bigl(\kappa^2(1+\ell)\bigr)\,\bb
    +O\bigl(\kappa(1+\ell)\bigr)\,\ee.
\label{eq:pcs-verifier-cost}
\end{equation}
The verifier bound includes base-field additions and coordinate
reorganization in the matrix recursion.
When \(m\geq\kappa\), schoolbook arithmetic gives
\(O(\kappa N)\) base-field operations for the prover.
Peak prover memory is \(O(N+\kappa^2)\) base-field coordinates,
or \(O(N)\) when \(m\geq\kappa\), excluding the workspace of the
underlying PCS.

With \(\ext v\) already part of the public statement,
prover-to-verifier communication, excluding the underlying
opening, consists of \(\kappa+2\ell+2\) extension-field elements.
The packed commitment inherits hiding from the underlying
commitment. The opening protocol reveals the column evaluations
\(\ext{\bm u}\), so this inheritance does not establish
zero knowledge of the opening transcript.
\end{proposition}

The proof is deferred to Appendix~\ref{appx:proof-pcs}.

\subsubsection{Given a Linear Code over the Base Field}

A base-field generator matrix commutes with basis packing and
therefore also encodes packed messages over \(\ExtensionField\).
We instantiate the preceding construction via Blaze's code-switching
protocol~\cite{eurocrypt/BrehmCFRRZ25}, with a BaseFold
code~\cite{crypto/ZeilbergerCF24} as the target, under Blaze's
source-encoding and proximity assumptions.

Let \(\Code:\Field^m\to\Field^n\) be a linear code with
generator matrix \(\bm G\in\Field^{n\times m}\).
Its \(\kappa\)-way interleaving encodes the columns of
\(\bm A\in\Field^{m\times\kappa}\) independently.
Packing the resulting codeword in the fixed basis gives:
\begin{equation}
\begin{aligned}
    \bm C&=\bm G\bm A\in\Field^{n\times\kappa},\\
    \ext{\bm c}&=\bm C\ext{\bm\beta}
       =\bm G(\bm A\ext{\bm\beta})
       =\Code_{\ExtensionField}(\ext{\bm a}).
\end{aligned}
\label{eq:pcs-code-packing}
\end{equation}
Here \(\Code_{\ExtensionField}\) is the scalar extension defined
by the same generator matrix; it has the same rate and relative
distance as \(\Code\).
Thus a row of \(\bm C\) represents one extension-field codeword
symbol. Committing to these rows supplies Blaze's source codeword
oracle for \(\ext{\bm a}=\bm A\ext{\bm\beta}\).

Apply Blaze's code-switching construction over \(\ExtensionField\)
to this oracle, switching to a BaseFold encoding of the same
length-\(m\) message \(\ext{\bm a}\).
The resulting MLE openings instantiate the underlying PCS \(\pc\)
in the preceding construction.
The packed opening reduction handles the challenge-dependent weights
\(\ext{\bm\rho}\), and code switching authenticates its terminal
fixed-basis opening. We include the costs and soundness error of
the chosen Blaze instantiation, including source-encoding
consistency, in the underlying PCS costs and
\(\epsilon_{\mathsf{PCS}}\), respectively.

If \(\Code\) is constant-rate and linear-time encodable, forming
\(\bm C\) costs \(O(N)\) base-field operations and its payload
occupies \(O(N\log_2|\Field|)\) bits, excluding Merkle-tree hashes.
The BaseFold stage works on \(m=N/\kappa\) extension-field message
symbols. Thus the source encoding retains the witness's small-field
representation, and the extension-field opening operates on the
shorter packed message.


\clearpage

\iffalse
\subsection{Complexity Bounds for the Hadamard Comparison}
\label{appx:hadamard-comparison}

Table~\ref{tab:hadamard-comparison} compares proofs of
\(\bm a\circ\bm b=\bm c\) for length-\(N\) vectors.
The constraint \(a_i b_i-c_i\) has degree two. Costs include equality
weights, the linear output check, witness-dependent preprocessing, and all reduction
rounds, but exclude commitments and openings. The derived upper bounds
use schoolbook arithmetic: multiplication in
\(\Field\times\Field\), \(\Field\times\ExtensionField\), and
\(\ExtensionField\times\ExtensionField\) costs \(O(1)\),
\(O(\kappa)\), and \(O(\kappa^2)\), respectively, where
\(\kappa=[\ExtensionField:\Field]\). RAM counts base-field elements
including inputs and read-only tables, not measured RSS.
Public field and basis preprocessing is separate; lookup-table
construction is accounted for below. Table accesses are charged
per base-field coordinate, as in Section~\ref{sec:notation}.

\smallskip\noindent\textbf{Optimized rounds and omitted terms.}
Let \(0\leq t\leq\log_2N\) count optimized or skipped rounds.
For small-field-first~\cite{iacr/BagadDT24,iacr/BagadDDT25,iacr/DaoDBDT26},
streamed prefix binding and a weighted \(3^t\)-entry quadratic grid give:
\begin{equation}
\begin{aligned}
\text{Arithmetic:}\quad&
O\!\left(\kappa N(3/2)^t+\kappa^2N/2^t+\kappa^2 3^t\right),\\
\text{RAM:}\quad& O\!\left(N+\kappa N/2^t+\kappa 3^t\right).
\end{aligned}
\label{eq:hadamard-sff-full-cost}
\end{equation}
The last terms account for grid contractions and storage.
For \(N\geq6^t\), they are absorbed by the corresponding
\(N/2^t\) terms, yielding the table entries.
For univariate skip~\cite{iacr/Gruen24}, including FFT interpolation,
weighting, binding, first-round processing, and the sumcheck tail gives:
\begin{equation}
\begin{aligned}
\text{Arithmetic:}\quad&
O\!\left(Nt+\kappa N+\kappa^2N/2^t+\kappa^2 2^t\right),\\
\text{RAM:}\quad& O\!\left(N+\kappa N/2^t+\kappa 2^t\right).
\end{aligned}
\label{eq:hadamard-skip-full-cost}
\end{equation}
Here \(N\geq4^t\) suffices to absorb the last terms.
The \(Nt\) term uses FFT interpolation on suitable subgroup
and off-domain evaluation sets. The protocol also admits general
interpolation sets~\cite[Section~5.1, footnote~18]{iacr/Gruen24},
with their interpolation cost accounted for separately. When supported,
\(2^t=\Theta(\kappa)\) and \(N=\Omega(\kappa^2)\) give
\(O(\kappa N)\) arithmetic and \(O(N)\) RAM.
For \(t\) fixed by \(\kappa\), the omitted terms are independent of
\(N\). If the absorption conditions fail, use the full bounds in
Eq.~\eqref{eq:hadamard-sff-full-cost} and
Eq.~\eqref{eq:hadamard-skip-full-cost}.
For degree-two constraints, the off-domain first-round values in
Gruen's protocol and the quotient coefficients in our Hadamard
reduction represent the same vanishing residual. Its known zeros
on the interpolation set determine the conversion between the two
representations~\cite[Section~6]{iacr/Gruen24}.

\smallskip\noindent\textbf{Flock's lookup-based skip.}
We re-cost the degree-two zerocheck of
Flock~\cite[Sections~4.2--4.4]{iacr/BunzRW26} with a conventional
sumcheck tail. Fix its byte-sized chunks and subfield
\(\mathbb F_{256}\subseteq\ExtensionField\), and let
\(3\leq t\leq7\) be the number of skipped variables.
Each row has \(2^t\) input bits and \(2^t\) off-domain samples.
Its LDE uses \(2^t/8\) lookups of length-\(2^t\) byte vectors,
so all rows cost \(O(N2^t)\) bit operations, including reads
and XORs. Byte-field products cost \(O(N)\).
In compatible \(\ExtensionField/\mathbb F_{256}\) coordinates,
weighted aggregation costs \(O(\kappa N)\), as do binding at
the skip challenge and the linear output check.
The conventional tail and its equality weights cost
\(O(\kappa^2N/2^t)\). Including parameter-only message
processing and fixed lookup tables gives:
\begin{equation}
\begin{aligned}
\text{Arithmetic:}\quad&
O\!\left(N2^t+\kappa N+\kappa^2N/2^t+\kappa^2 2^t\right),\\
\text{RAM:}\quad&
O\!\left(N+\kappa N/2^t+\kappa 2^t\right).
\end{aligned}
\label{eq:hadamard-flock-full-cost}
\end{equation}
The last terms are absorbed when \(N\geq4^t\).
This bound covers Flock's algebraic reduction to evaluation claims.
Its implementation uses \(t=6\), \(\kappa=128\),
a 16\,KiB compressed LDE table, and a 64\,KiB table for its
fixed equality weights. Compression improves table storage and
cache locality, not the asymptotic LDE operation count.
The admissible skip length is bounded by the byte-field domain,
independently of \(\kappa\).

\smallskip\noindent\textbf{Binary-tower instantiations.}
Assume \(N\geq\kappa\). The Dao--Thaler row pairs its reduction
with conventional sumcheck~\cite{iacr/DaoT24}.
The Liu--Zhang row uses every other level of a Wiedemann
tower~\cite{iacr/LiuZ25}, with mixed subfield/top-field multiplication
and streamed deterministic weights~\cite{iacr/DaoT24}.
RAM retains folded inputs and random-suffix weight tables.

\smallskip\noindent\textbf{Packing assumptions.}
The RMFE--Aurora bound~\cite{iacr/CascudoG21} includes sparse Boolean
R1CS and Reed--Solomon encoding. It assumes constant-rate packing and
codes, linear-space quasilinear FFTs, security \(O(\kappa)\), and a
packed instance large enough to absorb security padding.
Packed sumcheck~\cite{iacr/WeiWXZDWZ25} uses total packed length \(N\),
packing/repetition parameter \(r\), and \(|\Field|>3(2r-1)\).
Streaming forward differences retains \(O(r^2)\) first-round data.
This term remains in the table because \(r\) is independent of
\(\kappa\); it is absorbed if \(N\geq r^2\).
For the base-field \methodname{} row, \(|\Field|\geq\kappa\),
\(N=m\kappa\), \(m\geq\kappa\), and precomputed
bases give \(O(\kappa N)\) arithmetic and
\(O(N+\kappa^2)=O(N)\) RAM, including the linear output check.
All parameter choices must satisfy the target soundness bound.
In particular, \(r=\Theta(\kappa)\) is not assumed, and univariate
skip's error depends on \(2^t\).

\smallskip\noindent\textbf{Intermediate-field accounting.}
Let
\(t_{\mathbb K}=[\mathbb K:\Field]\) be the degree chosen in
Section~\ref{sec:hadamard}. Without lookup tables, multiplication in
\(\Field\times\mathbb K\) costs \(O(t_{\mathbb K})\)
base-field operations, whereas multiplication in \(\mathbb K\)
costs \(O(t_{\mathbb K}^2)\).
In Eq.~\eqref{eq:hadamard-small-field-aggregation}, computing
\(\bm A\bm P\) and \(\bm B\bm P\) therefore costs
\(O(N\kappa t_{\mathbb K})\), and their entrywise product costs
\(O(Nt_{\mathbb K}^2)\).
The contraction with \(\bm L\) uses
\(O(N\kappa/t_{\mathbb K})\) multiplications in \(\mathbb K\),
or \(O(N\kappa t_{\mathbb K})\) base-field operations.
Only the \(\kappa/t_{\mathbb K}\) aggregate coordinate tables
are interpolated. Dense interpolation and division on each cost
\(O(\kappa^2)\) operations in \(\mathbb K\), for a total of
\(O(\kappa^3t_{\mathbb K})\) base-field operations.
The column claims, subsequent inner-product reduction, and linear
output check cost \(O(N\kappa+\kappa^3)\).
Thus Eq.~\eqref{eq:hadamard-small-field-cost} includes the
online parameter-only work; \(N\geq\kappa^2\) gives the
untabulated bound \(O(N\kappa\log_q\kappa)\).

For the space bound, generate one column of \(\bm P\) at a time
using prefix and suffix products of the node differences, and
scan the original rows to accumulate its column of \(\bm U\).
With the public Lagrange denominators precomputed, generating all
columns costs \(O(\kappa^2)\) operations in \(\mathbb K\),
absorbed by the bounds above. The witness and \(\bm L\) use
\(O(N)\) base-field coordinates, and \(\bm U\) uses
\(O(\kappa^2)\). Interpolation basis polynomials can likewise
be generated one at a time from the public vanishing polynomial,
so no dense interpolation matrix is stored.
Fixed conversions between the tower coordinates and the original
commitment basis take at most \(O(N\kappa+\kappa^3)\) operations
and \(O(\kappa^2)\) additional coordinates.
The resulting RAM bound is \(O(N+\kappa^2)\), with the witness
still stored over \(\Field\).

\smallskip\noindent\textbf{Tabulated intermediate-field implementation.}
The second \methodname{} row uses
Eq.~\eqref{eq:hadamard-block-tabulated-cost}.
Choose \(\tau=\lfloor\log_q\kappa\rfloor\), so
\(q^\tau\leq\kappa\) and \(t_{\mathbb K}/\tau=O(1)\).
The block tables of Eq.~\eqref{eq:hadamard-block-table}
form all samples in \(O(N\kappa t_{\mathbb K}/\tau)
=O(N\kappa)\) base-field operations.
Each multiplication-table lookup reads and writes
\(O(t_{\mathbb K})\) coordinates. Thus the
\(O(N\kappa/t_{\mathbb K})\) products in the contraction
cost \(O(N\kappa)\), and aggregate interpolation and division
cost \(O(\kappa^3)\). Samplewise products and the remaining
reduction steps are absorbed in these bounds. Hence:
\begin{equation}
\begin{aligned}
T_{\mathrm{online}}&=O(N\kappa+\kappa^3),\\
T_{\mathrm{pre}}&=O(\kappa^3+t_{\mathbb K}^2q^{2t_{\mathbb K}}),\\
S_{\mathrm{total}}&=O(N+\kappa^3+t_{\mathbb K}q^{2t_{\mathbb K}}).
\end{aligned}
\label{eq:hadamard-tabulation-full-cost}
\end{equation}
The multiplication table has \(q^{2t_{\mathbb K}}\) entries;
schoolbook construction costs \(O(t_{\mathbb K}^2)\) per entry.
The block tables use \(O(\kappa^3)\) coordinates and operations:
precompute scalar multiples of the sample-map rows, then enumerate
input blocks by extending prefixes. Both tables depend only on
public field and interpolation parameters and are reused across
witnesses. Streaming witness rows leaves \(O(N+\kappa^2)\)
working coordinates. For \(N\geq\kappa^2\), the online time is
\(O(N\kappa)\), but the read-only tables remain in the RAM entry.
\fi
\section{Deferred Proofs}
\label{appx:proof}
\label{appx:auxiliary-proofs}

This appendix proves the claims in Sections~\ref{sec:blocks}
and~\ref{sec:together} and the packed PCS claim in
Appendix~\ref{appx:pcs}, using the notation of each statement.

\subsection{Proof of Proposition~\ref{proposition:ip-complexity}}
\label{appx:proof-ip}

\begin{proof}
\noindent\textbf{Completeness.}
Expand the Laurent polynomial in Eq.~\eqref{eq:ip-laurent-polynomial}:
\[
    f(X)=
    \sum_{\tau=-(\kappa-1)}^{\kappa-1}f_\tau X^\tau,
    \qquad
    f_\tau=
    \sum_{\substack{j,j'\in[\kappa]\\j-j'=\tau}}
       \langle\bm a^{(j)},\bm b^{(j')}\rangle.
\]
Its constant coefficient is \(f_0=\langle\bm a,\bm b\rangle\).
Thus the auxiliary polynomials in Algorithm~\ref{alg:ip}
satisfy Eq.~\eqref{eq:ip-correction-identity} for a correct
inner-product claim. Evaluation at \(\ext x\neq0\) gives the
length-\(m\) claim in
Eq.~\eqref{eq:packed-inner-product-check}.

The terminal values of Fact~\ref{fact:ip} have the matrix form:
\[
\begin{aligned}
    \ext v_1
      &=\tensorspan{\ext{\bm y}}^\top\bm A\ext{\bm\rho}_A,
    &\ext{\bm\rho}_A
      &=(1,\ext x,\ldots,\ext x^{\kappa-1})^\top,\\
    \ext v_2
      &=\tensorspan{\ext{\bm y}}^\top\bm B\ext{\bm\rho}_B,
    &\ext{\bm\rho}_B
      &=(1,\ext x^{-1},\ldots,\ext x^{-(\kappa-1)})^\top.
\end{aligned}
\]
Both have tensor-structured row weights and fit the
packed-opening relation in Eq.~\eqref{eq:pcs-opening-relation}.

\smallskip\noindent\textbf{Soundness.}
Suppose \(c\neq\langle\bm a,\bm b\rangle\), and fix the
degree-\((\kappa-2)\) polynomials \(h_+,h_-\) before sampling
\(\ext x\). The Laurent residual
\(f(X)-c-Xh_+(X)-X^{-1}h_-(X^{-1})\) has a nonzero
constant coefficient. Multiplying by
\(X^{\kappa-1}\) gives a nonzero polynomial of degree at most
\(2\kappa-2\), so the evaluation check accepts with probability
at most \((2\kappa-2)/(|\ExtensionField|-1)\).
Conditioned on a false evaluated claim, Fact~\ref{fact:ip}
contributes at most \(2\log_2m/|\ExtensionField|\).
A union bound gives the stated error, excluding the terminal
packed openings.

\smallskip\noindent\textbf{Efficiency.}
Each ordered pair of distinct columns contributes a
length-\(m\) base-field inner product to the correction
coefficients, for \(m\kappa(\kappa-1)=N(\kappa-1)\) base-field
multiplications. Given the powers of \(\ext x\) and
\(\ext x^{-1}\), packing both matrices uses \(2N\,\be\).
Generating these powers and evaluating the corrections costs
\(O(\kappa)\,\ee\); Fact~\ref{fact:ip} adds
\((4m+O(\log m))\,\ee\).
Thus the total arithmetic cost is:
\[
    O(N\kappa)\,\bb+O(N)\,\be+O(m+\kappa)\,\ee
       =O(N\kappa+\kappa^3)\,\bb.
\]
Streaming the cross-column inner products retains only
\(O(\kappa)\) base-field correction coefficients.
The witness matrices and packed vectors use \(O(N)\)
base-field coordinates, and the challenge powers use
\(O(\kappa^2)\). An in-place inner-product reduction therefore
has peak memory \(O(N+\kappa^2)\Field\).
\end{proof}

\subsection{Proof of Proposition~\ref{proposition:hadamard-complexity}}
\label{appx:proof-hadamard}

\begin{proof}
\noindent\textbf{Completeness.}
The honest column message is
\(\ext{\bm c}^{\top}
=\ext{\bm l}^{\top}(\bm A\circ\bm B)\), so the scalar check
\(\ext c=\langle\ext{\bm c},\ext{\bm r}\rangle\) accepts.
At every interpolation node, the residual satisfies:
\[
    f(\gamma_j)
      =\langle\ext{\bm l},
          \bm a^{(j)}\circ\bm b^{(j)}\rangle-\ext c_j
      =0,
    \qquad j\in[\kappa].
\]
Hence \(f=\Psi h\) for a polynomial \(h\) of degree at most
\(\kappa-2\). Evaluating this identity at \(\ext x\) gives:
\[
    \langle\ext{\bm a},\ext{\bm b}\circ\ext{\bm l}\rangle
      =\ext c_x+\Psi(\ext x)h(\ext x).
\]
The terminal values of Fact~\ref{fact:ip} satisfy:
\[
\begin{aligned}
    \ext v_1
      &=\tensorspan{\ext{\bm y}}^\top
          \bm A\ext{\bm\rho}(\ext x),\\
    \ext v_2
      &=\bigl(\ext{\bm l}\circ\tensorspan{\ext{\bm y}}\bigr)^\top
          \bm B\ext{\bm\rho}(\ext x).
\end{aligned}
\]
Writing the supplied row weight as
\(\ext{\bm l}=\bigotimes_{j=1}^{\log_2m}
(\ext l_{j,0},\ext l_{j,1})^\top\), its product with the MLE
weight factors as:
\[
    \ext{\bm l}\circ\tensorspan{\ext{\bm y}}
      =\bigotimes_{j=1}^{\log_2m}
        \begin{pmatrix}
            \ext l_{j,0}(1-\ext y_j)\\
            \ext l_{j,1}\ext y_j
        \end{pmatrix}.
\]
Both terminal claims fit Eq.~\eqref{eq:pcs-opening-relation},
with Lagrange column weights and tensor-structured row weights.

\smallskip\noindent\textbf{Soundness.}
Suppose the weighted Hadamard claim is false but the scalar
check accepts. Then the column message differs from
\(\ext{\bm l}^{\top}(\bm A\circ\bm B)\) in at least one
coordinate, and \(f(\gamma_j)\neq0\) for some \(j\).
For every degree-\((\kappa-2)\) polynomial \(h\) fixed before
\(\ext x\), the residual \(f-\Psi h\) is therefore nonzero
and has degree at most \(2\kappa-2\).
A random \(\ext x\in\ExtensionField\setminus\mathcal H\)
makes this residual vanish with probability at most
\((2\kappa-2)/(|\ExtensionField|-\kappa)\).
Fact~\ref{fact:ip} contributes at most
\(2\log_2m/|\ExtensionField|\), giving:
\[
    \epsilon_{\mathsf{Had}}
      \leq
      \frac{2\kappa-2}{|\ExtensionField|-\kappa}
      +\frac{2\log_2m}{|\ExtensionField|}
      \leq
      \frac{2\kappa-2}{|\ExtensionField|-\kappa}
      +\frac{2\log_2N}{|\ExtensionField|}.
\]
The first inequality bounds Algorithm~\ref{alg:hadamard};
the second gives the stated bound. Terminal packed openings
are accounted for separately.
The argument applies unchanged when the interpolation nodes
lie in \(\mathbb K\subseteq\ExtensionField\).

\smallskip\noindent\textbf{Efficiency over the base field.}
First suppose \(\mathbb K=\Field\).
Computing \(\ext{\bm c}\) uses \(O(N)\,\bb+O(N)\,\be\).
At each of the \(\kappa-1\) additional interpolation nodes,
evaluating both witness encodings takes \(O(N)\,\bb\).
Pointwise multiplication and contraction with \(\ext{\bm l}\)
then take \(O(m)\,\bb+O(m)\,\be\).
Over all nodes, these operations cost
\(O(N\kappa)\,\bb+O(N)\,\be\).
Evaluating \(g_C\), interpolating the residual, and dividing by
\(\Psi\) contribute \(O(\kappa^2)\,\be\).

After \(\ext x\) is sampled, packing both matrices costs
\(O(N)\,\be\).
Expanding the row factors, forming
\(\ext{\bm b}\circ\ext{\bm l}\), and running Fact~\ref{fact:ip}
cost \(O(m)\,\ee\).
The terminal row factors require only componentwise
multiplication of their two-entry representations.
With precomputed Lagrange coefficients, the public query
and polynomial evaluations cost
\(O(\kappa^2)\,\be+O(\kappa)\,\ee\).
The total is therefore:
\[
    O(N\kappa)\,\bb
      +O(N+\kappa^2)\,\be
      +O(m+\kappa)\,\ee
      =O(N\kappa+\kappa^3)\,\bb.
\]

\smallskip\noindent\textbf{Intermediate-field tables.}
Let \([\mathbb K:\Field]=t>1\), with \(t\mid\kappa\) and
\(|\mathbb K|\geq2\kappa-1\).
Use the addition and multiplication tables specified in
Section~\ref{sec:hadamard}, charging each lookup as \(O(1)\).
Evaluating the witness encodings at all additional nodes
requires \(O(N\kappa)\) table operations.
In \(\mathbb K\)-coordinates, each extension-field element
has length \(\kappa/t\), so weighting and accumulating the
\(O(N)\) local products uses \(O(N\kappa/t)\) lookups.
The \(O(\kappa^2)\) mixed operations for \(g_C\), interpolation,
and division cost \(O(\kappa^3/t)\).
The remaining packing and inner-product steps retain their
\(O(N\kappa+\kappa^3)\,\bb\) bound. Thus the same online
time bound holds, excluding construction of the fixed tables.

\smallskip\noindent\textbf{Memory.}
Processing rows and sample points sequentially retains
\(O(\kappa)\) extension-field residual values, followed by
\(O(\kappa)\) extension-field coefficients after interpolation.
The packed vectors and inner-product reduction use \(O(m)\)
extension-field elements. Including the witness matrices,
this is \(O(N+\kappa^2)\) base-field coordinates.
The fixed interpolation data use \(O(\kappa^2t)\)
coordinates, and the tables use \(O(tq^{2t})\) when \(t>1\);
no field tables are needed when \(t=1\).
The peak memory is therefore
\(O\bigl(N+\kappa^2t+(t-1)q^{2t}\bigr)\Field\).
\end{proof}

\subsection{Proof of Proposition~\ref{proposition:permutation-complexity}}
\label{appx:proof-permutation}

\begin{proof}
\noindent\textbf{Completeness.}
The coefficient tree in
Eq.~\eqref{eq:packed-permutation-product-tree} computes the
products of \(2^\tau\) shifted witness entries in each column.
Evaluating its last layer at \(\ext z\) gives \(m\)
extension-field entries, whose product is
\(\mathsf{GrandProd}(\ext z\bm1+\bm a)\).
The two grand products agree whenever \(\bm a\) and \(\bm b\)
contain the same multiset of entries.

For a transition from level \(\tau+1\) to level \(\tau\),
Lagrange interpolation gives:
\[
    \bm g_{\tau,U}(Z,\gamma_j)
      \circ\bm g_{\tau,D}(Z,\gamma_j)
       =\bm C_{\tau+1}(Z)_{*,j},
    \qquad j\in[\kappa].
\]
For the honest column claims
\(\ext{\bm u}_{\tau+1}^\top
=\ext{\bm l}_{\tau+1}^\top\bm C_{\tau+1}(\ext z)\),
the residual \(f(X)\) therefore vanishes on \(\mathcal H\).
Since its degree is at most \(2\kappa-2\), it satisfies
\(f(X)=\Psi(X)h(X)\) with \(\deg h\leq\kappa-2\).
At an additional node, coefficient aggregation commutes with
evaluation at \(\ext z\):
\[
\begin{aligned}
    F_{\tau+1,j}(\ext z)
      &=\sum_{\nu=0}^{2^{\tau+1}}
        \langle\ext{\bm l}_{\tau+1},
                 \bm p_{\tau+1,j,\nu}\rangle\ext z^\nu\\
      &=\left\langle\ext{\bm l}_{\tau+1},
             \bm g_{\tau,U}(\ext z,\gamma'_j)
               \circ\bm g_{\tau,D}(\ext z,\gamma'_j)
        \right\rangle.
\end{aligned}
\]
Algorithm~\ref{alg:permutation} therefore recovers \(h\) from
\(\kappa-1\) samples. Reencoding \(\bm C_\tau(Z)\) at each
level keeps the \(X\)-degree of \(\bm g_\tau(Z,X)\) at most
\(\kappa-1\).

Each evaluated relation is proved by weighted sumcheck.
For length-\(n\) child vectors and a tensor weight
\(\ext{\bm l}\), the sumcheck polynomial is:
\begin{equation}
    \widetilde{\ext{\bm l}}(\bm Y)\,
    \widetilde{\ext{\bm u}}(\bm Y)\,
    \widetilde{\ext{\bm v}}(\bm Y).
\label{eq:permutation-hadamard-as-ip}
\end{equation}
It has individual degree at most three.
The two child claims are evaluated at the same point
\(\ext{\bm y}\). Their affine combination with a fresh
\(\ext y_0\) is the MLE claim on the concatenated parent
vector at \(\ext y_0\mathbin\|\ext{\bm y}\).
This preserves the recursive matrix claim
\(\ext v_\tau=\ext{\bm l}_\tau^\top
\bm C_\tau(\ext z)\ext{\bm r}_\tau\).
At level zero, the row MLE weights and the column Lagrange
weights both sum to one. Removing the shift leaves:
\[
    \ext v_0-\ext z
       =\ext{\bm l}_0^\top\bm A\ext{\bm r}_0.
\]
This is the packed-opening interface of
Eq.~\eqref{eq:pcs-opening-relation}. The same argument applies
to \(\bm B\).

\smallskip\noindent\textbf{Soundness.}
For a weighted sumcheck on length-\(n\) vectors,
Eq.~\eqref{eq:permutation-hadamard-as-ip} contributes at most
\(3\log_2n/|\ExtensionField|\).
If at least one of the two terminal child claims is false,
their errors form a nonzero affine polynomial in the fresh
merge challenge \(\ext y_0\). The merged claim is then correct
with probability at most \(1/|\ExtensionField|\).
A transition thus preserves a false claim except with
probability at most \((3\log_2n+1)/|\ExtensionField|\).
The extension-field tree has child lengths
\(1,2,\ldots,m/2\). Summing its errors gives:
\[
    \sum_{i=0}^{\log_2m-1}
       \frac{3i+1}{|\ExtensionField|}
       =\frac{\frac32\log_2m(\log_2m-1)+\log_2m}
                    {|\ExtensionField|}.
\]

At coefficient level \(\tau+1\), a false incoming matrix
claim either fails the scalar check against
\(\ext{\bm u}_{\tau+1}\), or makes at least one column claim
false. In the latter case \(f(X)\) does not vanish on all of
\(\mathcal H\), and \(f(X)-\Psi(X)h(X)\) is nonzero for every
claimed \(h\) of degree at most \(\kappa-2\).
Since the message is fixed before \(\ext x\) is sampled,
this polynomial vanishes at \(\ext x\) with probability at
most \((2\kappa-2)/(|\ExtensionField|-\kappa)\).
The subsequent weighted sumcheck and merge have child length
\(m/2^{\tau+1}\) and error at most
\((3\log_2(m/2^{\tau+1})+1)/|\ExtensionField|\).
Summing over \(\tau=0,\ldots,\log_2\kappa-1\) and adding the
extension-field-tree error gives the stated
\(\epsilon_{\mathsf{GP}}\). This bound is conditional on
authentication of the terminal packed claim.

If the committed vectors are not permutations of one another,
their fingerprint difference is the nonzero polynomial:
\[
    \prod_{i=1}^N(Z+a_i)-\prod_{i=1}^N(Z+b_i).
\]
Its degree is at most \(N-1\), so the true fingerprints
coincide with probability at most
\(\min\{1,(N-1)/|\ExtensionField|\}\).
Otherwise, at least one claimed grand product is false.
A union bound over the two reductions and their terminal
authentication gives:
\[
    \epsilon_{\mathsf{Perm}}
       \leq\min\left\{1,\frac{N-1}{|\ExtensionField|}\right\}
          +2\epsilon_{\mathsf{GP}}+\epsilon_{\mathsf{PCS}}.
\]

\smallskip\noindent\textbf{Efficiency over \(\Field\).}
At level \(\tau\in[\log_2\kappa]\), the coefficient matrix
of \(\bm C_{\tau-1}(Z)\) has \(O(m)\) row-degree pairs.
Evaluating its \(\kappa\)-column encodings at \(\kappa-1\)
additional nodes costs \(O(N\kappa)\,\bb\).
Multiplying the resulting degree-\(2^{\tau-1}\) polynomials
costs \(O(m\kappa2^\tau)\,\bb\).
These products, including construction of the coefficient
tree, sum to \(O(N\kappa)\,\bb\).

The off-domain products and column claims contain \(O(N)\)
coefficients per level, whose aggregation costs \(O(N)\,\be\).
There are \(O(\kappa)\) aggregated polynomials of degree
at most \(2^\tau\), so their subsequent evaluation costs:
\[
    \sum_{\tau=1}^{\log_2\kappa}
        O(\kappa2^\tau)\,\ee
       =O(\kappa^2)\,\ee.
\]
The shared powers of \(\ext z\) require
\(O(\kappa)\,\ee\) once.
Forming the vectors for the short weighted sumchecks costs
\(O(N)\,\be+O(m)\,\ee\) per coefficient level.
Their total sumcheck length, together with the
extension-field tree, is \(O(m)\).
Interpolation and public queries contribute
\(O(\kappa^3)\,\bb\) per level.
Using \(m\geq\kappa\), the total is:
\[
    O(N\kappa\log_2\kappa+\kappa^4)\,\bb.
\]
The \(\kappa^4\) term accounts for evaluation of the
aggregated extension-field polynomials.

\smallskip\noindent\textbf{Intermediate-field implementation.}
For \(q<2\kappa-1\), use the addition and multiplication tables
of Section~\ref{sec:hadamard} over \(\mathbb K\), with the same
unit-cost lookup convention.
Precomputed Lagrange evaluations allow the
\(O(N\kappa)\) encoding operations at each level to be
performed over \(\mathbb K\).
The row-polynomial products likewise use
\(O(N\kappa)\) table operations in total.
The coefficients of \(\bm C_\tau(Z)\) remain in \(\Field\);
only the off-domain encodings and products use \(\mathbb K\).

In a compatible tower basis, a mixed
\(\mathbb K\times\ExtensionField\) product and its
accumulation cost \(O(\kappa)\) in base-field coordinates.
Schoolbook multiplication over \(\mathbb K\), using the
tables for coordinate products, costs
\(O(\kappa^2/t)\) per extension-field multiplication.
The preceding count becomes
\(O\bigl(N\kappa\log_2\kappa+\kappa^4/t\bigr)\,\bb\),
which implies the coarser bound in the proposition.

\smallskip\noindent\textbf{Space.}
One coefficient-tree level contains \(O(N)\) base-field
coefficients. The levels can be reconstructed from
\(\bm C_0(Z)\) during the reverse traversal:
reconstructing level \(\tau\) costs \(O(N2^\tau)\,\bb\),
and these costs sum to \(O(N\kappa)\,\bb\).
Thus the prover need not retain all coefficient levels.
Processing one additional node at a time uses
\(O(mt)\subseteq O(N)\) coordinates for its row polynomials.
The aggregated polynomials use \(O(\kappa^2)\)
extension-field coefficients, or \(O(\kappa^3)\)
base-field coordinates. The extension-field tree and
sumcheck tables use \(O(m\kappa)=O(N)\) coordinates.
Including interpolation data and the intermediate-field
tables, peak memory is \(O(N+\kappa^3+tq^{2t})\)
base-field coordinates.
When \(q\geq2\kappa-1\), no field tables are needed, giving
\(O(N+\kappa^3)\) base-field coordinates.
The time bounds exclude preprocessing, commitments, and
terminal openings.
\end{proof}

\subsection{Proof of Proposition~\ref{proposition:packed-pcs}}
\label{appx:proof-pcs}

\begin{proof}
\noindent\textbf{Binding and completeness.}
Each entry of \(\bm A\ext{\bm\beta}\) uniquely determines
the corresponding row of \(\bm A\). This injective encoding
preserves binding and, when provided by the underlying PCS,
commitment hiding.

For an honest prover, \(\bm U=\bm\Xi^\top\bm A\), so
Eq.~\eqref{eq:pcs-encoded-matrix} and its random projection hold.
We show that the verifier also computes the public terminal value
exactly.

By Eq.~\eqref{eq:pcs-multiplication-matrix}, multiplication by
\(\ext z\) acts on a column of base-field coordinates through
\(\bm M_{\ext z}\). This map is \(\Field\)-linear in \(\ext z\),
and \(\bm M_{1}=\bm I_\kappa\).
We prove by backward induction that
\(\bm W_j=\bm Y_j^\top\bm\Xi_j\).
Since \(\ext\beta_1=1\), the empty suffix has coordinate row
\(\bm e_1^\top\), giving
\(\bm Y_{\ell+1}^\top\bm\Xi_{\ell+1}
=\bm e_1\bm e_1^\top=\bm W_{\ell+1}\).
For the induction step, apply the block identities in
Eq.~\eqref{eq:pcs-query-blocks} and the induction hypothesis:
\[
\begin{aligned}
    \bm Y_j^\top\bm\Xi_j
    &=\bm M_{1-\ext y_j}\bm W_{j+1}
        \bm M_{\ext\xi^{(j)}_0}^\top\\
    &\quad{}+\bm M_{\ext y_j}\bm W_{j+1}
        \bm M_{\ext\xi^{(j)}_1}^\top\\
    &=\bm W_{j+1}\bm M_{\ext\xi^{(j)}_0}^\top\\
    &\quad{}+\bm M_{\ext y_j}\bm W_{j+1}
        \bm M_{\ext\xi^{(j)}_1-\ext\xi^{(j)}_0}^\top\\
    &=\bm W_j.
\end{aligned}
\]
The second equality uses linearity of \(\bm M\); the last
is Eq.~\eqref{eq:pcs-query-recurrence}.
Taking \(j=1\) gives \(\bm W=\bm W_1=\bm Y^\top\bm\Xi\)
and hence:
\[
    \ext{\bm\beta}^\top\bm W\ext{\bm r}
    =(\bm Y\ext{\bm\beta})^\top\bm\Xi\ext{\bm r}
    =\tensorspan{\ext{\bm y}}^\top\bm\Xi\ext{\bm r}.
\]
Thus Eq.~\eqref{eq:pcs-query-matrix} checks the query-side
terminal value exactly. Completeness follows from
Fact~\ref{fact:ip} and the underlying PCS opening.
For \(m=1\), the initialization suffices, with no recurrence
or sumcheck rounds.

\smallskip\noindent\textbf{Soundness.}
Fix the matrix \(\bm A\) bound by the commitment and suppose that
Eq.~\eqref{eq:pcs-opening-relation} is false.
If the first scalar check accepts, then
\(\ext{\bm u}^{\top}\neq\ext{\bm\xi}^{\top}\bm A\), hence
\(\bm U\neq\bm\Xi^\top\bm A\).
Injectivity of basis packing makes
\((\bm U-\bm\Xi^\top\bm A)\ext{\bm\beta}\) a nonzero
column of \(\ExtensionField^\kappa\), fixed before
\(\ext{\bm r}\) is sampled. Its inner product with the uniform
\(\ext{\bm r}\) vanishes with probability
\(1/|\ExtensionField|\). Outside this event,
Eq.~\eqref{eq:pcs-random-opening} is false.

Except with probability \(2\ell/|\ExtensionField|\),
Fact~\ref{fact:ip} then leaves a false terminal claim in
Eq.~\eqref{eq:pcs-terminal-claims}.
The recurrence computes \(\bm W\) from the public tensor factors
and the terminal point, so a false query-side claim is rejected
deterministically. If that check passes, the witness-side claim
must be false. The underlying PCS accepts this adaptively chosen
false MLE claim with probability at most
\(\epsilon_{\mathsf{PCS}}\).
A union bound gives Eq.~\eqref{eq:pcs-soundness}.

\smallskip\noindent\textbf{Prover complexity.}
Expanding \(\ext{\bm\xi}\) and reading its basis coordinates
gives \(\bm\Xi\) in \(O(m)\) extension-field operations.
The matrix product \(\bm U=\bm\Xi^\top\bm A\) takes
\(\kappa N\) base-field multiplications. Its columns give the
coordinates of \(\ext{\bm u}^{\top}=\ext{\bm\beta}^{\top}\bm U\),
so packing requires only a change of representation.

Computing \(\bm\Xi\ext{\bm r}\) takes \(N\) mixed-field
multiplications; the extension-field inner-product reduction takes
\(4m+O(\ell)\) extension-field multiplications.
Packing the rows of \(\bm U\) and taking their inner product
with \(\ext{\bm r}\) computes \(\ext c\) using
\(\kappa\) extension-field multiplications.
The prover need not construct \(\bm Y\) or \(\bm W\).
These costs give Eq.~\eqref{eq:pcs-prover-cost}.
Under schoolbook arithmetic they are
\(O(\kappa N+\kappa^3)\,\bb\), which is
\(O(\kappa N)\,\bb\) when \(m\geq\kappa\).

The matrices \(\bm A,\bm\Xi\) occupy \(O(N)\) base-field
coordinates, and \(\bm U\) and the length-\(\kappa\)
extension-field vectors occupy \(O(\kappa^2)\).
A constant number of length-\(m\) extension-field buffers also
occupy \(O(m\kappa)=O(N)\) base-field coordinates.
Thus peak prover memory, excluding the underlying PCS workspace,
is \(O(N+\kappa^2)\), or \(O(N)\) when \(m\geq\kappa\).

\smallskip\noindent\textbf{Verifier complexity and communication.}
The verifier implements Eq.~\eqref{eq:pcs-query-recurrence}
without forming the multiplication matrices.
Right multiplication by \(\bm M_{\ext z}^\top\) scales each
row, interpreted in the fixed basis as an extension-field element,
by \(\ext z\). Left multiplication by \(\bm M_{\ext z}\)
scales each column in the same way. Each scaling therefore costs
\(\kappa\) extension-field multiplications.
The two terms in each recurrence step require three such scalings,
followed by \(\kappa^2\) base-field additions.
Over \(\ell\) steps, the cost is at most
\(3\kappa\ell\) extension-field multiplications and
\(O(\kappa^2\ell)\) base-field operations, including coordinate
rearrangement. A constant number of intermediate matrices use
\(O(\kappa^2)\) additional base-field coordinates.

The initial scalar check costs \(\kappa\) extension-field
multiplications. To compute \(\ext c\) and the expected terminal
value \(\ext w\), the verifier packs the rows of \(\bm U\) and
the columns of \(\bm W\), respectively, and takes one
length-\(\kappa\) extension-field inner product per target.
Packing is coordinate reinterpretation, so these targets
cost \(O(\kappa)\) extension-field multiplications and
\(O(\kappa^2)\) coordinate operations in total.
The inner-product verifier costs \(O(1+\ell)\) extension-field
operations, including its terminal scalar check.
This gives Eq.~\eqref{eq:pcs-verifier-cost}.

The prover sends \(\kappa\) extension-field elements for
\(\ext{\bm u}\), followed by the \(2\ell+2\) extension-field
elements of Fact~\ref{fact:ip}, including both terminal values.
The verifier computes \(\bm W\) locally, so no matrix message or
base-field proof is sent. Excluding the underlying opening, total
prover-to-verifier communication is therefore
\(\kappa+2\ell+2\) extension-field elements.
\end{proof}

\subsection{Proof of Lemma~\ref{eq:permutation}}
\label{appx:proof-copy-constraint}

\begin{proof}
The radix-\(\underline q\) representations of \(0,\ldots,L-1\)
use at most \(\lceil\log_{\underline q}L\rceil\) digits.
The digit injection preserves their distinctness.
Since \(\theta_1,\ldots,\theta_d\) are linearly independent
over \(\Field\), and the last coordinate is reserved for the
witness value:
\[
    a_i=b_j
    \quad\Longleftrightarrow\quad
    \bigl(\sigma(i)=j\ \text{and}\ w_i=w_j\bigr).
\]
Every label occurs exactly once in each operand.
Thus equality of the multisets forces
\(a_i=b_{\sigma(i)}\), and hence \(w_i=w_{\sigma(i)}\),
for all \(i\in[L]\).
Conversely, these copy equalities give
\(a_i=b_{\sigma(i)}\), proving the permutation relation.
The encoding is deterministic and injective, so it adds no
soundness error. Appending the same zeros to both operands
preserves the equivalence by multiset cancellation.
\end{proof}

\subsection{Proof of Proposition~\ref{proposition:position-encoding-evaluation}}
\label{appx:proof-position-encoding}

\begin{proof}
For the chosen digit injection, each label is the sum of the
contributions of its nonzero position bits.
The recursion and interval decomposition in
Appendix~\ref{appx:mleiota} therefore give the ordinary MLE
evaluation in \(O(\log L)\) operations.

For a packed query, split the position bits into column bits
and row bits in column-major order.
The unpadded label is the sum of a column contribution
\(\bm\iota_{\mathsf{col}}\) and a row contribution
\(\bm\iota_{\mathsf{row}}\).
For arbitrary column weights \(\ext{\bm r}\) and tensor row
weights \(\ext{\bm l}\), the full-table contraction is:
\[
    \left\langle\bm\iota_{\mathsf{col}},\ext{\bm r}\right\rangle
       \left(\bm1^\top\ext{\bm l}\right)
    +
    \left(\bm1^\top\ext{\bm r}\right)
       \left\langle\bm\iota_{\mathsf{row}},\ext{\bm l}\right\rangle.
\]
Recursive doubling generates all \(\kappa_{\mathbb K}\)
column contributions in \(O(\kappa_{\mathbb K})\) additions.
For each row-bit suffix, compute both its total tensor weight
and its label-weighted sum.
Splitting off one two-entry tensor factor updates both
quantities in constant time, including when its entries do
not sum to one.

Zero padding keeps the first \(L\) positions.
In column-major order, these consist of complete columns
and at most one partially filled column.
All complete columns use the same two full-row contractions.
For the partial column, the binary interval decomposition
computes the two corresponding prefix contractions using the
precomputed suffix quantities.
One scan of the prefix length costs \(O(\log L)\) operations.
Combining the column contributions then costs
\(O(\kappa_{\mathbb K})\) operations.
The total is \(O(\kappa_{\mathbb K}+\log L)\), with the same
bound on working field elements and no division.
\end{proof}

\subsection{Proof of Theorem~\ref{theorem:single-instance-packing}}
\label{appx:proof-single-instance}

\begin{proof}
\noindent\textbf{Completeness.}
For a satisfying assignment, the five scalar claims in
Eq.~\eqref{eq:plonkWeightedHadamard1} and
Eq.~\eqref{eq:plonkWeightedHadamard2} sum to zero.
Four packed Hadamard proofs reduce the selector--wire product
claims to evaluations of the selectors and
\(\bm w_L,\bm w_R,\bm w_M,\bm w_O\).
The fifth authenticates the remaining evaluation of
\(\bm w_M=\bm w_L\circ\bm w_R\).
All terminal queries have arbitrary column weights and tensor
row weights, as required by Section~\ref{sec:pcs}.

Lemma~\ref{eq:permutation} proves the copy equivalence.
The two operands have length \(L=3N+n_{\mathsf{io}}\leq4N\);
append zeros to reach length \(4N\).
Their grand products then share the padding factor
\(\ext z^{4N-L}\), which is nonzero because
\(\ext z\sample\ExtensionField\setminus\{0\}\).
The two grand-product reductions therefore prove exactly the
original permutation relation.
Their packing width \(\kappa_{\mathbb K}=\kappa/d\) and
row count \(4N/\kappa_{\mathbb K}=4md\) are powers of two,
with \(4md\geq\kappa_{\mathbb K}\geq2\).
Thus the parameters satisfy the interface of
Section~\ref{sec:permutation}, and every honest check accepts.

\smallskip\noindent\textbf{Terminal copy queries.}
Each copy terminal is a weighted evaluation of a padded operand,
with \(\kappa_{\mathbb K}\) arbitrary column weights and
\(4N/\kappa_{\mathbb K}\) tensor row weights.
By linearity, the two terminals involve
\(\bm\varsigma,\bm\iota\), and \(\bm w\).
Restricting to one length-\(N\) wire segment fixes the first
two position bits.
The last \(\log_2m\) bits remain tensor row factors.
Absorb all other free factors into the column weights to obtain
a length-\(\kappa\) column query on the original
\(m\times\kappa\) wire or digit matrix.
This changes the query representation, not the committed data.

Proposition~\ref{proposition:position-encoding-evaluation}
computes the identity-label contribution in
\(O(\kappa_{\mathbb K}+\log L)\) extension-field operations.
The verifier evaluates the public-I/O values and their permuted
labels directly. Traversing their contiguous interval and
reusing tensor prefixes costs
\(O(\kappa_{\mathbb K}+n_{\mathsf{io}}+\log L)\) operations.
No commitment to either contribution is needed.

\smallskip\noindent\textbf{Batched terminal openings.}
Write \(h=\lceil\log_{\underline q}L\rceil\).
The three fixed-basis wire packings form one witness vector
of logical length \(3m\), committed before any challenge.
The public vector concatenates the five selector packings
and the three wire-segment restrictions of each of the \(h\)
permuted-position digit tables.
Its logical length is \(m_{\mathsf{pub}}=(5+3h)m\).
The public commitment, circuit tables, and packing polynomials
are preprocessed.

The gate argument leaves five selector and five wire queries.
The two copy terminals add six wire queries and \(3h\)
permuted-digit queries, giving \(16+3h\) packed queries.
Apply the column-projection step of Section~\ref{sec:pcs}
to each query and check its claimed scalar value.
Fix all column messages before sampling independent projection
vectors for the different queries.
Within each group, zero-extend the projected queries to their
committed blocks and sum the resulting inner-product claims.
One invocation of Fact~\ref{fact:ip} per group reduces their
sum to one MLE opening of the concatenated vector.
The verifier evaluates the combined public query exactly,
using the recurrence of Section~\ref{sec:pcs} on each block
and retaining its block-prefix weight.
Thus there are two underlying PCS openings.

If a column message is false, its coordinate error satisfies
\((\bm U-\bm\Xi^\top\bm A)\ext{\bm\beta}\neq\bm0\).
Consequently, a group containing a false claim has a nonzero
linear error polynomial in its independent projection vectors.
It vanishes with probability at most
\(1/|\ExtensionField|\) per group.
The two inner-product reductions add
\(O(\log N)/|\ExtensionField|\) error.
Here the padded group lengths have logarithm \(O(\log N)\),
since \(h<d\leq\kappa/2\).

\smallskip\noindent\textbf{Soundness.}
Condition on binding of the witness commitment.
Its fixed-basis packing determines all three wire vectors
and the virtual product before any challenge.
The selectors and position tables are fixed at setup.

A nonzero gate residual vanishes at the random MLE point with
probability at most \(\log_2N/|\ExtensionField|\).
Otherwise an accepting scalar sum contains a false claim.
Acceptance then requires an error in one of the five Hadamard
reductions or in terminal authentication.
For the nested product claim, its value is fixed before the
fresh challenges of the fifth reduction.

If a copy constraint fails, the unpadded fingerprint difference
is nonzero and has degree at most \(L-1\).
At the nonzero challenge \(\ext z\), its padded version vanishes
if and only if the unpadded difference does.
The fingerprint error is therefore at most
\((L-1)/(|\ExtensionField|-1)\).
Let \(\epsilon_{\mathsf{GP}}\) be the bound of
Proposition~\ref{proposition:permutation-complexity},
instantiated with length \(4N\), packing width
\(\kappa_{\mathbb K}\), and base field \(\mathbb K\).
A union bound gives:
\[
\begin{aligned}
    \epsilon_{\mathsf{single}}\leq{}&
    \epsilon_{\mathsf{PCS}}
    +\frac{L-1}{|\ExtensionField|-1}
    +2\epsilon_{\mathsf{GP}}\\
    &+5\left(
        \frac{2\kappa-2}{|\ExtensionField|-\kappa}
        +\frac{2\log_2N}{|\ExtensionField|}
      \right)
    +O\left(\frac{\log N}{|\ExtensionField|}\right).
\end{aligned}
\]
The last term includes the initial gate projection and the
two batched opening reductions.
The grand-product contribution is
\(O((\kappa_{\mathbb K}\log_2\kappa_{\mathbb K}
+(\log_2N)^2)/|\ExtensionField|)\).
Since \(N=m\kappa\geq\kappa^2\), \(L\leq4N\), and
\((\log_2N)^2=O(N)\), this proves
Eq.~\eqref{eq:together-soundness}.
The term \(\epsilon_{\mathsf{PCS}}\) includes joint binding
and the two adaptive underlying openings.

\smallskip\noindent\textbf{Prover time.}
The five Hadamard reductions cost \(O(N\kappa)\bb\).
For \(q\geq2\kappa-1\), interpolation uses base-field nodes.
If \(\kappa<q<2\kappa-1\), the intermediate-field variant
uses degree-two interpolation-field tables; its online cost
is unchanged and its table storage is \(O(q^4)=O(\kappa^4)\).
The power-of-two \(\kappa\) permits this degree-two subfield.
The \(16+3h\) column projections and the two grouped
inner-product reductions cost \(O(N\kappa h)\bb\).

Choose polynomial bases along \(\Field\subset\mathbb K
\subset\ExtensionField\) for the copy computation.
A \(\mathbb K\)-multiplication costs \(O(d^2)\bb\).
The two grand-product reductions therefore cost:
\[
    O\left(
       N\kappa_{\mathbb K}\log_2\kappa_{\mathbb K}
       +\kappa_{\mathbb K}^4
    \right)d^2\bb
    =
    O\left(
       N\kappa d\log_2\kappa_{\mathbb K}
       +\frac{\kappa^4}{d^2}
    \right)\bb.
\]
The \(\log_2\kappa_{\mathbb K}\) factor counts the packed
coefficient levels; evaluating their aggregated polynomials
accounts for the additive term.
The least-power-of-two choice gives
\(h+1\leq d<2(h+1)\).
Since \(\log\underline q=\Theta(\log q)\) and \(L=\Theta(N)\),
both \(h\) and \(d\) are \(\Theta(\log_qN)\).
Combining these contributions gives the stated
\(O(N\kappa\log_qN\log_2\kappa)\bb+O(\kappa^4)\) bound.
The copy computation stays in its tower basis; only interface
queries and messages are converted to the fixed commitment basis.
These conversions cost \(O(\kappa^2)\bb\) per extension-field
element and are covered by the same bounds.

\smallskip\noindent\textbf{Verifier time and communication.}
The \(O(h)\) packed column checks and combined query evaluations
cost \(O(\kappa h\log_2N)\ee\).
The two grand-product reductions add
\(O(\kappa_{\mathbb K}\log_2\kappa_{\mathbb K}
+(\log_2N)^2)\ee\).
Direct identity-label and public-I/O evaluations contribute
\(O(\kappa_{\mathbb K}+n_{\mathsf{io}}+\log N)\ee\).
These costs give the verifier bound.

The packed column messages contain \(O(\kappa h)\)
extension-field elements.
The two grand-product transcripts contain
\(O(\kappa_{\mathbb K}\log_2\kappa_{\mathbb K}
+(\log_2N)^2)\) such elements.
The remaining Hadamard and batched inner-product transcripts
are absorbed by these terms.
Using \(h=O(\log_qN)\) and
\(\kappa_{\mathbb K}\leq\kappa\) gives the stated proof size.

\smallskip\noindent\textbf{Prover memory and PCS costs.}
The grand-product workspace uses
\(O(dN+d\kappa_{\mathbb K}^3)
=O(dN+\kappa^3/d^2)\) base-field coordinates.
The witnesses, circuit tables, and grouped opening buffers
occupy \(O(Nh)\) coordinates.
Sequential execution reuses the gate and copy workspaces.
Including the optional interpolation tables and public I/O
gives \(O(N\log_qN)\Field+O(\kappa^4+n_{\mathsf{io}})\).

The four PCS functions bound aggregate resources for one online
witness commitment and one opening of each of the two vectors,
of logical lengths \(3m\) and \(m_{\mathsf{pub}}\).
They include domain padding and both retained commitment states;
public commitment generation is preprocessed.
Their total argument is therefore \(3m+m_{\mathsf{pub}}\).
This proves Eq.~\eqref{eq:together-single-efficiency}.
Unlabeled time and memory remainders count base-field operations
and base-field elements, respectively.
\end{proof}

\subsection{Proof of Theorem~\ref{theorem:multiple-instance-packing}}
\label{appx:proof-multiple-instance}

\begin{proof}
\noindent\textbf{Completeness.}
Let \(\bm R\in\Field^{N\times\kappa}\) contain the gate
residuals of the \(\kappa\) instances as columns.
The left-hand side of Eq.~\eqref{eq:together-multiple-gates}
is \(\bm R\ext{\bm\beta}\), including the term
\(\bm q_C\bm1_\kappa^\top\ext{\bm\beta}\).
Linear independence of the basis gives
\(\bm R\ext{\bm\beta}=\bm0_N\) if and only if
\(\bm R=\bm0_{N\times\kappa}\).
Similarly, packing commutes with the common row permutation,
so \(\sigma(\ext{\bm w})=\ext{\bm w}\) is equivalent to
the copy constraints in every instance.
This equivalence includes the packed public I/O.
Eq.~\eqref{eq:together-multiple-product-opening} authenticates
the virtual product vector from
Eq.~\eqref{eq:together-multiple-wires}.
Thus all scalar checks and tree relations accept for an honest prover.

\smallskip\noindent\textbf{Public packing.}
Let \(h=\lceil\log_{\underline q}L\rceil\).
Form \(\bm Q\in\Field^{N\times\kappa}\) with first five columns
\(\bm q_L,\bm q_R,\bm q_M,\bm q_O,\bm q_C\).
The next three blocks of \(h\) columns contain the position digits
\(\bm\sigma_{(1)},\ldots,\bm\sigma_{(h)}\), restricted to the
three length-\(N\) wire segments, respectively.
Fill the remaining columns with zeros.
The capacity condition \(5+3h\leq\kappa\) therefore gives:
\[
    \bm Q\ext{\bm\beta}\in\ExtensionField^N.
\]
A selector query uses its coordinate unit vector as the column
weight in Eq.~\eqref{eq:pcs-opening-relation}.
For the permuted-label query on segment \(\ell\in[3]\), use
\(\ext\rho_{5+(\ell-1)h+j}=\ext\beta_j\) for \(j\in[h]\),
and zero elsewhere.
Thus five selector queries and three permuted-label queries
are weighted evaluations of the same public matrix.
The row weights are tensor vectors, so Section~\ref{sec:pcs}
applies to all eight queries.
The public commitment is preprocessed.
The identity-label MLE is evaluated directly by
Proposition~\ref{proposition:position-encoding-evaluation};
it is not committed.

\smallskip\noindent\textbf{Gate soundness.}
Condition on binding of the witness commitment.
Its three length-\(N\) blocks have unique matrices of
\(\Field\)-coordinates in the fixed basis.
Thus the gate residual and \(\ext{\bm w}_M\) are determined
by the committed witnesses.
If some gate is invalid, the packed residual is nonzero.
Its random length-\(N\) tensor projection vanishes with
probability at most \(\log_2N/|\ExtensionField|\).

For a fixed projection point \(\ext{\bm x}\), the gate
sumcheck polynomial is the product of the multilinear equality
weight with the residual expression formed from selector and wire
MLEs. Each selector--wire product has individual degree two,
and multiplication by the equality weight raises this to three.
The constant-selector term has lower degree, giving sumcheck
error at most \(3\log_2N/|\ExtensionField|\).
The remaining virtual-product evaluation is checked by one
packed Hadamard proof on matrices of size \(N\times\kappa\).
Since the total number of entries is \(N\kappa\), its
reduction error is bounded by:
\[
    \frac{2\kappa-2}{|\ExtensionField|-\kappa}
       +\frac{2\log_2(N\kappa)}{|\ExtensionField|}.
\]
All terminal queries are authenticated by the batched openings below.
The bound applies jointly to all \(\kappa\) instances because
packing is injective.

\smallskip\noindent\textbf{Copy fingerprint.}
Suppose that the packed assignment violates a copy constraint.
Consider the bivariate difference polynomial:
\[
    P(Y,Z)=
       \prod_{i=1}^{L}(Z+\iota_i+Y\ext w_i)
       -\prod_{i=1}^{L}(Z+\varsigma_i+Y\ext w_i).
\]
If \(P\) were zero, unique factorization in
\(\ExtensionField[Y,Z]\) would equate the two multisets of
linear factors, each monic in \(Z\).
Since the labels are distinct and
\(\varsigma_i=\iota_{\sigma(i)}\), matching the constant
coefficients pairs position \(i\) with \(\sigma^{-1}(i)\).
Matching the \(Y\)-coefficients then forces
\(\ext w_i=\ext w_{\sigma^{-1}(i)}\), contradicting the copy
violation. Hence \(P\ne0\).

Both products have homogeneous degree-\(L\) part
\(\prod_{i=1}^{L}(Z+Y\ext w_i)\), so \(P\) has total degree
at most \(L-1\).
The independent uniform challenges \(\ext y,\ext z\)
therefore satisfy \(P(\ext y,\ext z)=0\) with probability at
most \(\min\{1,(L-1)/|\ExtensionField|\}\).
Padding each product-tree input with ones does not change
this fingerprint.

\smallskip\noindent\textbf{Product-tree soundness.}
Set \(\mu=\lceil\log_2L\rceil=\log_2(4N)\).
For one grand product, unfold the binary tree from its root
to the padded length-\(4N\) input vector.
At a level with \(2^\ell\) parent entries, where
\(\ell\in\{0,\ldots,\mu-1\}\), let
\(\ext{\bm x}\in\ExtensionField^\ell\) be the current
evaluation point.
The weighted product identity uses the polynomial:
\[
    \left\langle
       \tensorspan{\ext{\bm x}},\tensorspan{\bm X}
    \right\rangle
    \left\langle\ext{\bm c}_U,\tensorspan{\bm X}\right\rangle
    \left\langle\ext{\bm c}_D,\tensorspan{\bm X}\right\rangle,
\]
where \(\ext{\bm c}_U,\ext{\bm c}_D\) are the two child subvectors
and \(\bm X\) denotes \(\ell\) formal variables.
Its degree in each variable is at most three.
Sumcheck reduces its Boolean sum to child evaluations
\(\ext v_U,\ext v_D\) at one point \(\ext{\bm y}\), with
error at most \(3\ell/|\ExtensionField|\).
After both values are fixed, sample a fresh
\(\ext y_0\sample\ExtensionField\) and form the next claim:
\[
    (1-\ext y_0)\ext v_U+\ext y_0\ext v_D
       =
    \left\langle
       \ext{\bm c}_U\mathbin{\|}\ext{\bm c}_D,
       \tensorspan{(\ext y_0,\ext{\bm y})}
    \right\rangle.
\]
If either child claim is false, this merge accepts it with
probability at most \(1/|\ExtensionField|\).
Thus the two trees have total reduction error at most:
\[
    2\sum_{\ell=0}^{\mu-1}\frac{3\ell+1}{|\ExtensionField|}
       =\frac{3\mu(\mu-1)+2\mu}{|\ExtensionField|}.
\]
The \(\ell=0\) level has no sumcheck rounds but still uses
a fresh merge challenge.

Each tree ends in a single MLE claim on its padded input.
At its terminal point, the known shift and padding contribute:
\[
    \ext z\sum_{i=1}^{L}
       \bigl(\tensorspan{\ext{\bm x}}\bigr)_i
    +\sum_{i=L+1}^{4N}
       \bigl(\tensorspan{\ext{\bm x}}\bigr)_i
    =
    1+(\ext z-1)\sum_{i=1}^{L}
       \bigl(\tensorspan{\ext{\bm x}}\bigr)_i.
\]
The prefix sum takes \(O(\log N)\) operations by a binary
interval decomposition.
Each length-\(N\) wire segment fixes two prefix bits.
Retain their tensor weight and query the suffix MLE.
The public-I/O values and permuted labels at the remaining
\(n_{\mathsf{io}}\) positions are evaluated directly.
The identity-label MLE is computed by Appendix~\ref{appx:mleiota}.
Consequently, the two copy terminals leave six ordinary wire
queries and three weighted queries to the public matrix.

\smallskip\noindent\textbf{Batched openings and composition.}
There are nine ordinary wire queries: three from the gate check
and six from the copy check.
The two Hadamard terminals and eight public queries give
ten packed queries.
Apply the column-projection step of Section~\ref{sec:pcs}
to each packed query.
Fix all column messages and ordinary claimed values before
sampling independent projection vectors for the packed queries
and independent scalar weights for the ordinary queries.
Extend each query by zeros outside its committed block, and
sum the resulting inner-product claims within each group.

If an ordinary claim is false, its scalar weight has a nonzero
error coefficient.
If a packed column message is false and its scalar check passes,
injectivity gives
\((\bm U-\bm\Xi^\top\bm A)\ext{\bm\beta}\ne\bm0\).
Thus a group containing a false claim has a nonzero linear
error polynomial in its independent batching challenges.
Cancellation occurs with probability at most
\(1/|\ExtensionField|\) per group.

The witness concatenation has logical length \(3N\), padded
to \(4N\); the public packing has length \(N\).
One invocation of Fact~\ref{fact:ip} per group therefore adds
\(2\log_2(4N)/|\ExtensionField|\) and
\(2\log_2N/|\ExtensionField|\) error, respectively.
The verifier computes the combined query-side evaluation in
each group exactly: use the recurrence of Section~\ref{sec:pcs}
for packed queries and tensor contractions for ordinary
queries, retaining the witness block-prefix weights.
This leaves two underlying PCS openings.
Their joint binding and opening error is
\(\epsilon_{\mathsf{PCS}}\).

Combining these bounds gives:
\[
\begin{aligned}
    \epsilon_{\mathsf{multi}}
    \leq{}&\epsilon_{\mathsf{PCS}}
       +\frac{2\kappa-2}{|\ExtensionField|-\kappa}\\
       &+\frac{
          L-1+3\mu(\mu-1)+2\mu
          +8\log_2N+2\log_2(N\kappa)+6
       }{|\ExtensionField|}.
\end{aligned}
\]
Here the constant \(6\) includes the two aggregation checks
and the extra prefix rounds for the length-\(4N\) witness.
Since \(L\leq4N\), \(\kappa\leq N\), and
\((\log_2N)^2=O(N)\), this proves
Eq.~\eqref{eq:together-multiple-soundness}.

\smallskip\noindent\textbf{Efficiency.}
Each wire matrix has \(N\kappa\) entries.
Computing \(\bm W_L\circ\bm W_R\) costs \(O(N\kappa)\bb\).
The Hadamard reduction and ten packed column projections each
cost \(O(N\kappa^2+\kappa^3)\bb\), which is
\(O(N\kappa^2)\bb\) because \(N\geq\kappa\).
The assumption \(q\geq2\kappa-1\) supplies the interpolation
nodes over \(\Field\), so no intermediate-field table is needed.
The gate sumcheck and the two final inner-product reductions
use \(O(N)\) extension-field operations.
The product-tree layer lengths sum to \(O(N)\), and their
sumchecks cost \(O(N+\mu^2)\) extension-field operations.
Under schoolbook arithmetic, these contributions give
\(O(N\kappa^2)\bb\) total algebraic prover time.

The packed Hadamard check uses \(O(\kappa)\ee\) verifier time,
and the ten column projections and terminal query contractions
use \(O(\kappa\log_2N)\ee\).
The gate, tree, and final inner-product checks add
\(O((\log_2N)^2)\ee\).
Packing the public I/O reads \(O(\kappa n_{\mathsf{io}})\)
base-field coordinates, and its terminal contractions and
the direct identity-label evaluations take
\(O(n_{\mathsf{io}}+\log N)\) extension-field operations.
This gives the stated verifier time.

The ten column messages contain \(O(\kappa)\)
extension-field elements in total.
The Hadamard proof, gate sumcheck, tree reductions, and final
inner-product proofs add
\(O(\kappa+(\log_2N)^2)\) extension-field elements,
giving the proof-size bound.
The witnesses and public matrix occupy \(O(N\kappa)\)
base-field coordinates.
Sequential execution retains \(O(N)\) extension-field
tree and sumcheck entries and \(O(\kappa^2)\) coordinates
for the column messages.
Thus algebraic prover memory is \(O(N\kappa)\Field\).

The online PCS costs cover one witness commitment and two
openings on vectors of logical lengths \(3N\) and \(N\).
Their aggregate argument is therefore \(4N\), with the required
witness padding included in the PCS cost functions.
Public commitment generation is part of circuit preprocessing;
both retained commitment states count toward PCS memory.
These calculations prove
Eq.~\eqref{eq:together-multiple-efficiency}.
Dividing the algebraic prover time and peak memory by
\(\kappa\) gives the amortized per-instance bounds.
\end{proof}


\end{document}
